% =====================================================================
%  Sophus Heegaard, Professor R. Nielsens Lære om Tro og Viden. En
%  kritisk Afhandling. Kjøbenhavn: Forlagt af den Gyldendalske Boghandel
%  (F. Hegel), G. S. Wibes Bogtrykkeri, 1867. 98 pp.
%  DIPLOMATIC TRANSCRIPTION.
%
%  SOURCE. KB digitisation, copy DA 1.-2.S 3 8°, barcode 11030800293A
%    (KB online: https://www.kb.dk/e-mat/dod/11030800293A-bw.pdf).
%    Library file: Heegaard, Sophus/11030800293A_bw-2.pdf
%    sha256 b7ac81db04a9fa1bcd33fbc2a7de130c95aa9ffe7614697ee57b9d2427f6fef4
%    (SCANS.tsv key heegaard/nielsens-laere-tro-viden). 105 PDF pages,
%    A4; each page a 300 ppi 1-bit image (renders above 300 dpi are
%    interpolation and were used only to see glyph topology).
%    Page map: printed p. N = PDF N+5, uniform for pp. 1--98. The folio
%    was read at the image on every page PDF 9--103; the sheet signatures
%    "1*" p. 3, "2" p. 17, "3" p. 33, "6" p. 81, "7" p. 97 agree.
%    PDF 6 = title page [1]; PDF 7 = blank verso [2] (library stamp only);
%    PDF 8 = Forord p. [3], unfolioed; PDF 103 = p. 98. PDF 1 = KB notice,
%    2 = front board, 3 = pastedown (label, stamp), 4--5 blank flyleaves,
%    104 = blank end-paper, 105 = back board. No contents page, no running
%    heads, no errata. Structure: Forord pp. [3]--9; I. Metaphysiken
%    p. 10; II. Videnskabslæren p. 20; III. Ideelæren p. 39; IV. Den
%    philosophiske Methode p. 89 (table pp. 94--95); Efterskrift
%    pp. 96--98.
%
%  TYPE. Fraktur body. Letterspacing (Sperrsatz) → \emph{}; bold (fette)
%    Fraktur → \textbf{}; antiqua words (Latin tags) and antiqua letter
%    labels (A., B., a), b), c)) → \textit{}; antiqua numerals, Arabic or
%    Roman, are left plain in running text and footnotes (as the corpus
%    leaves antiqua "p." plain) and marked only in the section heads and
%    the table, where the change of face is part of the display. Emphasis was
%    judged by measured inter-letter advance within the line, not by
%    impression. Footnotes are marked *), **) and restart on every page,
%    hence \setcounter{footnote}{0} at each \opage; a note that runs over
%    onto the next page is kept whole at its anchor. Quotation marks „ “
%    as printed. The Fraktur „ꝛc.“ → \&c. Numeric ranges are printed with
%    the same long rule as the em-dash (measured: 87--99 px in the body,
%    69--79 px in footnotes, at 300 ppi) and are set --- accordingly.
%    Fraktur I/J → I for I-words, J for J-words; long s → s; ß kept in
%    German. \opage{N} stands at the exact word where printed page N
%    begins, mid-word where the page turns mid-word; a page turn is never
%    a paragraph break. Words broken at a line end are joined; printed
%    suspended hyphens are kept.
%
%  STRAY MARKS. Settled by one measured rule, applied to every site in
%    the book (2026-09-23): a point-shaped mark of at least half the ink
%    area of the page's median full stop is a printed point and is
%    transcribed where it prints, as · (raised) or . (on or below the
%    baseline), with a % PRINTED AS IS: comment giving the measurement;
%    smaller marks are specks, not transcribed, and noted (% NOTE:) only
%    where they could be taken for punctuation. Transcribed: pp. 13, 29,
%    30, 45, 67n, 68n, 69 (points); p. 81 (an apostrophe-shaped sort).
%    Marks called points in the first reading and found to be specks
%    (not transcribed, NOTE at the site): pp. 14, 19, 65 ("hvad der",
%    "den, at"), 88. Semicolons found to be comma + speck: pp. 19, 29;
%    p. 17 "uindskrænket," likewise; p. 93 "Ulige" takes no comma.
%
%  QUOTE BALANCE. 350 „ against 349 “: p. 41, the quotation opened at
%    „Gribe vi Ideen is never closed in print (PRINTED AS IS). Do not
%    "fix" it.
%
%  PRINTER'S ERRORS kept as printed and marked % PRINTED AS IS: at the
%    site: p. 12 "Interesseu"; p. 13n "Kbhvn, 1865"; p. 32n
%    "nogetsom-/somhest"; p. 35 "Virkelig-/ligheden"; p. 36n "R Nielsen.";
%    p. 44n "-- hvad"; p. 47 "giennemtrængt"; p. 52 "Lvgik"; p. 54n
%    "afgjorte"; p. 67/68 "be-/temmes"; p. 80 "Bevidsthed- er"; p. 81 bold
%    "saa"; p. 94 stray "-" before "="; p. 95 "(Grundideernes Logik II.)"
%    without full stop after Logik, and "Grundidecrnes" (a complete c);
%    p. 95 "ndgivet"; p. 98 "til ikke udtale"; plus the stray points above.
%
%  THE TABLE (pp. 94--95) is reproduced as a table: label column with the
%    printed number of spaced leader points (counted at 600 dpi), "="
%    column aligned as printed, fractions as built-up fractions, ditto
%    marks as '', the right brace joining rows a./b., the full-width rule
%    and the totals line; the display equation keeps its printed left
%    indent. Layout comments before each table give the measurements.
%
%  METHOD (2026-09-23). Transcribed from the page images in 12 batches
%    (pp. [1]--9; 10--19; 20--29; 30--38; 39--48; 49--58; 59--68; 69--78;
%    79--88; 89--93; 94--95, the table; 96--98), each by an independent
%    subagent working from 300 ppi renders and 600--900 dpi crops.
%  Word-accuracy: every batch diffed against the scan's embedded ABBYY
%    layer before splicing (ocrdiff.py --offset 5). 12 batches, 248
%    candidates, 0 transcription errors corrected by the diff at the batch
%    stage, 248 judged the witness's fault (the layer reads æ as ce/oe, ø
%    as o/e/v, sk as st, and splits letterspaced words), 0 unresolved.
%    One of those 248 was misjudged: p. 29 "ifleng", where the layer's
%    "iflæng" was right; found by the control sample below and corrected.
%    See TRANSCRIPTION-PLAYBOOK.md §3 step 4.
%  Second reading: every page was then read again at the image, word by
%    word and mark by mark, by a different agent (12 readers), and every
%    disagreement (72 in all) was settled at higher resolution by a third
%    reader. Words: 0 wrong or dropped words found in 98 pp.; one silent
%    normalization of a printer's error corrected (p. 95 heading
%    "Grundideernes" → "Grundidecrnes" as printed). Emphasis: 5
%    corrections (p. 15 "Troen" not spaced; p. 76 "Grund" spaced; p. 78
%    first "en" spaced; p. 90 "og" not spaced; p. 91 "er" spaced).
%    Antiqua: the numerals I.--IV. of the four section heads and the
%    table labels (c), I., II., B., a., b., (I), (II), (R), (I+II)).
%    Table: 5 leader counts (each one too many), leader and column
%    positions, display indent. Punctuation and stray marks: 15 sites
%    re-decided under the rule above. Layout: p. 3 \opage before the
%    heading; pp. 69--70 list items set on their own lines as printed;
%    double rules at the end of the Forord and of the book. Paragraphing
%    and page-marker positions confirmed on every page.
%  Control and targeted screens (after the second reading). A seeded
%    random sample of 10 pp. (6, 29, 34, 40, 47, 50, 56, 66, 69, 70) read
%    at the image by a fresh reader found 2 errors: p. 29 "ifleng" →
%    "iflæng", p. 70 "af" letterspaced. Two whole-book screens followed:
%    (1) every place where the ABBYY layer has æ/ø/å, a letter or a word
%    that the transcription lacks (the layer loses such things, it does
%    not invent them): 1 real error (p. 29, above); (2) every word the
%    layer splits into letters (its sign of letterspacing) against the
%    \emph markup, settled by measuring inter-letter gaps: 4 emphasis
%    corrections (p. 9 "den videnskabelige", p. 70 "udledes af", p. 77
%    "Charakteer af", p. 96 "ikke nøies med"); and a sweep of letter
%    labels, 7 marked antiqua (pp. 32n, 69, 70). A second seeded sample of
%    8 pp. (8, 15, 21, 39, 53, 81, 89, 95) then found 0 errors.
%  Accuracy: full second reading of all 98 pp. at the image, third-reader
%    adjudication, two whole-book machine screens, and 18 random pp. of
%    control reading. Word-level errors found after the batches: 2 (p. 29
%    "iflæng"; p. 95 "Grundidecrnes" normalized), both corrected; emphasis
%    errors found: 9, corrected. None known outstanding. Residual rate
%    after correction: 0 errors in the final 8-page control sample.
%    [2026-09-23] See ACCURACY-LEDGER.tsv.
%
%  OPEN ITEMS (see the % UNSURE: notes in the text). p. 23 whether the
%    "i" in "er i og under" is spaced; p. 30 "Braaden" (â-like mark,
%    damage?); pp. 33--34n a heavier "i" before "et" and bold "Over" in
%    "Overvidenskabeligt"; p. 37 "der" with a damaged d (or "ter"); p. 52
%    "Lvgik" (wrong sort or damaged o); p. 58 "At" spaced, "den" not;
%    p. 72 blot at "ogsaa" (not taken as a full stop); p. 79 "alene"
%    printing as "aleue"; p. 87 "Philosophi" with the second o open
%    ("Philoscphi"), taken as damaged o.
% =====================================================================

\documentclass[12pt,a4paper]{book}
\usepackage[utf8]{inputenc}
\usepackage[T1]{fontenc}
\usepackage{libertinus}
\usepackage{libertinust1math}
\usepackage{amsmath}     % the table pp. 94--95
\usepackage[danish]{babel}
\usepackage[protrusion=true,expansion=true]{microtype}
\usepackage{setspace}
\usepackage[margin=1.2in]{geometry}
\usepackage{enumitem}
\usepackage{fancyhdr}
\setlength{\headheight}{28pt}
\addtolength{\topmargin}{-13.5pt}
\usepackage{hyperref}

\hypersetup{
  pdftitle={Professor R. Nielsens Lære om Tro og Viden},
  pdfauthor={Sophus Heegaard},
  hidelinks
}

% Footnotes are printed *), **), ***) and restart on every page.
\renewcommand{\thefootnote}{\ifcase\value{footnote}\or*\or**\or***\or****\or*****\fi)}

\pagestyle{fancy}
\fancyhf{}
\fancyhead[LE,RO]{\thepage}
\fancyhead[RE]{\textit{Professor R. Nielsens Lære om Tro og Viden}}
\fancyhead[LO]{\textit{Sophus Heegaard}}

\setstretch{1.3}

\usepackage{marginnote}
\newcommand{\opage}[1]{\marginnote{\footnotesize\textit{[#1]}}}

\begin{document}
% ===== batch b01_pp01-09 =====
\thispagestyle{empty}
% TITLE PAGE (printed [1], PDF 6). No folio. Printed layout, top to bottom:
%   1. "Professor R. Nielsens Lære" -- large textura (black-letter display), centred.
%   2. "om Tro og Viden." -- same large textura, centred.
%   3-5. Epigraph, small Fraktur, set right of centre, three lines:
%        "En Philosophie, der er i aabenbar Modsigelse" /
%        "med „den sunde Forstand“ er falsk." /
%        attribution, indented: "R. Nielsen. Philosophisk Propædeutik. S. 41."
%        (the l of "Philosophie" in line 3 is damaged, reads like a broken i/l).
%   6. "En kritisk Afhandling" -- medium Fraktur, centred.
%   7. "af" -- small Fraktur, centred.
%   8. "Dr." in antiqua + "P. S. V. Heegaard." in textura; "Heegaard" letterspaced.
%   -- A circular stamp (belt-and-buckle garter enclosing a script monogram,
%      lettering round the ring) sits in the middle of the page: a library /
%      ownership stamp, NOT part of the print; not transcribed.
%   9. "Kjøbenhavn." -- Fraktur, centred.
%  10. "Forlagt af den Gyldendalske Boghandel (F. Hegel)." -- heavier
%      Schwabacher-style display face, centred; the closing parenthesis
%      is malformed (prints like a broken "}" / ")"), read as ")".
%  11. "G. S. Wibes Bogtrykkeri." -- small Fraktur; "Wibes Bogtrykkeri" letterspaced.
%      NOTE: speck below the line under „B“/„o“ of „Bogtrykkeri“, 51 px vs full stop F=140 px, top 22 px below the baseline; not transcribed.
%  12. "1867." -- antiqua numerals, centred.
\begin{center}
{\Huge Professor R. Nielsens Lære}\par
\bigskip
{\Huge om Tro og Viden.}\par
\bigskip
\hfill\begin{minipage}{0.62\textwidth}\small
En Philosophie, der er i aabenbar Modsigelse\\
med „den sunde Forstand“ er falsk.\\
\hspace*{1.5em}R. Nielsen. Philosophisk Propædeutik. S. 41.
\end{minipage}\par
\vspace{2\baselineskip}
{\Large En kritisk Afhandling}\par
\medskip
af\par
\medskip
{\LARGE \textit{Dr.} P. S. V. \emph{Heegaard}.}\par
\vspace{4\baselineskip}
{\large Kjøbenhavn.}\par
\medskip
Forlagt af den Gyldendalske Boghandel (F. Hegel).\par
\smallskip
G. S. \emph{Wibes Bogtrykkeri}.\par
\smallskip
\textit{1867.}\par
\end{center}
\clearpage
% Printed [2] (PDF 7): blank verso, no folio; only a circular library stamp
% (BIBLIOTHECA REGIA HAFNIENSIS, with crown), which is not print. Not transcribed.
% --- p. 3 ---
% (PDF 8; no folio printed; signature mark "1*" at foot)
\opage{3}\setcounter{footnote}{0}%
\begin{center}\textbf{Forord.}\end{center}
% Head "Forord." set in a heavier, letterspaced Schwabacher-style display
% face, centred; a short single rule is printed centred beneath it.
% The text opens with an enlarged initial N (not marked up).
Naar man seer tilbage paa Vendepunkterne i Philosophiens
historiske Udvikling for deraf at drage Slutninger
med Hensyn til vor egen Tid, saa tør man vistnok sige,
at der forestaaer en gjennemgribende Revision af Vidensideen
og dens Betydning for Personligheden. Philosophiens
Historie viser, at ethvert afgjørende Fremskridt
er betegnet ved to Perioder: en Forberedelsens Tid og
selve Gjennembrudet. I den forberedende Periode har
det Gamle tabt Evnen til at vække den friske umiddelbare
Begeistring, som efterhaanden er vegen for Kritiken;
de Kampe, som Kritiken kalder tillive, kunne vel
% NOTE: speck („i Lidenskab“), 60 px vs full stop F=145 px, below the baseline (top 19 px below) after „i“; not transcribed.
føres i Lidenskab, men Lidenskabeligheden i denne Periode
er i Almindelighed Kritikens, ikke Begeistringens
Lidenskab. Ere vi nu i den forberedende Periode eller i
Gjennembrudet? Vi ere ikke i Gjennembrudet. Vel
er Kampen igang, men der savnes et almindeligt
Feldtraab; det er ikke to Armeer, der kæmpe i sluttet
Slagorden, men større og mindre Afdelinger, der udfægte
Striden paa egen Haand. I Gjennembrudets
Øieblik er Talen kort og bestemt og samler sig gjerne i
% --- p. 4 ---
% (PDF 9)
\opage{4}\setcounter{footnote}{0}%
en Prægnants, som har uendelig befrugtende Magt, og
som formaaer at aabne en heel ny Verden for Bevidstheden;
ved Navnet Plato f. Ex. tænke vi uvilkaarligt
paa „Ideen“, ved Navnet Kant paa „synthetiske Domme
\textit{a priori}“. At fremkunstle noget Saadant lader sig aldeles
ikke gjøre, dertil er Sandhedens Magt for stor; vel kunne enkelte
Partier i kortere eller længere Tid vække Opmærksomhed
ved den bestandige Gjentagelse af et og samme Stikord,
men man kan altid forlade sig paa, at Sagen kun
bestaaer Prøven ved sig selv. Saavist som Viden er
Magt, saavist vil den ogsaa forstaae at gjøre Magten
gjældende.

Befinde vi os nu i den \emph{forberedende} Periode, saa
skade vi kun os selv og Sagen, naar vi ikke overholde
den for Forberedelsen gjældende Lov. Gjennembrudets
Charakteermærke er Produktivitet, Forberedelsens Charakteermærke
er Kritik. Blandt dem, der hos os i Henseende
til Fremskridt og Udvikling staae i første Række,
maa nævnes Professor R. Nielsen. Man dømme nu
om hans Philosophi, som man vil, man være uenig
med ham i Enkeltheder eller i Hovedsagen, dette er dog
et Faktum, at han allerede længe har holdt og fremdeles
holder Bevægelsen igang; det er mod ham, at Angrebene
rettes, og det er mod ham, at de fremdeles \emph{bør}
% NOTE: speck („ligger der“), 22 px vs full stop F=198 px, at mid x-height right after „ligger“; not transcribed.
rettes; men i dette „bør“ ligger der et Dobbelt, deels
nemlig, at han yder noget Værdifuldt, forsaavidt Samtiden
vedblivende anseer det for Umagen værd at bekæmpe
det, deels at han betegner et Gjennemgangspunkt,
hvorved man ikke kan blive staaende. At ville hævde
Professor Nielsens Philosophi Betydningen af et Gjen%
% --- p. 5 ---
% (PDF 10)
\opage{5}\setcounter{footnote}{0}%
nembrud i samme Forstand, som naar vi tale f. Ex.
om den Kantiske Philosophi, gaaer selvfølgelig ikke an.
Vel indeholder den værdifulde, ja geniale Enkeltundersøgelser,
men den kan neppe gjøre Fordring paa noget
overraskende nyt Princip, der som et Banner rager
høit iveiret og samler Alt omkring sig. Vel indeholder
den frugtbare Elementer, navnlig for nye Bearbeidelser
af Naturphilosophien, men man kan ikke med
Rette paastaae, at den bringer et nyt Foraar i den hele
Aandens Verden, eller at den lader alt det Gamle
foreløbigt glemmes. De, som paastaae dette, maatte i
al Fald føre Beviset ved i faa Ord at vise, hvori
den nye Opdagelse bestaaer og nævne den ved Navn.
Det er tildeels en saadan Paastand, som har fremkaldt
den her foreliggende Kritik, idet det er blevet Forfatteren
indlysende, at Professor Nielsen staaer som
en af Anførerne i den hele Bevægelse \emph{henimod} et
Gjennembrud, men at dette selv ingenlunde er indtraadt.
Det er en besynderlig Lyst, der undertiden
griber Videnskabsmændene til at stille sig under Dilemmaet
\emph{Alt eller Intet}, og til at see ned paa de Hæderspladser,
de med Rette kunne indtage, selv om de
ikke ere en Kant eller en Newton.

At give en udtømmende Kritik af Professor Nielsens
Philosophi ligger ikke i dette Arbeides Plan; dets Hovedformaal
er at kritisere hans Lære om Tro og Viden,
og Undersøgelserne berøre kun forsaavidt hans Philosophi
i videre Forstand, som den staaer i Forbindelse med
hiint Hovedspørgsmaal. Kritiken gaaer da ud paa at
vise, at Professor Nielsens Lære om \emph{Tro og Viden}
% --- p. 6 ---
% (PDF 11)
\opage{6}\setcounter{footnote}{0}%
fører ind i Modsigelser, hvad enten vi henvende os
til Kategoribestemmelserne i hans egen \emph{Metaphysik},
eller til hans \emph{Videnskabslære}, eller til hans \emph{Ideelære},\footnote{Vi have valgt disse Synspunkter for Udviklingen, fordi de
hos Professor Nielsen selv træde frem som Hovedpunkter.}
eller endelig til hans hele philosophiske \emph{Methode}.

I Professor Nielsens \emph{Metaphysik} faaer den hele
Lære om \emph{Grændsen} en fremragende Betydning, hvad
der jo iøvrigt ogsaa stemmer med de Fordringer, Begrebet
% NOTE: speck („Natur maa“), 36 px vs full stop F=163 px, at upper x-height in the word space; not transcribed.
selv ifølge sin egen Natur maa gjøre. Er nu Grændsen, som
hos Professor Nielsen, bestemt dialektisk som den, der
paa eengang er Eet med og forskjellig fra det Begrændsede,
saa er det en Selvmodsigelse at sætte \emph{Mysteriet}
som Grændse mellem Tro og Viden, navnlig naar Mysteriet
\emph{hæves til Kategori}; thi mellem Mysteriet og
\emph{Troen} mangler \emph{Forskjellen} (den ene Bestemmelse i
Grændsen), mellem Mysteriet og \emph{Viden} mangler \emph{Eenheden}
(den anden Bestemmelse i Grændsen), saa at
Mysteriet hverken i den ene eller i den anden Henseende lader
sig sætte som Grændsekategori. Overhovedet er det en
Selvmodsigelse at \emph{vide} Noget om et \emph{Mysterium}.

Den abstrakte Charakteer, som de metaphysiske Undersøgelser
om Grændsen ifølge Sagens Natur maatte
have, bliver konkret i \emph{Videnskabslæren}. Idet der her
stilles en skarp Modsætning mellem Videnskabens „upersonlige“
og Religionens „personlige“ Sandheder, bliver
\emph{Grændsen} mellem begge det egentlige \textit{punctum quæstionis}.
Men inden dette Problem kan blive undersøgt, maa først
det Punkt være klaret, om der overhovedet existerer reent
% --- p. 7 ---
% (PDF 12)
\opage{7}\setcounter{footnote}{0}%
„upersonlige“ Sandheder. Dette benegte vi, fordi Sandheden
ikke kan existere objektivt, men kun \emph{er}, forsaavidt
den \emph{vorder}. Som Udtryk for et Vexelforhold mellem
en erkjendende Personlighed og et Objekt, som bliver
erkjendt, bærer Sandheden altid Mærket af sin Tilblivelse,
og falder som subjektiv-objektiv aldrig udenfor
Personligheden. Mod den Antagelse, at de logisk-physiske
Sandheder kun erkjendes ad den rene Tænknings
Vei, gjøre vi gjældende, at \emph{al} Erkjendelse, ogsaa den
logisk-physiske, forudsætter den \emph{hele} Personligheds Virksomhed
og involverer ligesaavel Villie som Tænkning,
skjøndt Tænkningen her bliver det afgjørende Moment.
Naar Professor Nielsen vil opstille Ethiken som Grændsevidenskab
mellem hine skarpe Modsætninger, der repræsenteres
af de „upersonlige“ og de „personlige“ Sandheder,
saa gjøre vi gjældende \emph{deels}, at Grændsen heller
ikke her er klaret, fordi man savner al Oplysning om
selve Overgangen fra de logisk-physiske Videnskaber til
Ethiken og atter fra Ethiken til Religionslæren; \emph{deels},
at en Overgang her er umulig, naar Modsætningerne:
Personligt og Upersonligt, Tro og Viden ere stillede absolute.

At de Grundmodsigelser, der findes i Videnskabslæren,
maae komme igjen i \emph{Ideelæren}, skjøndt under
andre Former, er let forklarligt af Systemets indre
Sammenhæng. Til Kundskaben om Professor Nielsens
Ideelære have vi et vigtigt Indlæg i „Grundideernes
Logik.“ Hvad der imidlertid noget vanskeliggjør Undersøgelsen,
er deels den Omstændighed, at „Grundideernes
Logik“ er et endnu ufuldendt Værk --- skjøndt Forfatteren
% --- p. 8 ---
% (PDF 13)
\opage{8}\setcounter{footnote}{0}%
selv i Forordet erklærer, at han ikke paa denne Grund vil
afvise Kritiken --- deels dette, at \emph{Viden}, \emph{Magt} og
\emph{Sandhed} opstilles som Grundideerne uden gjennemført
Beviis for, at de virkelig ere det.

Vort første Indlæg i Problemet: Tro og Viden
gjælder her Ideens Forhold til Virkeligheden. Er Ideen
det Oprindelige, Virkeligheden det Afledede, saa kan man
ikke uden Modsigelse forudsætte Eenhed af Villie og
Tanke i Ideen og samtidigt bestemme Villie og Tanke
samt Tro og Viden som \textbf{absolut} \emph{ueensartede} i den
menneskelige Bevidsthed; thi i saa Tilfælde er Overgangen
fra det Oprindelige til det Afledede brudt. Skal
Forklaringen søges deri, at Grundideerne selv ere \emph{absolut}
ueensartede og just derfor i \emph{absolut} Eenhed,
maa dette Absolute først udvikles og paavises; men en
saadan Paaviisning mangler.

Vil „Grundideernes Logik“ \emph{begynde} Philosophien,
saa kan den ikke indskrænke sig til „den rene Viden“, men maa
i Udgangspunktet have Villien med; thi Villien er ligesaa
\emph{oprindelig} som Viden, og Villien er ligesaa \emph{virkelig}
som Viden. Vi søge derfor at paavise, at „Grundideernes
Logik“, naar den først har begyndt med „den rene Viden“,
paa intet Punkt af Systemet kan naae Villien
uden ved et Spring, og det saameget mere, som Forfatteren
selv bestemmer Viden og Villie som \textbf{absolut}
\emph{ueensartede}. Det er atter her \emph{Grændsen}, som aabenbarer
Modsigelsen i Forholdet mellem Tro og Viden.
Det er Professor Nielsens Tanke, at Viden vel ikke
ligefrem kan erstatte det, den maa opgive i Troen, men
at den dog i en „Videnstheisme“ kan „afspeile“ Troes%
% --- p. 9 ---
% (PDF 14)
\opage{9}\setcounter{footnote}{0}%
theismen og saaledes give „rationelle Ækvivalenter“ for
det Tabte. Vi see ikke rettere, end at dette er en Umulighed,
% NOTE: the second n of „Tænkning“ prints as two separate stems (3 px clear between them at 300 ppi), with neither the top arch of n nor the foot join of u, so it can look like „Tænkniug“; a broken n. Read n (an absent hairline is not evidence for another letter).
fordi der i den rationelle Tænkning Intet findes,
som med „logisk Nødvendighed“ fører til de positiv-religiøse
Bestemmelser i Gudsbegrebet. Omvendt kan
Troestheismen ligesaalidt afgive noget Grundlag for
% NOTE: speck („ingen Overgang“), 41 px vs full stop F=167 px, at x-height top in the word space before „Overgang“; not transcribed.
Videnstheismen; thi deels gives der ingen Overgang fra
Tro til Viden, som jo ere \textbf{absolut} ueensartede Principer;
deels ville de positiv-religiøse Begreber ganske
tabe deres Charakteer, naar de overføres i Videns
Sphære.

Lader det sig i de efterfølgende Undersøgelser godtgjøre,
at de antydede Modsigelser virkelig ere tilstede,
saa følger det ganske af sig selv, at \emph{den videnskabelige
Methode} ikke kan være uberørt deraf. Ere Tro
og Viden, Villie og Tænkning \emph{absolut} ueensartede
Principer, saa er \emph{Dialektiken} med det Samme umuliggjort.
Dialektikens væsentligste Bestemmelser, \emph{Grændse}
og \emph{Overgang}, ere ved Læren om det Absolut-Ueensartede
stillede saaledes, at den første ophæver den sidste. Heraf
bliver da Følgen, at Dialektikeren Professor Nielsen
af og til slaaer over i Dogmatisme.
% End of the Forord (p. 9): a double rule, centred, set well below the last
% line: two thin parallel rules of equal length (794 and 792 px at 300 ppi,
% c. 67 mm in the scan, 0.32 of the 2497 px measure), 16 px apart centre to
% centre, both equally ragged/broken along their length; skewed in the scan.
\begin{center}\rule[2.2pt]{0.32\textwidth}{0.4pt}\hspace*{-0.32\textwidth}\rule{0.32\textwidth}{0.4pt}\end{center}
% ===== batch b02_pp10-19 =====
% --- p. 10 ---
% (PDF 15)
\opage{10}\setcounter{footnote}{0}%
% [centred bold head, a short centred rule below it; enlarged initial „E“ not marked up]
\begin{center}\textbf{\textit{I.} Metaphysiken.}\end{center}

Et Hovedpunkt i al Metaphysik danner Læren om
% NOTE: speck („At det“), 37 px vs full stop F=176 px, at x-height top in the word space; not transcribed.
\emph{Grændsen}. At det forholder sig saaledes indsees let
deraf, at „Tanke ikke kan være uden Tænkning“, altsaa
heller ikke uden Bevægelse; men allerede i Bevægelsen
er Grændsebegrebet indeholdt, fordi Bevægelse forudsætter
Overgange. Som nu Alt det, der i et philosophisk
System fremkommer paa de senere Stadier i Udviklingen,
forudsætter og viser tilbage til Metaphysiken,
saaledes er der ogsaa Anledning til ved Troens og Videns
indbyrdes Begrændsning at spørge om deres Forhold
til Grændsens Udvikling i Metaphysiken.

Men inden vi besvare dette Spørgsmaal, maa det
først staae klart for Bevidstheden, hvad Professor Nielsen
i det Hele lærer om Troens Forhold til Viden. I
Korthed lader dette sig fremstille saaledes: Professor
Nielsen udvikler Tro og Viden som \textbf{absolut} \emph{ueensartede}
Principer, der i væsentlig Forstand ikke kunne
gjøre hinanden Afbræk, og som uden Modsigelse kunne
forenes i een Bevidsthed, ikke \emph{endskjøndt}, men \emph{fordi}
% --- p. 11 ---
% (PDF 16)
\opage{11}\setcounter{footnote}{0}%
de ere absolut ueensartede Principer.\footnote{R. Nielsen. Philosophisk Propædeutik. Kbhvn. 1857. S. 77.} Hvorledes denne
Forening skal tænkes, angives paa følgende Maade:
% The quotation pp. 11--12 is set as an indented block in the print.
\begin{quote}
„At Tro og Viden lade sig forene i een Bevidsthed,
ikke \emph{endskjøndt}, men \emph{fordi} de ere absolut
ueensartede Principer, vil da med andre Ord sige,
at de kun lade sig forene i \emph{et absolut Mysterium}.
Det absolute Mysterium er paa een Gang
exclusiv Grændse og negativ Midte. Som exclusiv
Grændse sætter Mysteriet Skilsmisse imellem Tro
og Viden som imellem to Regioner. Hver Region
har sit Ubetingede, sit Midtpunkt: Troens Midtpunkt
er Friheden, den almægtige, hellige Villie;
Videns Midtpunkt er Nødvendigheden, den rationelle,
i uforanderlige Love sig aabenbarende Nødvendighed.
Naar nu den Troende seer mod Troens,
den Vidende mod Videns Midtpunkt, ere Extremerne
spændte til det Yderste: hvad der er sandt i
Troen, at nemlig det Ubetingede er Miraklet, er
usandt i Viden; hvad der er sandt i Viden, at
nemlig det Ubetingede maa søges i Fornuftnødvendigheden,
er usandt i Troen. Som negativ Midte
forener Mysteriet derimod de adskilte Extremer og
viser sig som Regionernes Gjennemgangsmedium:
gjennem Mysteriet gaaer Veien fra Tro til Viden,
fra Viden til Tro. Idet den Troende og den Vidende
begge paa een Gang rette deres Blik mod
Mysteriet, udlignes Striden, og Extremerne enes.
Seet fra Videns Standpunkt er nu Miraklet mu%
% --- p. 12 ---
% (PDF 17)
\opage{12}\setcounter{footnote}{0}%
ligt. Thi vel er det umuligt, at der nogensinde
skulde kunne skee et vilkaarligt Brud paa
de evige med Fornuftnødvendigheden identiske Naturlove;
men med Lovenes Uforanderlighed kan
jo, naar de i Mysteriet skjulte Mellembestemmelser
underforstaaes, den Aabenbarelse af Villiens,
Viisdommens og Kjærlighedens Magt, hvori
den Troende seer Miraklet, meget vel tænkes
at bestaae. Seet fra Troens Standpunkt kan Fornuftnødvendigheden
med de uforanderlige Naturlove
forenes med Forudsætningen af Miraklet. Thi vel
er det umuligt, at Nødvendigheden skulde kunne
hemme Friheden, og Naturens Love indskrænke Guds
ubetingede Villie; men naar de i Mysteriet skjulte
Mellembestemmelser underforstaaes, lader det sig jo
tænke, at Naturlovene, langt fra at brydes, tvertimod
i allerhøieste Forstand opfyldes, hver Gang
en guddommelig Villiesact fuldbyrdes“.\footnote{R. Nielsen. Om „Den gode Villie“ som Magt i Videnskaben. Kbhvn. 1867. S. 33.}
\end{quote}

Denne Udtalelse kan med Rette ansees for det Centrale
og Nye i Professor Nielsens Skrift: Om „Den
gode Villie“ som Magt i Videnskaben. Han har her
for første Gang indladt sig paa en ligefrem Besvarelse
af det omtvistede Problem, hvorledes Tro og Viden,
netop \emph{fordi} de ere \textbf{absolut} \emph{ueensartede} Principer,
uden Modsigelse lade sig forene i een og samme Bevidsthed.
% PRINTED AS IS: „Interesseu“ (last letter a turned n, read at 300 ppi: stems joined at the foot, as u)
Det Begreb, hvorom Interesseu her samler
sig, er „\emph{Mysteriet}“, som danner Grændsen mellem Tro
% --- p. 13 ---
% (PDF 18)
\opage{13}\setcounter{footnote}{0}%
og Viden, og som faaer sin Betydning netop ved at
være \emph{Grændsebestemmelse}.

Om Grændsen\footnote{Om Behandlingen af Grændsebegrebet see følgende Skrifter af Professor Nielsen:\\
Philosophie og Mathematik. Kbhvn. 1857. S. 19 flg.\\
Mathematik og Dialektik. Kbhvn. 1859. S. 28---38 og S. 56---76, navnlig S. 70 og 76.\\
% PRINTED AS IS: „Kbhvn, 1865.“ with a comma (elsewhere „Kbhvn.“)
Forelæsninger over „Philosophisk Propædeutik“ (fra Universitetsaaret 1860---61, o. s. v.). Andet Cursus. Kbhvn, 1865. S. 53 flg.\\
At man paa de anførte Steder kun træffer Grændsebegrebet i Sammenhæng med mathematiske Udviklinger betyder her Intet, da vi kun benytte de Grundbestemmelser i Grændsen, som høre hjemme i dens Metaphysik, og som derfor gjælde \emph{overalt}, hvor der er Tale om Grændse.} lærer Professor Nielsen i det Metaphysiske,
at den er af dialektisk Natur; den er paa
eengang Eet med og forskjellig fra det, den begrændser,
og er forsaavidt paa eengang respektive \emph{positiv} og \emph{negativ}.
„I Grændsen“, hedder det, „reflecteres det Begrændsede.
At tænke Grændsen uden at underforstaae et
derved Begrændset, er ligesaa umuligt, som at fatte Betydningen
af en Negtelse uden at tænke paa det, som derved
bliver negtet. . . . Men naar Grændsen nu ikke destomindre
adskiller sig fra det Begrændsede som Negationen fra
det Positive, naar denne Adskillelse er i lige Grad gyldig
% PRINTED AS IS: stray point („.saa“), 127 px vs full stop F=175 px, sunk just below the baseline, immediately before „saa“.
for Tanken og Anskuelsen: .saa indeholdes heri den Opgave
at følge det, som skal begrændses, til Grændsen“.\footnote{R. Nielsen. Forelæsn. over „Philos. Propæd.“ Andet Cursus. Kbhvn. 1865. S. 56.}

Det er denne Opfordring, vi følge, idet vi vende
% --- p. 14 ---
% (PDF 19)
\opage{14}\setcounter{footnote}{0}%
os mod Mysteriet som Grændse mellem Tro og Viden. Det
absolute Mysterium er paa eengang \emph{exklusiv Grændse}
og \emph{negativ Midte}, hvilket med andre Ord vil sige, at
% NOTE: speck („paa begge“), 50 px vs full stop F=161 px, at mid x-height just before „begge“; not transcribed.
Grændsen paa \emph{begge} Maader er bestemt som \emph{negativ},
som den, der er af modsat Natur af det, den begrændser.
Dette udtrykkes for det Første ved Ordet „exklusiv“, men
det ligger ogsaa i Udtrykket „negativ Midte.“ Naar nu
Mysteriet som „negativ Midte“ skal \emph{forene} Tro og
Viden, da er dette ligefrem mod Forfatterens egen
Sprogbrug. Vi see imidlertid bort fra Sprogbrugen og
vende os mod de mere reelle Bestemmelser i Grændsen.
Forsaavidt Mysteriet sætter Grændse mellem Tro og
Viden, er her, ligeledes ifølge Professor Nielsens metaphysiske
Udviklinger, en dobbelt Opgave at løse: Som
\emph{Grændse for Viden} maa Mysteriet nemlig indeholde
noget fra Viden Forskjelligt, ellers kunde det jo ikke
kaldes Grændse, nemlig \emph{negativ} Grændse. At paavise
denne Forskjel volder ingen Vanskelighed, da den forstaaer
sig af sig selv. Men nu skal Mysteriet ogsaa
have noget med Viden Fælles; thi uden dette kunde det
heller ikke være Grændse, nemlig \emph{positiv} Grændse. Men
om et saadant Fællesskab mangler der enhver Oplysning,
og dog er just det \emph{positive} Element i Grændsen
udtrykkelig Betingelse for, at den kan være Gjennemgangsmedium.
Forsaavidt Mysteriet skal være \emph{Grændse
for Troen}, bliver Opgaven ligeledes dobbelt: for det
Første maa man see det for Mysteriet og Troen Fælles,
hvad der ikke er synderlig vanskeligt; men man skal tillige
erkjende deres indbyrdes Forskjel, det \emph{negative}
Grændsemoment i Mysteriet, og dette er meget vanske%
% --- p. 15 ---
% (PDF 20)
\opage{15}\setcounter{footnote}{0}%
ligt, fordi der ingen Forskjel er; Mysteriet er nemlig
ikke Grændse for Troen, men Troens eget Væsen, Troen
selv. Det vil saaledes sees, at Forsøget med at stille
Mysteriet som Grændse lider af tvende Skrøbeligheder:
I Forholdet til \emph{Viden} mangler Grændsens \emph{positive}
Moment, det vil sige, Mysteriet kan ikke \emph{forbinde} Tro
og Viden \emph{og er saaledes ikke Grændse}; i Forholdet
til \emph{Troen} mangler Grændsens \emph{negative} Element,
det vil sige, Mysteriet kan ikke \emph{adskilles fra} Troen
\emph{og er saaledes heller ikke Grændse}.

Vistnok er der Intet tilhinder for i al Almindelighed
at sætte et Mysterium som Grændse for Viden;
thi der gives ingen menneskelig Viden, uden at den jo
fører ud i den mørke Nat, hvad enten man saa gaaer
tilbage i Kausaliteternes Række og spørger om Oprindelsen,
eller man gaaer fremad fra Konsekvents til Konsekvents.
Men hvad der ikke gaaer an, er \emph{at hæve
Mysteriet til Kategori}, af Mysteriet som Grændsebestemmelse
at ville uddrage \emph{videnskabelige Konsekventser}.
I dette Tilfælde melder Kritiken sig af sig selv,
og Resultatet falder naturligviis ikke ud til Fordeel for
Mysteriet; thi om et \emph{Mysterium} kan der kun \emph{vides},
% „at“ is letterspaced (measured inter-letter gap 18 px against 4 px in an unspaced „at“ on this page); „det“ is not
\emph{at} det er et Mysterium, ellers Intet. Hvorledes Mysteriet
er hævet til Kategori, viser sig i den første Hovedsætning:

„Som exclusiv Grændse sætter Mysteriet Skilsmisse
imellem Tro og Viden som imellem to Regioner.“

Om Tro og Viden hedder det iforveien, at de ere
\textbf{absolut} \emph{ueensartede} Principer; men \textbf{absolut} \emph{ueensartede}
Principer kunne ikke begrændse hinanden, først
% --- p. 16 ---
% (PDF 21)
\opage{16}\setcounter{footnote}{0}%
fordi de Intet have tilfælles, dernæst fordi de \emph{ikke
engang ere} „\emph{forskjellige}“; thi for at være „forskjellige“
maatte de overhovedet have Noget med hinanden
at gjøre: Forskjel forudsætter Lighed, --- og dog skulle
de have en fælles Grændse i Mysteriet, altsaa dog begrændse
hinanden. Herved bliver Selvmodsigelsen eklatant;
skulde dette endnu ikke ansees for tilstrækkeligt, saa
ville i ethvert Tilfælde det citerede Skrifts egne Ord
levere det Manglende:

„Men Tro og Viden kunne --- paa Grund af
% „\&c.“ in the footnote renders the Fraktur ꝛc. abbreviation (read „rc.“ by the text layer); so also p. 19n
Mysteriet --- umulig indskrænke hinanden“.\footnote{R. Nielsen. Om „Den gode Villie“ \&c. S. 35.}

Analysis: Grændse betegner Indskrænkning; hvad der
indskrænker, begrændser. Naar nu Mysteriet opstilles
som Grændse, saa lyder Sætningen opløst saaledes:
„Men Tro og Viden kunne --- paa Grund af den
Bestemmelse, hvorved de \emph{udtrykkelig} indskrænke hinanden
--- \emph{umulig} indskrænke hinanden!“

Det hedder fremdeles i den ovenfor anførte Udtalelse:

„Hver Region har sit Ubetingede, sit Midtpunkt:
Troens Midtpunkt er Friheden, den almægtige, hellige
Villie; Videns Midtpunkt er Nødvendigheden, den rationelle,
i uforanderlige Love sig aabenbarende Nødvendighed.“
Hertil svarer i samme Skrift S. 35: „Vel
er der i Troen en Grændse for Viden, forsaavidt det,
der kan troes, ikke kan vides; men denne Grændse er
Eet med Videns egen iboende Grændse. Med sin iboende
Grændse, sin indre Begrændsning, er Viden lige%
% --- p. 17 ---
% (PDF 22)
\opage{17}\setcounter{footnote}{0}%
overfor Troen aldeles uindskrænket, saa uindskrænket, som
om der slet ingen Tro var.“

Hvorledes gaaer det nu med Grændsen? Ja
nu ere vi komne til et Punkt, hvor Tro og Viden
\emph{slet ikke begrændse hinanden}, hvor altsaa
Mysteriet bliver ganske overflødigt: der siges nemlig,
at Videns Grændse er iboende, altsaa er det
ikke Mysteriet, langt mindre Troen (som tænkes at ligge
paa den anden Side af Mysteriet), der danner Grændsen
for Viden. Vistnok er der et Tilløb til en Dialektik,
naar der siges, at der i Troen dog er en Grændse for
Viden, forsaavidt det, der kan troes, ikke kan vides; men
Dialektiken er igjen umuliggjort ved det \textbf{Absolut}-\emph{Ueensartede},
som lyder gjennem den Udtalelse, at Viden ligeoverfor
% NOTE: stray tick after the comma („uindskrænket, som“), comma-shaped, 12x23 px, 129 px vs comma 363 px on the page; not a sort, not transcribed.
Troen er aldeles uindskrænket, som om der slet
ingen Tro var. Naar enhver af Sphærerne har sit
eget \emph{ubetingede} Midtpunkt, naar Viden ingensinde er
istand til at angribe Troen eller omvendt, saa kan der
heller aldrig blive Tale om Grændse, hverken paa endelig
eller paa uendelig Maade. Med andre Ord: Læren
om to \textbf{absolut} \emph{ueensartede} Principer, som paa
nogensomhelst Maade skulde kunne begrændse hinanden,
staaer i den meest skjærende Modsigelse med Professor
Nielsens egen Lære om Grændsen i det Metaphysiske.

„Naar nu den Troende seer mod Troens, den Vidende
mod Videns Midtpunkt, ere Extremerne spændte
til det Yderste: hvad der er sandt i Troen, at nemlig
det Ubetingede er Miraklet, er usandt i Viden; hvad der
er sandt i Viden, at nemlig det Ubetingede maa søges i
Fornuftnødvendigheden, er usandt i Troen.“

% --- p. 18 ---
% (PDF 23)
\opage{18}\setcounter{footnote}{0}%
Vi svare: Dersom det, der er \emph{sandt} i Viden,
ligefrem og uden videre Indskrænkning er \emph{usandt} i
Troen og omvendt, saa kan den samme Bevidsthed umuligt
optage begge Bestemmelser i sig; thi selv om det
Sande i Viden og det Sande i Troen bestemmes paa
en nok saa forskjellig Maade, \emph{saa maa Sandhedens
Princip dog blive det samme i begge Sandheder}.

Idet vi nu gaae over til den anden Hovedsætning,
tiltage Vanskelighederne saaledes, at Modsigelsen mellem
Religionsphilosophien og Metaphysiken bliver aldeles slaaende.
Den anden Hovedsætning lyder saaledes:

% NOTE: the same sentence quoted on p. 11 reads „de adskilte Extremer“; here „Elementer“ as printed
„Som negativ Midte forener Mysteriet derimod de
adskilte Elementer og viser sig som Regionernes Gjennemgangsmedium:
gjennem Mysteriet gaaer Veien fra
Tro til Viden, fra Viden til Tro.“

Vi have allerede paaviist, at den negative Grændse,
ifølge Professor Nielsens egen Metaphysik, netop \emph{adskiller}
de Begrændsede, medens den her skal \emph{forbinde} dem;
men ikke nok hermed, den skal forbinde tvende \textbf{absolut}
\emph{ueensartede} Sphærer, af hvilke enhver, ifølge sin
egen iboende Grændse, er saa uindskrænket, som om den
anden slet ikke var! . . . . og dog er dette for Intet at
regne mod det, som følger:

„Idet den Troende og den Vidende begge paa een
Gang rette deres Blik mod Mysteriet, udlignes Striden,
og Extremerne enes. Seet fra \emph{Videns} Standpunkt, er
nu \emph{Miraklet muligt}. \textbf{Thi} vel er det umuligt, at
der nogensinde skulde kunne skee et vilkaarligt Brud paa
de evige med Fornuftnødvendigheden identiske Naturlove;
men med Lovenes Uforanderlighed kan jo, \emph{naar de i
% --- p. 19 ---
% (PDF 24)
\opage{19}\setcounter{footnote}{0}%
Mysteriet skjulte Mellembestemmelser underforstaaes},
den Aabenbarelse af Villiens, Viisdommens
% NOTE: speck („Miraklet, meget“), 41 px vs full stop F=172 px, at mid x-height after the comma; not transcribed.
og Kjærlighedens Magt, hvori den Troende seer Miraklet,
meget vel \emph{tænkes} at bestaae.“

For det Første forbliver man ganske uvidende om,
hvilke de i Mysteriet skjulte Mellembestemmelser vel kunne
være, dernæst om, hvorledes man skal bære sig ad med
at underforstaae dem. Iøvrigt er dette Mysterium af et
mærkeligt \textit{genus}; lærer man nemlig \emph{ikke} de skjulte
Mellembestemmelser at kjende, saa er det jo umuligt at
underforstaae dem, da man saa ikke veed, hvad man skal
underforstaae, eller kan underforstaae alt Muligt; der
% NOTE: „Tro,“ -- above and right of the comma a flat sliver, 69 px vs full stop F=172 px, 39 px above the comma head (the semicolon points of „genus;“ and „Muligt;“ are 150-165 px, 22-24 px above theirs); not a semicolon, not transcribed.
skeer da ingen Bevægelse fra Viden henimod Tro, og
Mysteriet bliver uanvendeligt. \emph{Lærer} man dem at
kjende, saa er Mysteriet ikke længere Mysterium, og saa
bliver det ligeledes uanvendeligt. Og endda skal denne
Proces finde Sted fra \emph{Videns} Standpunkt\footnote{I Skriftet: Om „Den gode Villie“ \&c. hedder det S. 89---90: „For den troende Bevidsthed er Grundmiraklets Modsigelse ganske vist løst; men ikke derved, at den selv \emph{tænker} Løsningen, nei derved, at den \emph{troer}, at Løsningen er tænkt i Mysteriet.“ Overfor denne Udtalelse suspendere vi enhver Dom, da vi ikke alene ved „Den gode Villie“, men selv ved Hjælp af den bedste Villie ikke formaae at fatte Meningen. Hvorledes kan Noget \emph{tænkes i et Mysterium}? Af \emph{hvem} er det tænkt? Og hvad i al Verden betyder det for en religiøs Bevidsthed, at den \emph{troer}, at en Løsning er \emph{tænkt} i et Mysterium?}.

„Seet fra Troens Standpunkt kan \emph{Fornuftnødvendigheden}
med de uforanderlige Naturlove forenes
% ===== batch b03_pp20-29 =====
% --- p. 20 ---
% (PDF 25)
\opage{20}\setcounter{footnote}{0}%
med Forudsætningen af \emph{Miraklet}. \textbf{Thi} vel er det
umuligt, at Nødvendigheden skulde kunne hemme Friheden,
og Naturens Love indskrænke Guds ubetingede
Villie; \emph{men naar de i Mysteriet skjulte Mellembestemmelser
underforstaaes}, lader det sig
jo \emph{tænke}, at Naturlovene, langt fra at brydes, tvertimod
i allerhøieste Forstand opfyldes, hver Gang en
guddommelig Villiesact fuldbyrdes.“

Herom gjælder selvfølgelig, hvad vi ovenfor paaviste
om Mysteriet, og vi skulle kun tilføie: „Naar Tro
og Viden, hvad det absolute Mysterium absolut forhindrer,
hverken kunne gavne eller skade hinanden“, hvad
Magt ligger der da paa, at de blive forenede?

Selvmodsigelsen lod sig videre forfølge gjennem „Approximationen“,
„Bevægelsen“ o. s. v., ja udstrække til
hele „Grundideernes Logik.“ Disse Undersøgelser vilde
imidlertid falde udenfor den her lagte Plan; vi gaae
derfor over til at betragte Religionsphilosophien i Forhold
til Videnskabslæren.

% Section head, centred bold Fraktur, set mid-page on p. 20, with a
% short rule above it and a shorter rule below it.
\begin{center}\textbf{\textit{II.} Videnskabslæren.}\end{center}

I sin philosophiske Propædeutik (S. 74) lærer
Professor Nielsen, at Liv og Videnskab, paa Grund af
begges inderlige Sammenhæng med Kultur og Humanitet,
ikke kunne forholde sig udvortes og ligegyldig til hin%
% --- p. 21 ---
% (PDF 26)
\opage{21}\setcounter{footnote}{0}anden; men at de desuagtet ere selvstændige ifølge en
principiel Adskillelse af det Subjektive og det Objektive.
Mod det første af disse Synspunkter kan der neppe
reises nogen Indvending; med Hensyn til det sidste maa
man derimod være yderst forsigtig. Det kommer nemlig
an paa, hvorledes Ordet „selvstændig“ skal forstaaes.
At „subjective Stemninger: Tilbøielighed, Utilbøielighed,
Frygt, Haab o. desl. Intet have at betyde i Videnskaben“,
kan vel tildeels indrømmes, forsaavidt hermed
er sagt, at Vilkaarligheden Intet har at betyde i Videnskaben.
Men drager man nu uden videre den Slutning,
at Personligheden er udelukket af Videnskaben, og hævder
de „upersonlige Sandheder“ som Kategori, saa begaaer
man en Feil ved at forvexle \emph{Stemning} og \emph{Personlighed};
med andre Ord, man gjør en Side af Personligheden
til den hele Personlighed. Ved at charakterisere
f. Ex. Geometriens Sandheder som „upersonlige“, Ethikens
som „personlige“ overseer man, at Sandhedens
Væsen bestemmes, ikke alene efter Erkjendelsens \emph{Gjenstand}
(Geometri og Ethik), men ogsaa efter den erkjendende
\emph{Bevidstheds} Natur. At lade Stemningen være
det Dominerende i den videnskabelige Erkjendelse er naturligviis
meningsløst, fordi Forholdet mellem Bevidstheden
og dens Gjenstand derved vilde blive reent \emph{tilfældigt}.
Under denne Forudsætning gjælder det ganske vist, at „det
ikke er subjective Følelser, ikke fromme Ønsker, men
Indsigt i Tingenes Væsen, der afleder Lynstraalen, og
sætter Dampvognen i Bevægelse.“ Men denne Indsigt i
Tingenes Væsen beroer i sin dybeste Grund paa en
Congruents mellem det Aprioriske (Fornuftnødvendige) i
% --- p. 22 ---
% (PDF 27)
\opage{22}\setcounter{footnote}{0}%
Tingenes Væsen og det Aprioriske i Personligheden, det
Aprioriske i Villien ikke mindre end det Aprioriske i
Tænkningen. Til nærmere Forstaaelse indføre vi den af
% The final e of „ofte“ and of „fremhævede“ prints without its upper
% stroke (worn sort, looks like c); read e.
Professor Nielsen saa ofte fremhævede Forskjel mellem
Forgrund, Mellemgrund og Baggrund\footnote{Jfr. R. Nielsen. Forelæsninger over „Philosophisk Propædeutik“.
Første Cursus. Kbhvn. 1865. Forord. S. IV.}. Hvad der i
% The m of „Sammenhæng“ (twice in the next two lines) is broken.
een Sammenhæng ifølge Sagens Natur maa træde i
Forgrunden, kan i en anden Sammenhæng, ligeledes
ifølge Sagens Natur, enten have Betydning som Bestemmelse
i Mellemgrunden eller endog træde tilbage i
Skyggen. For den \emph{videnskabelige} Erkjendelses Vedkommende
er det uimodsigeligt, at Tænkningen træder i
Forgrunden og giver Vexelforholdet mellem Bevidstheden
og Gjenstanden sit Præg; men derfor ophører Bevidstheden
ikke at være Villie, saalidt som den efter Behag kan
lægge Villien fra sig og atter optage den. Saavist som
Personlighedens Eenhed aldrig kan brydes, saavist er Villien
ogsaa medvirksom i en \emph{hvilkensomhelst} videnskabelig
Erkjendelse. Derom vidner allerede den \emph{Tilfredsstillelse},
som ledsager enhver Erkjendelse af Sandheden; i
al Tilfredsstillelse er der et Villiesmoment, fordi Tilfredsstillelsen
\emph{begjæres}. Og nu \emph{Opmærksomheden};
paa intet Punkt i Erkjendelsen bliver Villiens Nærværelse
saa aabenbar som her. Der er Intet i den rene
Tænkning, som ifølge sin Natur \emph{vækker} eller \emph{forøger}
Interessen eller \emph{foranlediger} Tænkningen til at gaae
videre; uden Villiesmomentet vilde Tænkningen optage
sit Indhold som Speilet optager sit Billede, --- eller
% --- p. 23 ---
% (PDF 28)
\opage{23}\setcounter{footnote}{0}%
rettere, Tænkningen vilde slet ikke existere. Men vi kunne
gaae videre: selv Stemningerne høre med i de videnskabelige
Sandheder, ja endog --- er det nu ikke forfærdeligt
--- i selve de geometriske Sandheder! \emph{Interessen}
for de videnskabelige Sandheder, den \emph{Spænding}, hvormed
Problemernes Løsning imødesees, den \emph{Tilfredsstillelse},
som føles, naar Maalet er naaet, viser hen
til de fineste psychologiske Overgange, hvor Villie og Stemning
bevæge sig paa Grændsen af Bevidsthedens Mellemgrund.
Hvad man saa ellers kan paavise om Sandheden,
den ophører dog ikke at være \emph{menneskelig}
Sandhed. Man kan i Philosophien stille sig paa hvilket
Standpunkt, man vil, man kommer her ligesaalidt udenfor
det Psychologiske, som man i Astronomien kommer
% The M of the first „Maanen“ is damaged (upper part broken off).
op paa Maanen, fordi man „stiller sig i Maanen“.
Allerede af denne Grund negte vi derfor, at der gives
reent upersonlige Sandheder. Man komme nu ikke med
den fortærskede Indvending, om da f. Ex. de geometriske
Sandheder ere afhængige af vore Stemninger og kunne
forandres med dem; for ikke at tale om, at man derved
overseer, hvad der nylig er sagt om Forgrund, Mellemgrund
og Baggrund, vilde man tillige røbe en mærkelig
Mangel paa Kundskab om Sandhedens \emph{Natur}. Sandheden
existerer ikke objektiv som en Steen, men \emph{er} kun,
idet den \emph{erkjendes}; at tage Sandheden alene objektivt
existerende vilde være ligesaa naivt som at forvexle den
Bog, der handler om Geometrien, med Geometrien selv.
% UNSURE: whether the single letter „i“ in „er i og under“ is spaced
% cannot be measured; „er“ and „under“ are spaced, „og“ is not.
Idet Sandheden kun \emph{er} i og \emph{under} sin Vorden, bevarer
den Charakteren af sin Oprindelse og er altid subjektiv-objektiv.
Udaf denne Dobbelthed kan den ikke ud%
% --- p. 24 ---
% (PDF 29)
\opage{24}\setcounter{footnote}{0}løses. Naar det i Propædeutiken (S. 75) hedder: „Vistnok
er Videnskab umulig uden Subjectivitet; men Subjectiviteten
% NOTE: stray grave-shaped tick before the second „her“ („føles, her tænkes“), 72 px vs full stop F=169 px, at x-height top; not transcribed.
er tjenende. Her sandses, her føles, her tænkes;
men det i Sandsning og Tænkning blot Subjective maa
overalt udrenses som det Tilfældige, det Forstyrrende, det
Usande“, --- da kan dette indrømmes, dog med den
Indskrænkning, at det \emph{blot} Subjektive ikke sættes identisk
med \emph{alt} det Subjektive. Men at dette er Meningen
sees af de følgende Ord: „\emph{hver} Forudsætning er \emph{objectiv}.“
Er da Bevidstheden ikke en Forudsætning? At
sige: „\emph{hver} Forudsætning er \emph{objektiv}“ i samme Aandedræt,
som der siges: „Vistnok er Videnskab umulig uden
Subjectivitet“, er en Modsigelse. At sige: „\emph{hver} Forudsætning
er \emph{objectiv}“ og saa begynde „Grundideernes
Logik“ med den Bemærkning, at man kun kan lære, hvad
Tænkning er, ved at tænke; og at Tænkningen ikke alene
er Tankens Begyndelse, men ogsaa dens Fortsættelse og
Fuldendelse, --- er atter en Modsigelse.

% The m of „Bestemmelsen“ and the U of „Undervurdering“ below are broken.
I Bestemmelsen af den videnskabelige Sandheds
\emph{Natur} finder her i det Hele en dobbelt Miskjendelse
Sted: en Undervurdering af det Subjektive og en Overvurdering
af det Objektive.

Først altsaa \emph{Undervurderingen af det Subjektive}.
Er Subjektiviteten med som „tjenende“, saa
er den jo dog med, og giver altsaa Sandheden sit
Præg. Den Paastand, at \emph{hver} Forudsætning er
objektiv, overseer, at den erkjendende Bevidsthed er
en ligesaa væsentlig Forudsætning som Erkjendelsens
Objekt. Skulde endelig den Omstændighed, at Geometriens
Sandheder ikke kunne forandres ved Villie,
Stemninger o. s. v., være Beviis for, at disse slet ikke
% --- p. 25 ---
% (PDF 30)
\opage{25}\setcounter{footnote}{0}%
ere med, saa kunde man jo ogsaa paastaae, at Tænkningen
ikke var med, thi Tænkningen kan jo heller ikke
forandre Geometriens Sandheder. Fordi Villie og
Stemning \emph{kunne} virke forstyrrende, deraf følger ikke, at
de \emph{maae} gjøre det.

Undervurderingen af det Subjektive beroer altsaa
derpaa, at Personligheden (i den videnskabelige Erkjendelse)
identificeres med de subjektive Stemninger, og
atter med denne Indskrænkning kun opfattes som „det
Tilfældige, Forstyrrende og Usande.“ Subjektivitetens
Nærværelse opfattes som noget Abnormt, som Noget,
der \emph{ikke bør være}; vi opfatte dens Nærværelse som
naturnødvendig og normal, som Noget, der nødvendigviis
\emph{maa være}.

Den anden Miskjendelse: \emph{Overvurderingen af
det Objektive} staaer i nøie Sammenhæng med den
foregaaende. Det er det i de videnskabelige Sandheder
\emph{Uimodsigelige}, hvoraf man har ladet sig blænde og
forlede til at bestemme dem som upersonlige og reent objektive.
Men det Uimodsigelige (Aprioriske, Fornuftnødvendige)
er ogsaa tilstede i den hele Personlighed. I
Begribningens Akt --- og kun i denne Akt existerer
Sandheden --- coinciderer det „Uimodsigelige“ i Objektet
med det „Uimodsigelige“ i Bevidstheden, i Villien
saavelsom i Tænkningen, \emph{ja i den hele, udeelte Personlighed}.
Det er ved at fastholde Resultatet (f. Ex.
den geometriske Sætning) værende, og fordi Villien o. s. v.
ikke \emph{umiddelbart} efterlader noget Spor af sin Virksomhed,
ikke findes, saa at sige, grovkornet og haandgribelig, at
man kan komme til at udelukke dens Medvirken; men det er
% --- p. 26 ---
% (PDF 31)
\opage{26}\setcounter{footnote}{0}%
jo ofte Tilfældet, at Analysen ikke kan nøies med Resultatet
alene, men maa søge sin Oplysning netop i
de Processer, hvorigjennem Resultatet er naaet. Og
det er ligeledes kun ved at fastholde Resultatet abstrakt,
at det for en „populær“ Bevidsthed kan tage sig ud
som Galmands-Snak, at Villie og Stemning ere med
i Geometriens Sandheder. Da det imidlertid ikke kommer
an paa at underholde den „populære“ Bevidsthed,
men paa at undersøge et videnskabeligt Problem, saa ville
vi til yderligere Forstaaelse til den foregaaende Argumentation
% NOTE: stray single bar after „videnskabelige“ at the line end, 18x6 px, 90 px (no single-stroke hyphen sort in this fount); not transcribed.
endnu føie et Billede: I den reent videnskabelige
Erkjendelse er Tænkningen den fremherskende Faktor,
Villie, Stemning o. s. v. „symphoniske Elementer.“
Den Umusikalske hører i en Symphoni kun den fremtrædende
Melodi, hvori de forskjellige Stemmer smelte
sammen i en Eenhed; det uddannede Øre derimod hører
paa eengang den symphoniske Eenhed og de enkelte Stemmer,
der afvexlende træde frem og tilbage. Paa lignende
Maade i Erkjendelsen: her dominerer Tænkningen, medens
Villie og Stemninger medlyde. Som nu en Tilhører
mangen Gang ikke ahner, hvormange Instrumenter
han i Virkeligheden har hørt, saaledes kan Villiens og
Stemningernes Virksomhed i den videnskabelige Erkjendelse
ogsaa undgaae Opmærksomheden. Naar endelig
dertil føies, at de ofte misbruges og anvendes vilkaarligt,
lader det sig forstaae, at man er kommet til reent
at ville afskaffe dem; men \textit{abusus non tollit usum}.

Vi fastholde altsaa som Resultat:

1) I Erkjendelsen af enhver videnskabelig Sandhed
% --- p. 27 ---
% (PDF 32)
% Items 1)--5) are set with a hanging indent (numeral at the paragraph
% indent on p. 26, flush left on p. 27; run-over lines indented).
\opage{27}\setcounter{footnote}{0}%
ere ikke blot Sandsning og Tænkning, men den
hele, udeelte Personlighed \emph{tilstede}.

2) Den er tilstede med \emph{Nødvendighed}.

% The final s of „Personlighedens“ is a damaged sort.
3) Ved Personlighedens Nærværelse \emph{kan} Sandheden \emph{ikke}
krænkes, fordi her kun er Tale om det i Personligheden
(i Sandsning, Tænkning, Villie o. s. v.)
Aprioriske, det med Sandheden Adækvate, medens
al Vilkaarlighed er udelukket.

4) Enhver Sandhed er \emph{personlig}; den er subjektiv-objektiv,
skjøndt paa \emph{relativ} forskjellig Maade i
det Theoretiske og i det Praktiske.

5) I Erkjendelseslæren bliver det en Opgave at tilstræbe
Efterviisningen af \emph{alle} Personlighedens Momenter
i Erkjendelsen af \emph{enhver} Sandhed.

Efterat vi foreløbigt have gjort Rede for vor egen
Opfattelse af Sandhedens Forhold til Personligheden,
vende vi tilbage til Professor Nielsens nærmere Begrundelse
af de upersonlige Sandheder: „Kan en Bevidsthed
ikke halveres af Venskab og Mathematik, af
Kjærlighed og Chemie, saa kan den endnu mindre halveres
af Ethik og Astronomie, af Religion og Geologie.
Der gives mathematiske Venner, men der gives ingen
venskabelig Mathematik; der gives forelskede Chemikere,
men ingen chemisk Forelskelse, ingen forelsket Chemie.
Der kan gives ethiske og uethiske Astronomer, men Astronomien
er hverken ethisk eller uethisk. Videnskaben er
med andre Ord hverken ethisk eller uethisk, hverken
christelig eller uchristelig, hverken troende eller vantro:
% Footnote: „rc.“ is the Fraktur et-cetera sign, rendered \&c.
Kjærlighed er praktisk, men Videnskab er theoretisk“\footnote{R. Nielsen. Om „Den gode Villie“ \&c. S. 36.}.

% --- p. 28 ---
% (PDF 33)
\opage{28}\setcounter{footnote}{0}%
Ved første Øiekast er denne Argumentation slaaende,
man kunde fristes til at sige: bedaarende. Det gjælder
nu, om vor almindelige Undersøgelse af Personlighedens
Forhold til Sandheden kan bestaae sin Prøve i det konkrete
Tilfælde; thi det kan let skee, at en Argumentation
seer plausibel ud, saalænge den holdes i reen Almindelighed,
medens Manglerne komme tilsyne ved dens Anvendelse
i Virkeligheden; og det hænder ligesaa ofte, at
man først kan føle sig slaaet af den generelle Principudvikling,
men senere blive betænkelig og faae Skrupler,
naar det enkelte Tilfælde, udstyret med Billeder og stærke
Farver, lokker Phantasien paa Glatiis. „En chemisk
Forelskelse“ er unegtelig et komisk Begreb, egnet til at
illustrere, hvor meningsløse Resultater der vil fremkomme,
naar man indfører det Vilkaarlige i Videnskaben
og sammenblander Personligt og Upersonligt. Men see
vi lidt nærmere paa den „chemiske Forelskelse“, mon
saa den Omstændighed, at Forelskelsen Intet har at
gjøre i Chemien, berettiger til den Slutning, at Personligheden
ingen Betydning faaer i Erkjendelsen af
Chemiens Sandheder?

% The M of „Mathematik“, „Modsætningen“, „Maade“ and the m of
% „Komiske“, „Kæmpe“ on this page are broken sorts.
% „forelsket Chemi“: the quotation on p. 27 reads „forelsket Chemie“.
% Both as printed.
Det Bestikkende ved saadanne Exempler som „venskabelig
Mathematik“ og „forelsket Chemi“ ligger deri,
at Modsætningen er komisk. Det Komiske fremkommer
i Almindelighed --- og flere Bestemmelser behøve vi ikke
her --- naar Modsætninger, helst stærke Modsætninger,
umiddelbart forbindes paa vilkaarlig eller tilfældig
% footnote continues onto p. 29 (from „sluttende uvidenskabelig“; the
% word „Afsluttende“ is divided „Af-/sluttende“ across the page turn).
Maade\footnote{Det vil her være oplysende at eftersee Johannes Climacus. Afsluttende uvidenskabelig Efterskrift til de philosophiske Smuler.
Kbhvn. 1846. S. 394 flg.}. Naar en Kæmpe og en Dværg spadsere
% --- p. 29 ---
% (PDF 34)
\opage{29}\setcounter{footnote}{0}%
jevnsides, naar et høitideligt Foredrag afbrydes af
en triviel Bemærkning, føler Enhver det Komiske.
I det nævnte Exempel ere Modsætningerne „venskabelig“
og „Mathematik“ valgte vilkaarligt og tilfældigt;
de ere fremdeles valgte i de yderste Extremer,
det Første i Villiens (eller Følelsens), det Sidste i Videns
Sphære; endelig ere de forbundne ganske umiddelbart.
Naar nu Bevidstheden, lokket af det Komiske i
Sagen, ikke kan vende Opmærksomheden bort derfra,
kommer den let til at indrømme som almeengyldigt
Noget, som den, med mindre bestikkende Exempler for
% NOTE: „Omstændighed,“ -- above the comma a wedge-shaped mark, 145 px vs full stop F=170 px, wholly above the x-height and 44 px above the comma head (the „Sphære;“ point on this page sits at x-height top, 22 px above its comma); not a semicolon, not transcribed.
Øie, aldeles ikke vilde indrømme. Den Omstændighed,
at Venskabet, en enkelt Yttring af Personligheden, ikke
umiddelbart kan forbindes med Mathematiken, beviser Intet
om Personlighedens Forhold til de mathematiske Sandheder.
Det er en gammel Witz, at man, for at modbevise
den dialektiske Methodes almene Gyldighed, vælger
% NOTE: stray single bar after the closing quote of „Fastelavn“ („Fastelavn“ og), 19x8 px, 103 px, at mid x-height (no single-stroke hyphen sort in this fount); not transcribed.
Modsætninger som „Kjøbenhavn“ og „Fastelavn“ og
derefter udbeder sig deres „dialektiske Eenhed.“ Det
% PRINTED AS IS: stray point („os ·at“), 94 px vs full stop F=170 px, raised to mid x-height, immediately before „at“.
være nu langt fra os ·at sammenstille Professor Nielsens
Raisonnement med dette, som i overfladisk Tankeløshed
vanvittigt sammenblander Væsentligt og Uvæsentligt og
hjemfalder til det rene Vaas; men med dette Forbehold
bliver der dog altid den Lighed tilbage, at Modsætningerne
paa begge Steder vælges iflæng og hvile paa eensidige og
vilkaarlige Abstraktioner. Seer man paa Videns Indhold i
de forskjellige Videnskaber, saa lader dette sig ordne saaledes,
% ===== batch b04_pp30-38 =====
% --- p. 30 ---
% (PDF 35)
\opage{30}\setcounter{footnote}{0}%
at det danner en Skala, bestemt ved Forholdet til det
Personlige. Ved Skalaens Udgangspunkt begynde vi
med den Videnskab, hvis Gjenstand er Personligheden
selv, nemlig Ethiken; derfra komme vi gjennem Psychologi,
Anthropologi, Æsthetik, Historie o. s. v.
% PRINTED AS IS: stray point („hvor· Personligheden“), 188 px vs full stop F=146 px, raised to x-height top and above, right after „hvor“.
ned til Mathematiken, hvor· Personligheden tilsyneladende
ganske er forsvunden. Tager man nu et enkelt Led ud
af Rækken, navnlig et saadant, hvor Personlighedens
Nærværelse og hele Betydning først kan erkjendes gjennem
flere Mellemled (Mathematiken), og stiller det umiddelbart
sammen med en Villies- eller Følelsesbestemmelse,
hvor Personligheden er stærkt fremtrædende (Venskabet),
saa fremkommer der vistnok en overraskende Modsætning,
men man har med det Samme brudt Rækkens Kontinuitet,
og kun gjennem en omhyggelig Fastholden af
denne kan Personlighedens Nærværelse i alle de „upersonlige“
Sandheder erkjendes. Tager man altsaa det
komiske Element ud af Exemplet, saa betager man det
% UNSURE: „Braaden“ — the second a carries a circumflex-like mark (â), probably damage to the sort; transcribed plain a (the text layer reads „Branden“, which the image does not support)
Braaden, og Fortryllelsen hæves. Som illustrerende Pendant
til den „venskabelige Mathematik“, der gjør sin
Virkning ved en egen Anvendelse af Diskretionen og
Kontinuiteten, anføre vi et Indfald af en Kunstner, som
paa det første Blad i en Bog aftegnede Venus fra Milo,
paa det næste atter Venus fra Milo, men med en Afændring
saa ringe, at den kun opdages af et meget skarptseende
Øie. Paa denne Maade fortsættes der gjennem
smaae Ændringer, indtil man paa det sidste Blad finder
en Ourang-Outang. Man behøver nu blot at sammenholde
det første og det sidste Blad i Bogen for at forstaae
% --- p. 31 ---
% (PDF 36)
\opage{31}\setcounter{footnote}{0}%
Tilblivelsen af en Modsætning som „venskabelig Mathematik.“

Vi komme herfra til Spørgsmaalet om Videnskabslærens
Forhold til Problemet om Tro og Viden.
I det Logisk-Physiske have vi altsaa reent upersonlige
Sandheder; i Religionsphilosophien reent personlige Sandheder.
Men hvad have vi saa paa Overgangen mellem
begge? Professor Nielsen svarer: „Den vide Kløft
mellem Philosophie og Religion udfyldes af Ethiken. I
Ethiken mødes Livstankerne med de videnskabelige Grundtanker;
den rationelle Ethik staaer under Videnskabens,
den religiøse Ethik under Troens Indflydelse: Forskjellen
bestemmes ved den Accent, der lægges paa Subjectiviteten.
„Det er den videnskabelige Kritik ligesaa umuligt
at anerkjende Evangeliet om den Opstandne, som det er
Videnskaben umuligt at anerkjende Mythen om Herkules.“
Hvad kan da hindre Videnskaben i at erklære Evangeliet
om den Opstandne for en Mythe? Den religiøse Ethiks
Overlegenhed over den rationelle! Er der paa det rationelle
Standpunkt ingen Modsigelse i at gjøre Forskjel
imellem Physikens upersonlige og Ethikens personlige
Sandheder, saa er der paa Personlighedens Standpunkt
heller ingen Modsigelse i at fastholde Philosophiens
videnskabelige og Religionens overvidenskabelige Sandheder
i deres indbyrdes, for den i Timelighed existerende
Subjectivitet afgjørende, Modsætning“\footnote{R. Nielsen. Grundideernes Logik. I. Kbhvn. 1864. Indledning. S. XXVII.}.

% --- p. 32 ---
% (PDF 37)
Denne Udtalelse, der tillige skal være en Forkla\opage{32}\setcounter{footnote}{0}ring,
nemlig Forklaring af, hvorledes Ethiken skal tænkes
som \emph{Grændse}videnskab, giver Anledning til en Mængde
Spørgsmaal, som lades heelt ubesvarede, skjøndt Besvarelsen
af dem er det Eneste, der kan give Forklaringen
% NOTE: speck („nogen Værdi“), 34 px vs full stop F=150 px, just below the baseline after „nogen“; not transcribed.
nogen Værdi. Det, som i denne Sag interesserer, er
naturligviis allerførst Spørgsmaalet om Grændsen, d. v. s.
om Overgangen fra de reent upersonlige Sandheder til
de reent personlige. Først hedder det, at den vide Kløft
mellem Philosophi og Religion „udfyldes“ af Ethiken.
Men Bestemmelsen „udfyldes“ er kun Billede, ikke Begreb.
Fremdeles: Livstankerne „mødes“ med de videnskabelige
Grundtanker; men i Udtrykket „mødes“ kan der, naar
al Oplysning mangler, lægges aldeles modsatte Betydninger.
Dette Samme gjælder, naar den rationelle Ethik
siges at staae under Videnskabens „Indflydelse“; „at
staae under Indflydelse af“ er et vagt og taaget Udtryk,
som kan underlægges saamange Fortolkninger man vil;
for ikke at tale om, at rationel og religiøs Ethik opstilles
som bekjendte og givne i \emph{Indledningen} til en
„Grundideernes Logik“, der giver sig ud for \emph{principielt}
at \emph{begynde} Philosophien heelt forfra. Man kan
% footnote continues onto p. 33 and p. 34
fortsætte det ovenfor Citerede\footnote{Den sidste af de ovenfor citerede Sætninger fortjener særlig
Omtale paa Grund af den besynderlige Feilslutning, som
deri forekommer. Vi anføre den igjen: „Er der paa
det rationelle Standpunkt ingen Modsigelse i at gjøre
Forskjel imellem Physikens upersonlige og Ethikens personlige
Sandheder, saa er der paa Personlighedens Standpunkt
heller ingen Modsigelse i at fastholde Philosophiens
videnskabelige og Religionens overvidenskabelige Sandheder i
% (footnote text, continued at the foot of p. 33)
deres indbyrdes, for den i Timelighed existerende Subjectivitet
afgjørende, Modsætning.“ Her er Feilslutning Sætning
for Sætning, Led for Led, Ord for Ord. Nemlig:
\par 1) Det \emph{er}, som vi ovenfor have viist, paa det \emph{rationelle}
Standpunkt en Modsigelse at gjøre Forskjel mellem Physikens
upersonlige og Ethikens personlige Sandheder, naar
denne Forskjel sættes absolut.
\par 2) Der opstilles en Modsætning mellem det \emph{rationelle}
Standpunkt og \emph{Personlighedens} Standpunkt; er
Personligheden da ikke rationel? Man skulde troe det,
siden der paa det rationelle Standpunkt tales om Ethikens
\emph{personlige} Sandheder. Altsaa: \emph{Enten er} Personligheden
rationel, men i saa Fald kan man ikke tale
om \emph{overvidenskabelige} Sandheder paa Personlighedens
Standpunkt; \emph{eller} ogsaa er Personligheden \emph{ikke}
rationel, men saa kan man ikke tale om personlige Sandheder
paa det \emph{rationelle} Standpunkt; \emph{eller} endelig
falder Personligheden baade indenfor og udenfor det Rationelle,
men i saa Fald maatte man have fundet en
begrebsmæssig Bestemmelse heraf, da Slutningen jo ellers
bliver et Postulat.
\par 3) \emph{Almindelig} Feilslutning. Man kan ikke slutte saaledes:
Fordi der paa det \emph{rationelle} Standpunkt ingen
% UNSURE: the „i“ before „et“ is printed heavier than the type around it (bold sort?); transcribed plain, with only „et“ letterspaced
Modsigelse er i \emph{et} Forhold, derfor er der paa \emph{Personlighedens}
Standpunkt heller ingen Modsigelse i et
\emph{ganske andet} Forhold.
\par 4) \emph{Speciel} Feilslutning. Skulde Slutningen gjælde, maatte
de korrelate Modsætninger nøiagtigt svare til hinanden.
Men:
\par \textit{a)} Paa det rationelle Standpunkt forholder \emph{Physiken}
sig aldeles ikke til \emph{Ethiken}, som paa Personlighedens
Standpunkt \emph{Philosophien} til \emph{Religionen}.
% (footnote text, continued at the foot of p. 34)
\par \textit{b)} Modsætningen: Upersonlig---Personlig er en ganske anden
end Modsætningen: Videnskabelig---Overvidenskabelig;
thi:
\par $\alpha$) Man \emph{veed} jo for det Første slet ikke, om det Overvidenskabelige
overhovedet \emph{er}; man faaer altsaa en
Modsætning, hvori man kjender begge Led (nemlig:
Upersonlig---Personlig), og slutter derfra til en anden
Modsætning, om hvis ene Led man videnskabeligt
Intet veed.
\par $\beta$) Dernæst staaer i Modsætningen: Upersonlig---Personlig
det Upersonlige i et ganske andet Forhold til
det Personlige end i den tilsvarende Modsætning
det Videnskabelige til det Overvidenskabelige.
\par $\gamma$) Endelig har det sin Vanskelighed at forstaae, hvorledes
man tør gjøre en formeentlig stringent \emph{videnskabelig
Slutning} med Hensyn til nogetsomsomhest
% PRINTED AS IS: „nogetsom-/somhest“ (for „nogetsomhelst“): „som“ repeated across the line break, and no „l“ in „hest“
% UNSURE: „Over“ is set in bold sorts and not visibly letterspaced; the rest of the word is letterspaced in ordinary weight
\textbf{Over}\emph{videnskabeligt}.}, Linie for Linie; det er
den ene dogmatiske Sætning, som følger ovenpaa den
% --- p. 33 ---
% (PDF 38)
\opage{33}\setcounter{footnote}{0}%
anden, og som Enhver udenvidere kan benegte i samme
Øieblik den udtales; thi Begrundelse findes der ikke.
% --- p. 34 ---
% (PDF 39)
\opage{34}\setcounter{footnote}{0}%
Kun dette fastholde vi som vitterligt: Der paavises ingen
Overgang fra de „upersonlige“ Videnskaber til den rationelle
Ethik, ingen Overgang fra den rationelle til den
religiøse Ethik og endelig ingen Overgang fra den religiøse
Ethik til Religionsphilosophien. Og dog var det
om Alt dette man ønskede Oplysning, da Udtalelsen
ellers bliver en „fortrolig Meddelelse“, som vistnok kan
have sin Interesse --- kun ikke i Videnskaben. „Hvad vil
nu skee,“ spørge vi med Professor Nielsen, „naar en
philosophisk Tænker, en virkelig Logiker, engang faaer
isinde at revidere en saadan Fremstilling?“\footnote{R. Nielsen. Om „Den gode Villie“ \&c. S. 128---129.}

Vende vi os fra denne specielle Udtalelse og
betragte Sagen i Almindelighed, saa bliver Spørgsmaalet
om Overgangen fra Viden til Tro. Det er det
gamle Grændsespørgsmaal fra Metaphysiken, som igjen
% --- p. 35 ---
% (PDF 40)
\opage{35}\setcounter{footnote}{0}%
vender tilbage i Videnskabslæren. I „Grundideernes
Logik“ tages Udgangspunktet i „den rene Viden“; hvorledes
er det saa muligt nogensinde at komme ud over
Videns Kreds? Ved Viden selv kan det naturligviis
ikke skee, da „Viden kun kan bevirke Viden“, aldrig Tro.
Men maaskee lader det sig gjøre ved Hjælp af Villien;
thi om denne lærer Professor Nielsen, at den, uafhængigt
af Viden, sætter sig selv et nyt Udgangspunkt.
Dette gaaer imidlertid heller ikke an, da det jo netop
hedder, at Tanke aldrig kan forvandles til Villie, og da
tilmed Viden og Tro ere \textbf{absolut} \emph{ueensartede} Principer.
Ikke destomindre tales der om en Videnstheisme
og en Troestheisme\footnote{R. Nielsen. Om „Den gode Villie“ \&c. S. 135 flg.} med tilsvarende Guds\emph{erkjendelse},
hvori enhver af de to Theismer har sin egen
\emph{Erkjendelses}sphære. Skjøndt nu disse Sphærer ere
\textbf{absolut} \emph{ueensartede}, skulle dog Charakteermærkerne i
den ene „speile“ Charakteermærkerne i den anden. Men
% PRINTED AS IS: „Virkelig-/ligheden“ (for „Virkeligheden“): „lig“ repeated across the line break
hvorledes dette er muligt, eller hvorledes det i Virkeligligheden
skal gjennemføres, derom modtage vi ikke den
ringeste Oplysning. Istedetfor en streng Begrebsudvikling
møde vi igjen den rene Dogmatisme. Hvad betyder
„speile“ i denne Sammenhæng? Det er kun et Billede,
men her kommer det netop an paa Begrebet. At „Religionsphilosophien,
\emph{forsaavidt} den er en Videnskab, „naturligviis“
maa høre under Videns Kreds“, forstaaer sig af sig
selv; og at den, „\emph{forsaavidt} den har Religionen til Indhold,
maa orientere os i Troens Sphære“, er ogsaa vist; men
% --- p. 36 ---
% (PDF 41)
hvorledes den \emph{kan} gjøre begge Dele, hvorledes „Spei\opage{36}\setcounter{footnote}{0}lingen“
tager sig ud, hvorledes det samme Begreb genetisk
og \textit{in concreto} bliver til som Vidensbegreb og
som Troesbegreb, derom faae vi ikke det Mindste at
vide. „At de \textbf{absolut} \emph{ueensartede} Begreber i Grændsen
selv maae blive formelt eensartede“ er en eklatant Modsigelse,
fordi Bestemmelsen „absolut“ energisk og ubetinget
udelukker endog den blot formelle Parallelisme;
men selv afseet herfra, hvilken Betydning kan det have
enten for den videnskabelige eller for den religiøse Erkjendelse,
naar hos Forfatteren den formelle Eensartethed ikke
løber ud paa Andet end „at Princip jo \textit{in abstracto} altid er
% PRINTED AS IS: footnote „R Nielsen.“ without the full stop after R
Princip, Erkjendelse = Erkjendelse, Begreb = Begreb“?\footnote{R Nielsen. Om „Den gode Villie“ \&c. S. 136---137.}
Kan denne formelle Parallelisme berettige til den Antagelse,
at Vidensbegreberne „speile“ Troesbegreberne;
thi „Afspeilingen“ skulde jo dog være mere end et Spil
med Ord?

Vælger man sit Udgangspunkt i Viden, saa bliver man
indenfor Viden; vælger man sit Udgangspunkt i Troen, saa
bliver man indenfor Troen. I første Tilfælde vil man snart
opdage Videns Utilstrækkelighed og ende i det Dunkle;
hvad Vei man da vil gaae, bliver den Enkeltes Sag og
beroer paa en fri Beslutning; men er man, netop ved at
gaae ud fra Viden, naaet ud i det Dunkle, saa kan man
ikke, ved blot at gaae videre ud i Dunkelheden, komme
til en ny Art Erkjendelse, naar der da overhovedet skal
menes Noget ved Erkjendelse. Beviset for en saadan
Erkjendelses Mulighed hviler heelt igjennem paa de
mærkeligste Argumenter, hvor det „Absolut-Ueensartede“
gjør En haabløs i Retning af Forstaaelse. Blandt disse
% --- p. 37 ---
% (PDF 42)
\opage{37}\setcounter{footnote}{0}%
Argumenter findes der imidlertid et, som i mere end
een Henseende kan kaldes „enestaaende“: „Den tilsyneladende
Modsigelse, at Religionsphilosophien paa de opstillede
% UNSURE: „der“ — the d is damaged in the print (upper stroke missing, it looks like t); read der (alternative: a misprint „ter“)
Vilkaar maa blive en Videnskab, der handler
baade om det, der falder \emph{udenfor}, og det, der falder
\emph{indenfor} Videnskaben, er let at hæve. Vil Nogen
paastaae, at der umulig kan gives en Viden om en
Troeserkjendelse, naar Troeserkjendelsen efter sin Natur
skal falde udenfor Viden, lad ham da begynde med Begyndelsen
og paastaae Umuligheden af, \emph{at der kan
gives en Bevidsthed om Jorden, naar den virkelige
Jord efter sin Natur skal falde udenfor
Bevidstheden}.“\footnote{R. Nielsen. Om „Den gode Villie“ \&c. S. 137.}

Dette Argument hviler mærkelig nok paa den selvsamme
Feilslutning, som Professor Nielsen i sin Philosophiske Propædeutik
forekaster Theologerne under Navn af den „Anden
Sophisme.“\footnote{Philosophisk Propædeutik. S. 79.} At Jorden falder \emph{udenfor} Bevidstheden
er efter Art og Væsen noget ganske Andet,
end at Troeserkjendelsen falder \emph{udenfor} Viden. Atter
her drages der Slutninger fra et Forhold til et andet,
som er grundforskjelligt fra det første.

I \emph{første} Tilfælde nemlig (Jordens Forhold til
Bevidstheden) er „udenfor“ lokalt, og Problemet et rationelt
Erkjendelsesproblem, bekjendt nok fra Fichtes
„Wissenschaftslehre“; i \emph{sidste} Tilfælde (Troeserkjendelsens
Forhold til Viden) er „udenfor“ en reen Væsensbestemmelse,
og Problemet et suprarationelt Troesproblem.

% --- p. 38 ---
% (PDF 43)
\opage{38}\setcounter{footnote}{0}%
I \emph{første} Tilfælde refererer „udenfor“ sig til Forholdet
mellem Sandsning og Tænkning; i \emph{sidste} Tilfælde
til Forholdet mellem Troeserkjendelse og Viden.

I \emph{første} Tilfælde sigter „udenfor“ til noget \emph{ikke}
Absolut-Ueensartet; i \emph{sidste} Tilfælde til noget \emph{netop}
Absolut-Ueensartet.

I \emph{første} Tilfælde --- og dette er næsten det
Mærkeligste --- angiver „udenfor“ et Forhold, som efter
sin Natur heelt og holdent falder „\emph{indenfor}“ Videnskaben
(Bevidsthedens Opfattelse af den virkelige Jord);
i \emph{sidste} Tilfælde angiver „udenfor“ et Forhold, hvis
ene Led (Viden) heelt og holdent falder indenfor Videnskaben,
men hvis andet Led (Troen) heelt og holdent
falder udenfor den --- nemlig i Kraft af det „Absolut-Ueensartede.“

Der var endnu Anledning til at undersøge Professor
Nielsens dobbelte Videnskabsbegreb (Videnskab i
% NOTE: speck („jfr. Philosophisk“), 50 px vs full stop F=161 px, at x-height top in the word space; not transcribed.
videre og i snevrere Forstand; jfr. Philosophisk Propædeutik
S. 78), men det allerede Udviklede vil, for Videnskabslærens
Vedkommende, formeentlig være tilstrækkeligt.

Kun Eet staaer nu tilbage, at Professor Nielsen
giver os en i systematisk Sammenhæng udviklet Religionsphilosophi,
der ikke dogmatisk forsikkrer, \emph{at} en med
Videnstheismen \textbf{absolut} \emph{ueensartet} Troestheisme lader
sig gjennemføre, \emph{men gjennemfører den}.
% a short rule is printed below, closing section II; p. 38 ends here
% ===== batch b05_pp39-48 =====
% --- p. 39 --- (PDF 44)
\opage{39}\setcounter{footnote}{0}%
% Section head: centred bold Fraktur, numeral "III." in bold antiqua; one short rule below it.
\begin{center}\textbf{\textit{III.} Ideelæren.}\end{center}

Ideens første og almindeligste Kjendemærke er
Totalitet; derfor har ogsaa den philosophiske Ideelære,
i en ganske anden Betydning end nogen anden
Deel af Philosophien, Totaliteter at behandle; ja, man
% NOTE: speck („at den“), 60 px vs full stop F=163 px, in the word space just below mid x-height; not transcribed.
kunde i en vis Forstand sige, at den paa ethvert Punkt
har den hele Tilværelse til Gjenstand. Det lader sig
derfor tænke, at Behandlingen af Problemerne, saalænge
de endnu kun undersøges i en enkelt Disciplin af Philosophien,
kan indeholde endog væsentlige Feil, uden at
disse just saa lige opdages, fordi Synspunktet ifølge
Sagens Natur er valgt relativt. I Ideelæren derimod,
hvor Synspunktet er omfattende, kunne Feil og Selvmodsigelser
vanskelig skjules, og de faae her en, for det
hele System gjennemgribende, organisk Betydning.

Sin Ideelære har Professor Nielsen udviklet i
„Grundideernes Logik“, der, skjøndt den endnu langt fra
er afsluttet, dog ifølge Forfatterens egen Meddelelse
% footnote continues onto p. 40
danner et Hele for sig\footnote{Grundideernes Logik. I. Forord: „Endskjøndt dette Skrift
kun er en lille Deel af et større Hele --- neppe en Totrediedeel
af den første Niendedeel af den hele Logik ---, er det
dog ikke et Brudstykke, men et Hele for sig. Har den heri
angivne Plan og Methode nogen videnskabelig Betydning,
da maa Principet for Arbeidet i det Hele tydelig kunne erkjendes
af den foreliggende Deel. Men skulde det ved grundig
Undersøgelse lade sig godtgjøre, at Plan og Fremgangsmaade
i nærværende Skrift enten er uden al Betydning for
en videre Udvikling af Philosophien i Nutiden eller dog
lider af betydelige Grundmangler, da skal Forfatteren idetmindste
ikke forsøge at undskylde sig med, at Kritiken har
fældet Dom om den enkelte Deel uden at kjende det Hele.“}.

% --- p. 40 --- (PDF 45)
\opage{40}\setcounter{footnote}{0}%
„\emph{Grundideerne}“ bestemmes som de tre: \emph{Viden},
\emph{Magt} og \emph{Sandhed}. „Fordres der nu til en sand
Idee-Erkjendelse, at Modsætningen mellem den vidende
Subjectivitet og den objective Virkelighed maa udvikles,
for at Eenheden med sit rige Indhold ret kan erkjendes,
saa sees heraf, at en Ideens Logik maa forvandle sig til
en \emph{Grundideernes Logik}. Idet Virkelighedsideen
antager paa den ene Side Subjectivitetens, paa den anden
Side Objectivitetens Form, deler den sig i to oprindelige
Ideer: \emph{Videns Idee} for den subjective,
\emph{Magtens Idee} for den objective Virkelighed. Først
naar Extremerne af Begreb og Existens, af Ideebevægelse
og Realbevægelse ved denne Ideernes Modsætning
ere tilstrækkelig udviklede, faaer de modsatte Ideers concrete
Eenhed sin fulde Forklaring i \emph{Sandhedens Idee}.
Den i Typen angivne Fordobling vil typisk gjentage sig
i enhver Afdeling..... Endvidere kan den logiske
Type siges at reflectere Typen af det hele System saaledes,
at naar Videns Idee i sin Kreds giver os Logikens
Logik, saa giver Magtens Idee i sin Kreds Naturens
Logik, og Sandhedens Idee i sin Kreds Aandens
Logik. \emph{Saavidt om Anlæget}“\footnote{Grundideernes Logik. I. Indledning. S. XXXIII.}.

Vi anmode derefter Læseren om at henvende sig i
% --- p. 41 --- (PDF 46)
\opage{41}\setcounter{footnote}{0}%
Grundideernes Logik II. S. 369, altsaa langt inde i
Systemet, midt inde i Læren om Slutningerne; her hedder
det: „Gribe vi Ideen som Indhold i Formen, ville
de tre Grundbevægelser i den foregaaende Formudvikling
give os Ideen i Skikkelse af tre \emph{Grundideer}: Sandhed,
% PRINTED AS IS: the quotation opened at „Gribe vi Ideen“ is never closed: „Godhed*). Her“ prints no “, and nothing later on the page closes it.
Skjønhed og Godhed\footnote{Jfr. R. Nielsen. Grundideernes Logik. II. Kbhvn. 1866.
S. 383. Anm.}. Her overraskes man ved at
finde tre nye „Grundideer“, af hvilke kun een kjendes fra
Anlæget: Sandhedsideen. Naar nu Skriftet bærer Titelen:
% Only the ending „nes“ is letterspaced in „Grundideernes“ here (measured); „Grundideer“ is set solid.
Grundideer\emph{nes} Logik, saa maa det fremkalde
Forvirring, naar langt inde i Systemet nye Grundideer
paa eengang dukke op, uden at der gjøres Rede for,
hvorfra de komme; de tages som umiddelbart givne, indføres
i den formelle Analyse af det Almene, det Særskilte og
det Enkelte, og opstilles udenvidere, saa at Ingen ahner,
om de staae sideordnede med „Anlægets“ Grundideer ---
i saa Fald fik vi fem Grundideer --- eller om de „typisk
forekomme“ under enhver af de oprindeligt nævnte Grundideer,
eller om de maaskee henhøre under en enkelt
af dem.

Men afseet herfra, \emph{med hvilken Ret} kunne nu
Viden, Magt og Sandhed opstilles som Tilværelsens tre
Grundideer? Hvorfor just disse og ingen andre? Disse
Spørgsmaal opkaster enhver tænkende Læser, som tager
„Grundideernes Logik“ i Haanden og seer, at den stiller sig
den Opgave, at \emph{begynde} Philosophien. I Indledningen
% --- p. 42 --- (PDF 47)
finder man ingen Oplysning, men derimod den Med\opage{42}\setcounter{footnote}{0}delelse:
„Om det her angivne Schema iøvrigt vil føre til
sand Systematik eller gold Schematisme, maa Udførelsen
vise.“ Men i „Udførelsen“ faaer man heller
ingen Oplysning. I „Grundideernes Logik“ I. S. 434
opkaster Forfatteren vel det Spørgsmaal: „Hvilke og
hvormange Grundideer gives der?“ --- og kommer S.
450 til det Resultat, at Grundideerne ere de tre: Aandens
Idee, Naturens Idee og Sandhedens Idee; men
dette fremkommer som \emph{Konsekvents} af de foregaaende
Udviklinger og kan saaledes ikke gjælde som Beviis;
det vilde i ethvert Tilfælde blive en \textit{circulus in
probando}, da jo Indledningen opstiller de samme Ideer
uden at godtgjøre deres Gyldighed som Grundideer.
\emph{At} Indledningen gjør dette, vil man overbevise sig
om ved at vende tilbage til det ovenfor S. 40 Anførte.
Fordi nemlig „Virkelighedsideen paa den ene Side antager
Subjectivitetens, paa den anden Side Objectivitetens
Form“, følger da heraf, at Subjektivitetssidens
\emph{adækvate} „Grundidee“ er \emph{Viden}? Ingenlunde; den
har et ligesaa subjektivt, virkeligt, oprindeligt, kortsagt
adækvat Udtryk i \emph{Villien}. Og nu Objektiviteten, \emph{med
hvilken Ret} fremstilles den som \emph{Magtens} Idee? Endelig
indseer man ikke, hvorfor de modsatte Ideers
konkrete Eenhed skulde „faae sin \emph{fulde} Forklaring“ netop
i \emph{Sandhedens} Idee. I en Ideelære ere Betegnelserne
\emph{mere} end Betegnelser; de ere typiske og prægnante
for den følgende Udvikling; men i „Sandheden“ ligger
der Intet, som fører, endsige tvinger, til deri at søge
Videns og Magtens konkrete Eenhed. Som \emph{konkret}
% --- p. 43 --- (PDF 48)
\opage{43}\setcounter{footnote}{0}%
Eenhed skulde Sandheden være rigere end Viden og
rigere end Magten, men selve Kategorien „Sandhed“
lyder, hvor indholdsriig man saa gjør den, paa Noget,
som efter sin Charakteer er \emph{abstrakt}, i \emph{Sammenligning
med Magten}. Vel har Philosophen en relativ
Berettigelse til at vælge sine Betegnelser, selv om disse
ikke ganske svare til den almindelige Sprogbrug, men i
saa Fald er han endmere forpligtet til nøie at gjøre
Rede for Anvendelsen deraf.

De Gamles Trilogi, det Gode, det Sande og det
Skjønne, peger umiddelbart hen paa omfattende Sphærer
i Virkeligheden og har en saa stor Betydning i Philosophiens
historiske Udviklingsgang, at \emph{enhver} Grundideernes
Logik, selv om den ikke tager denne Trilogi som
Princip for Inddelingen, ja netop naar den ikke gjør dette,
strax i Begyndelsen maa paa det Nøieste gjøre Rede for sit
Forhold dertil. At der i Indledningen henvises til Hegels
Encyclopædi, hvis Inddelingsprincip kritiseres, og at der
i Forbigaaende tales om den videnskabelige Sandhed og
Troessandheden, gjør ingenlunde en slig Udvikling overflødig.

Men, afseet fra Sagens historiske Side, fordrer
„Grundideernes Logik“ selv en saadan Analyse; thi naar
man f. Ex. her (II. S. 376) erfarer, at „Sandhedsideen,
Ideen for den væsentlige Objectivitet, selv vilde
blive eensidig og usand, dersom den ikke havde Skjønhedsideen,
Ideen for den phænomenale Objectivitet, til
absolut Modsætning“, og naar man tilmed erindrer, at
begge Ideer ere \emph{Grundideer}, saa forudsætter en saadan
% --- p. 44 --- (PDF 49)
Udtalelse, om den ellers skal blive forstaaet og over\opage{44}\setcounter{footnote}{0}hovedet
have Betydning, at Forholdet mellem disse
Grundideer i Forveien har fundet sin Udvikling.

Idet der ud fra disse Forudsætninger bliver Spørgsmaal
om Ideelærens Forhold til Problemet Tro og
Viden, tage vi Udgangspunktet i den Hovedtanke, som
bestandigt vender tilbage i Professor Nielsens Philosophi,
at \emph{Ideen} som en \emph{Eenhed} behersker og forsoner de
Modsætninger, der ere traadte ud i Virkeligheden paa endelig
og brudstykkeagtig Maade. Men naar nu Mennesket
opfattes som det \emph{Afledede} i Modsætning til Ideen som
det \emph{Oprindelige}, hvorledes lader det sig da tænke,
% NOTE: speck („optræde som“), 50 px vs full stop F=158 px, at upper x-height right after „optræde“; not transcribed.
at Tro og Viden kunne optræde som \textbf{absolut} \emph{ueensartede}
Principer? For det Første bestemmes Virkeligheden
netop ved Endelighed, Udstykning og Relativitet,
saa at man allerede af denne Grund maatte opgive at
tillægge nogen „menneskelig“ Bestemmelse Prædikatet absolut;
men naar Mennesket tilmed sees under Synspunktet
af det \emph{Afledede}, bliver Selvmodsigelsen aldeles
evident, fordi der ingen Overgang lader sig paapege fra
Ideen til Mennesket eller omvendt. Vel tales der af og
til om, at der i Ideen selv finder \emph{absolut} Modsætning
% footnote continues onto p. 45
Sted mellem Grundideerne\footnote{Saaledes f. Ex. i Grundideernes Logik. II. S. 376: „At
den absolute Modsætning“ (mellem Sandhedsideen og Skjønhedsideen)
„maa grunde sig paa absolut Eenhed, vil i nærværende
Sammenhæng endvidere fremlyse deraf, at Skjønhedsideen
kun er sand ved at være Sandhedsideen modsat.“
For det Første maatte her nu konsekvent staae „absolut modsat;“
men af den Omstændighed, at disse to Ideer ere \emph{modsatte},
følger deraf, at de ere \emph{absolut modsatte}? Her
opnaaes ikke det Mindste ved at gjøre Modsætningen absolut;
det Selvsamme vilde være opnaaet ved at gjøre den
% PRINTED AS IS: the first dash is made of two short rules („-- hvad“), the second is a normal dash.
--- hvad den er --- relativ.}, og at den \emph{absolute}
% --- p. 45 --- (PDF 50)
\opage{45}\setcounter{footnote}{0}%
Modsætning ret maa bringes til Bevidsthed, for at den
\emph{absolute} Eenhed kan blive desto sandere, men hvad nytter
det bestandigt at bruge dette skjærende, ethvert Forhold
opløsende Udtryk „absolut“, naar der dog ikke bringes
Andet ud deraf, end hvad der ligesaa let kunde fremkomme
uden dette Udtryk? Aldrig har endnu nogen Philosoph
villet arbeide sig tildøde paa at vælte den Sisyphus-Steen:
i en \emph{absolut} Eenhed at ville forene \emph{absolut}
Modsatte, en Opgave, der --- naar man ret
gjør sig den levende --- kun har sit Sidestykke i Cirklens
Kvadratur. I og for sig kunde man gjerne betragte
denne mærkelige Anvendelse af det Absolute som
en Egenhed i Professor Nielsens Sprogbrug; men naar
det udtrykkelig bliver fremhævet som Kategori, som Noget,
hvorpaa der lægges det stærkeste Eftertryk og hvoraf der
drages Konsekventser i Philosophiens vigtigste Spørgsmaal,
saa maa Kritiken fra alle Sider søge at paavise
det Selvmodsigende deri.

Naar en saa mærkelig Lære som Troens og Videns
absolute Ueensartethed fremstilles, tilmed af en saa begavet
Mand som Professor Nielsen, saa indstiller det
Spørgsmaal sig ganske af sig selv, hvorledes vel hans
Begrundelse lyder. I Skriftet: Om „Den gode Villie“
% Fraktur abbreviation „rc.“ (= etc.) set as \&c.
\&c. findes S. 35, under Overskriften: „At Adskillelsen
% PRINTED AS IS: stray point („Viden ·ikke“), 410 px vs full stop F=173 px, raised to mid x-height, immediately before „ikke“.
af Tro og Viden ·ikke kan halvere den menneskelige Bevidsthed“,
følgende Forklaring: „Bevidstheden halveres
% --- p. 46 --- (PDF 51)
\opage{46}\setcounter{footnote}{0}%
ved relativt eensartede, men ikke ved absolut ueensartede
Factorer. Naar en Bevidsthed halveres i Tænken, faaer
den en halv Tanke istedenfor en heel; det, der mangler,
er Tanke; halveres den i Villie, faaer den en halv Villie
istedenfor en heel; det, der mangler, er Villie. Derimod
kan man ikke i strengere Forstand sige, at en Bevidsthed
halveres af Tænken og Villen saaledes, at Tankens
Halvdeel bliver Villie, Villiens Halvdeel Tanke; thi
Tanke kan kun i høist uegentlig Mening forvandles til
Villie, og Villie til Tanke\footnote{Men hvorledes kunne \textbf{absolut} \emph{ueensartede} Principer,
selv „i høist uegentlig Mening“, forvandles til hinanden?}. Staaer det da fast, at
Troen paa ingen Maade kan forvandles til Viden eller
Viden til Tro, saa staaer det ogsaa fast, at Troen ligesaa
umulig kan halveres ved Viden som Viden ved Tro.“

Men mod denne Forklaring gjøre vi \emph{for det Første}
gjældende, at Ingen paastaaer eller fornuftigviis kan paastaae,
at Tankens Halvdeel skulde være Villie eller omvendt;
at indføre en saa udvortes Betragtningsmaade i
Aandsvirksomheden vilde jo ligefrem være Galskab. Derimod
har man gjort gjældende og vil fremdeles gjøre
gjældende, at Tanke aldrig forekommer i Bevidstheden
uden Villie eller omvendt. Professor Nielsen gaaer ud
fra Tanke og Villie som \emph{absolut ueensartede} og
seer derfor i Villiens Nærværelse i Tankevirksomheden
en Halvering af Bevidstheden og af Personligheden. Vi
gaae tvertimod ud fra Personlighedens Eenhed som et
Givet og slutte derfra, at Tanke og Villie maae være
\emph{relativt eensartede}, og at den ene af disse Virksomheder
aldrig kan forekomme i Bevidstheden uden den anden.

% --- p. 47 --- (PDF 52)
\opage{47}\setcounter{footnote}{0}%
Anvende vi dette paa Tro og Viden, og antage vi
for et Øieblik, at de, som absolut ueensartede Principer,
kunne forefindes i den ene og samme Bevidsthed,
saa maa Bevidstheden \emph{heelt og udeelt} være gjennemtrængt
af Troen, være væsentligt bestemt som Tro, og
ligeledes \emph{heelt og udeelt} være gjennemtrængt af Viden.
% PRINTED AS IS: „giennemtrængt“ (i, not j; elsewhere „gjennemtrængt“).
Vælge vi nu en Bevidsthed, som er giennemtrængt
af det i sig absolute og ubetingede Vidensprincip, hvor
bliver saa Troen af? Er der endog blot Glimt af
Tro, hvorledes er da Overgangen mulig fra Viden til
% The line-end hyphen in „absolut-ueensartede“ is kept: elsewhere the two words stand apart, so it is not a word-break.
Tro, hvorledes er Overgangen mulig mellem to \textbf{absolut}-\emph{ueensartede}
Principer? Den er simpelthen \emph{umulig}.
Maaskee vil Professor Nielsen svare: Ganske vist, der
er ingen Overgang, og jeg vil heller ingen Overgang have.
Men hvorledes kan han saa tale om en Videnstheisme,
der \emph{afspeiler} Troestheismen og omvendt? Hvorledes
kan han fremdeles sige, at „gjennem Mysteriet \emph{gaaer
Veien} fra Tro til Viden, fra Viden til Tro“?

Om Professor Nielsens ovenfor anførte Begrundelse
gjælder det \emph{fremdeles}, at den beviser Formeget.
Den kan nemlig ikke indskrænkes til kun at gjælde Villie
og Tanke, forsaavidt der herved sigtes til Modsætningen
% NOTE: specks on the baseline at the line ends after „Ret“, „sig“ and „mange“ (p. 47), 24, 8 and 31 px vs full stop F=156 px; not transcribed.
mellem \emph{Tro og Viden}, men gjælder med samme Ret
Villie og Tanke indenfor det \emph{Rationelle}, ja lader sig
overhovedet udstrække til \emph{alle} de menneskelige Evner:
Følelse, Indbildningskraft, Abstraktion, Hukommelse, Dømmekraft
o. s. v., saa at vi kunne faae ligesaa mange
absolut ueensartede Principer i Bevidstheden, som der er
Evner. \textit{Ex. gr.}: „Bevidstheden halveres ved relativt
eensartede, men ikke ved absolut ueensartede Factorer.
% --- p. 48 --- (PDF 53)
\opage{48}\setcounter{footnote}{0}%
Naar en Bevidsthed halveres i Følelse, faaer den en
halv Følelse istedenfor en heel; det, der mangler, er
Følelse; halveres den i Dømmekraft, faaer den en halv
Dømmekraft istedenfor en heel; det, der mangler, er
Dømmekraft. Derimod kan man ikke i strengere Forstand
sige, at en Bevidsthed halveres af Følelse og
Dømmekraft saaledes, at Følelsens Halvdeel bliver Dømmekraft,
Dømmekraftens Halvdeel Følelse; thi Følelse
kan kun i høist uegentlig Mening forvandles til Dømmekraft,
og Dømmekraft til Følelse.“ \textit{Qui nimium probat
nihil probat.}

Der er af Biskop Martensen gjort den Indsigelse
mod „Grundideernes Logik“, at dens Gudsbegreb „er et
blot \emph{logisk-physisk}, der forholder sig i dogmatisk
Nægtelse til det \emph{ethiske} Gudsbegreb, der er den eneste
utvetydige Theisme“; og at Forfatteren, naar han udgiver
„Grundideernes Logik“ for „et Forsøg til en udpræget
Theisme“, „maa have taget Theisme i en heelt
anden Betydning end det i Almindelighed tages“.\footnote{H. Martensen. Om Tro og Viden. Kbhvn. 1867. S. 33.}
Hertil svarer Professor Nielsen: „Enten tør Gudsideen
slet ikke optages i en Ideernes Logik, eller, hvis Ideen
tør optages i en Logik, da maa Gudsbegrebet udvikles
logisk. En Udvikling af det logiske Gudsbegreb er saa
langt fra at beroe paa nogen Fornegtelse af det ethiske,
at den tvertimod indeholder alle for det ethiske Gudsbegreb
nødvendige Momenter. Vel er Logik ikke Ethik,
og de logiske Begreber ikke ethiske; men derfor maa
Ethiken ligefuldt have sin Logik, og de ethiske Begreber
% ===== batch b06_pp49-58 =====
% --- p. 49 ---
% (PDF 54)
\opage{49}\setcounter{footnote}{0}%
reflectere sig i den logiske Idee. Den Viden, jeg vil
indføre, er, det paastaaer Biskoppen, „i Virkeligheden
Pantheisme“; hvorfor? Fordi Logik skal være Ethik:
en ret interessant Confusion.“\footnote{R. Nielsen. Om „Den gode Villie“ \&c. S. 118.}

Dersom Biskop Martensen havde hentet sin Argumentation
fra Dogmatiken, vilde den ikke være kommen
paa Tale i disse Undersøgelser; vi maatte nemlig da i
Forveien have analyseret de eiendommelige Forudsætninger,
hvorpaa Dogmatik i det Hele hviler, for derefter
at bestemme Argumentationens Charakteer, hvilket ganske
falder udenfor vor Undersøgelses Plan; men Biskop
Martensens Argumentation er netop paa dette Punkt af
reent philosophisk Natur og holdes reent philosophisk,
saa at man altsaa her slet ikke har med en Dogmatiker
at gjøre.\footnote{Vi bemærke her, at vi ikke kunne slutte os til den af Biskop
Martensen i Skriftet „Om Tro og Viden“ givne Udvikling
af Theologiens Forhold til de rationelle Videnskaber.}

Efterat Biskop Martensen har henviist til den ontologiske
Subjektivitet, som er udviklet i „Grundideernes
Logik“, gjør han opmærksom paa, at denne ontologiske
Subjektivitet ikke kan danne Grundlaget for Theismens
Gud. Hans Argumentation lyder saaledes: „Men hvor
priisværdigt nu end dette er, saa opstaaer dog det
Spørgsmaal, om han (Professor Nielsen) ogsaa tilfulde
har benyttet sit psychologiske Udgangspunct, eller som vi
ogsaa kunne udtrykke det, om Subjectivitetsproblemet
% --- p. 50 ---
% (PDF 55)
\opage{50}\setcounter{footnote}{0}%
af ham er stillet i dets hele Omfang og Dybde. Thi at
spørge om den overempiriske (transscendentale) Grund til
vor psychologiske Bevidsthed er at spørge ikke blot om
Grunden til vor theoretiske, men ogsaa til vor \emph{practiske}
Erkjendelse, efterdi vi ikke kunne tænke, allermindst
videnskabeligt, uden at ville, og der ikke kunne være to
forskjellige, endsige to absolut ueensartede guddommelige
Principer, af hvilke det ene begrunder vor theoretiske,
det andet vor practiske Erkjendelse; det er at spørge om
Grunden ikke blot til vor Viden, men til vor \emph{Samvittighed};
ikke blot til vor theoretiske og speculative
Vished, eller \emph{vor Vished om at Verdenslovene
ere samstemmende med vor egen Fornufts Love
og derfor af os objectivt kunne erkjendes}, men
ogsaa til vor practiske Vished, eller \emph{vor Vished om,
at den objective Verden er modtagelig for vor
Villies Indvirkninger, og at vore fornuftige
Formaalsbegreber deri kunne realiseres}; eller
hvad der er det Samme, \emph{vor Vished om Samstemningen
mellem Sædelighedsloven i vort
Indre og den objective Verdensorden}, hvorved
da Begreberne Vished, Overbeviisning, Tro ville komme
under særlig Forhandling. Men tages Subjectivitetsproblemet
saaledes i sin Heelhed, tages den psychologiske
Bevidsthed i uopløselig Eenhed af Tænken og Villen,
Viden og Samvittighed, da opstaaer atter det store
Spørgsmaal, om den psychologiske Bevidsthed er mulig
under Forudsætning af et abstract og upersonligt Princip,
selv om dette kaldes et Subjectivitetsprincip, eller om
% --- p. 51 ---
% (PDF 56)
\opage{51}\setcounter{footnote}{0}%
der ikke her maa forudsættes et Princip, der ikke blot
er logisk, men ethisk, d. v. s. et virkeligt \emph{Personlighedsprincip}\footnote{H. Martensen. Om Tro og Viden. S. 37.}.“

Saaledes lyder den af al Dogmatik uafhængige
Indvending. Det er den rene Philosophies to Hovedproblemer,
som her indføres; det er det Grundlag,
hvorpaa Kants og Fichtes Philosophi hviler; Fichte begyndte
jo endog sin praktiske Videnskabslære med det
Spørgsmaal, hvorledes vi komme til at tilskrive os en
Virksomhed i en Sandseverden udenfor os. Det er
Spørgsmaal, som paatrænge sig ikke blot den videnskabelige
Philosoph, men ethvert tænkende Menneske; og
naar vi have arbeidet os gjennem en Tænkers System
og anvendt Tid og Kræfter paa en uendelig, ofte temmelig
formalistisk Detail, saa er det dog i sidste Instants
om disse Spørgsmaal, vi ønske Oplysning; det er
dem, \emph{der ligge os paa Hjerte}. Paa en Besvarelse
af disse Spørgsmaal har „Grundideernes Logik“ ikke
givet noget bestemt, afgjørende og utvetydigt Svar. I
Første Deel af „Grundideernes Logik“ træder vistnok
Objectiveringsproblemet stærkt frem, medens Forfatteren
under Udviklingen kritiserer de tidligere Løsninger af
dette Problem; men i Anden Deel, hvor man kunde
vente en mere sammentrængt og energisk Besvarelse fra
Forfatterens eget Standpunkt, forlades den tidligere
fulgte Plan, og Undersøgelserne antage saa kæmpemæssige
Dimensioner, at selv den meest opmærksomme Læser har
% --- p. 52 ---
% (PDF 57)
\opage{52}\setcounter{footnote}{0}%
Vanskelighed ved at følge med; Hovedproblemet bliver
her seet, saa at sige, gjennem en Spredelindse, der i den
Grad forstørrer (og forstyrrer), at Omridsene tabe
sig i Taager. Dette blot med Hensyn til Sagens
formelle Side; angaaende Realiteten henvises til det
Følgende.

Om det \emph{andet} Hovedspørgsmaal, som er den
\emph{praktiske} Philosophies Grundproblem (vor Vished om,
at den objektive Verden er modtagelig for vor Villies
Indvirkninger), finde vi slet ingen Oplysning i „Grundideernes
Logik“. Vel tør vi ikke sværge paa, at der
aldeles ikke lod sig paavise en eller anden Sætning, som
kunde tydes i den Retning; men, hvad det her kommer
an paa, Spørgsmaalet er intetsteds blevet Gjenstand for
en gjennemgaaende Undersøgelse og forekommer navnlig
ikke i Systemets Udgangspunkt; man sige nemlig ikke,
at det vil komme frem under „Sandhedens“ eller maaskee
% NOTE: speck („Idee. Det“), 67 px vs full stop F=168 px, at upper x-height in the word space; not transcribed.
allerede under „Magtens“ Idee. Det kan slet ikke
komme frem paa et senere Stadium --- uden ved et
Spring --- naar det ikke er med i Begyndelsen. Af
Intet kommer Intet, og af en Viden, der er absolut
ueensartet med Villie, kan Villien aldrig deduceres; men
er dette Tilfældet, hvorledes kan saa „Grundideernes
Logik“ ville afgive Grundlaget for en Villiesethik og
udgive sig for et Forsøg til en udpræget Theisme? Et
er det at forvexle Logik og Ethik, d. v. s. forlange, at
Villien skal udvikles \emph{paa samme Maade} i en Logik
som i en Ethik; et ganske Andet er det at fordre \emph{Villiens}
Logik udviklet i en \emph{Grundideernes}
% PRINTED AS IS: "Lvgik" -- the second letter of this "Logik" carries the v's top-left flag (cf. "udviklet", same line), a wrong sort for o (or a damaged o); UNSURE between the two.
Lvgik,
naar denne tilmed udtrykkelig sætter sig selv som Grundlag
% --- p. 53 ---
% (PDF 58)
\opage{53}\setcounter{footnote}{0}%
for en Villiesethik. At det netop er dette, som Biskop
Martensen har fordret, kan Enhver see ved at gjennemlæse
de ovenfor S. 49---51 citerede Linier.

Hvad det \emph{førstnævnte} Hovedproblem (det Theoretiske)
angaaer, begynder „Grundideernes Logik“ med en
Kritik af Hegels Udgangspunkt „den rene Væren“ og
sætter istedetfor denne „den rene Viden“. Hvad er nu
herved egentlig vundet? Hegels „rene Væren“ er Dogmatisme,
siges der; men er den „rene Viden“ Andet?
Vi sætte, at en Skeptiker melder sig --- og en saadan
melder sig altid --- og gjør følgende Indvending: Ved
at begynde Philosophiens System med „den rene Viden“
er man vel gaaet bagved Hegel, men hvad berettiger
overhovedet til at begynde paa denne Maade? Beviis
først, at Viden \emph{er}, ja endnu mere, beviis, at Viden
\emph{kan være}! --- hvad vil saa „Grundideernes Logik“
svare? Vi nære ingen Tvivl om, at Professor Nielsen
veed at svare; men Et er „Grundideernes Logik“, et
Andet er Professor Nielsen. At vi ikke opstille slige
Indvendinger som tilfældige Paastande, der udenvidere
kunne afvises, men at de ere begrundede i Sagens
Natur, i Videns eget Problem, vil, foruden af den
Omstændighed, at Spørgsmaalet gjør sig gjældende hos
Enhver, der begynder at bruge sin Forstand, ogsaa kunne
sees af den Omhyggelighed, hvormed Kant gjør Spørgsmaalet
om \emph{Videns Mulighed} til Gjenstand for Undersøgelse
i Indledningen til „Critik der reinen Vernunft“:
„Wie ist Metaphysik als Naturanlage \emph{möglich}?
Wie ist Metaphysik als Wissenschaft \emph{möglich}? Wie ist das
Vermögen zu Denken selbst \emph{möglich}? Dogmatism ist also
% --- p. 54 ---
% (PDF 59)
\opage{54}\setcounter{footnote}{0}%
das dogmatische Verfahren der reinen Vernunft, \emph{ohne
vorangehende Critik des eigenen Vermögens}.“

Det er i det Hele mærkeligt, at „Grundideernes
Logik“, som lægger an paa at \emph{begynde} den hele Philosophi
forfra, ikke saaledes optager Skepticismen som
Moment i sig, at den atter og atter opkaster Tvivl om
\emph{Muligheden} af sit hele Standpunkts Ugyldighed, og
gjendriver disse Tvivl. At igjen dette ikke er en tilfældig
Bemærkning, men en Hovedindvending, vil indsees
deraf, at Philosophien, til Forskjel fra de andre
Videnskaber, ikke blot reflekterer, men reflekterer over sin
egen Reflexion. Den kan derfor ikke nøies med, at
Tænkeren blot udvikler Skepsis under den Form, som
falder sammen med almindelig Selvkritik, og som
forbliver en privat Sag mellem Autor og ham selv,
men den fordrer udtrykkelig, at Tænkeren holder Dom
over sit eget System, at han lader Tvivlen om dets
\emph{mulige} Ugyldighed og Gjendrivelsen af disse Tvivl
% footnote continues onto p. 55, p. 56 and p. 57
blive til for vore Øine\footnote{Det er den her anførte Mangel paa Skepsis, som ofte forleder
% PRINTED AS IS: „afgjorte“ (for afgjørte): a clean o with an open counter, no trace of the slash that crosses the counter of ø in „anførte“ on the line above; a wrong sort, or an ø whose slash failed entirely.
Professor Nielsen til at slaae over i den meest afgjorte
Dogmatisme. Vi anføre et Par Exempler, der, om det
ønskes, yderligere kunne suppleres.
\par\noindent 1) I Indledningen til „Grundideernes Logik“ I. Side XXVII
hedder det: „Det er den videnskabelige Kritik ligesaa
umuligt at anerkjende Evangeliet om den Opstandne, som
det er Videnskaben umuligt at anerkjende Mythen om
Herkules. Hvad kan da \emph{hindre} Videnskaben i at erklære
Evangeliet om den Opstandne for en Mythe? \emph{Den
religiøse Ethiks Overlegenhed over den rationelle}!“
\par Hvoraf følger det udenvidere, at den religiøse Ethik er
% (footnote text on printed page 55 begins here)
den rationelle overlegen? Man kan jo ligesaavel antage det
Modsatte, uden derfor at opgive Videnskaben. Strauß,
Feuerbach o. s. v. øine her aldeles ingen Hindring:
% NOTE: speck („kan bevæge“), 59 px vs full stop F=158 px, just below the baseline after „kan“; not transcribed.
„Hvad kan \emph{bevæge} Videnskaben til at erklære Evangeliet
om den Opstandne for en Mythe? \emph{Den rationelle
Ethiks Overlegenhed over den religiøse}!“
\par\noindent 2) Philosophisk Propædeutik. S. 77: „At sætte Troen øverst
og dog beholde Videnskaben, er muligt; at sætte Videnskaben
øverst og dog bevare Troen, er umuligt.“
\par Hvorfor?
\par\noindent 3) Grundideernes Logik. I. S. 53: „Afskrækket ved Dogmatikens
mislykkede Forsøg med Testimoniet, kunde man
næsten føle sig fristet til med eet Slag at gjøre Ende paa
alle Qvaklerier med Subjectiviteten og dens Viden ved
at opstille og fastholde den Forudsætning, at den absolute
Aand er i sin Alvidenhed ligesaa ophøiet over al menneskelig
Tænkning som over al menneskelig Sandsning.“
\par Her kan den Troende svare: Fra mit Standpunkt seer
jeg heri intet saa Forfærdeligt; min Forudsætnings Gyldighed
gjendrives jo ikke derved, at den henstilles under
Synspunktet af det, „man næsten kunde føle sig fristet til
at gjøre“. --- Jeg gjør det udenvidere.
\par\noindent \emph{Fremdeles}: „Men, ikke at tale om det Uvidenskabelige i
at antage en guddommelig Viden, der skal falde saa
ganske udenfor vor menneskelige, at vi om en saadan,
strengt taget, ingen Viden kunne have“ . . .
\par Saa maa jeg i al Beskedenhed henregne mig til de
Uvidenskabelige, høre vi Kant sige; den Tanke, at vi om
den guddommelige Viden, strengt taget, ingen Viden kunne
have, har jeg hidtil holdt for \emph{videnskabeligt} udviklet i
„Critik der reinen Vernunft“; at den \emph{dogmatisk} henstilles
% (footnote text on printed page 56 begins here)
som \emph{uvidenskabelig} beviser jo Intet. Den maa gjendrives
\emph{med Grunde}.
\par\noindent \emph{Fremdeles}: . . . „saa lægger \emph{jo} den opstillede Forudsætning
et saa svælgende Dyb, ikke blot imellem den menneskelige
og den guddommelige Viden, men imellem den menneskelige
Viden og Sandheden, at det paa slige Vilkaar maa
blive Videnskaben umuligt at danne sig et sandt Begreb
om, hvad Sandhed er.“
\par Men, Herre min Gud, saa maa Videnskaben finde sig
i at opgive „at danne sig et sandt Begreb om, hvad Sandhed
er“. $\alpha$) Den Omstændighed, at Videnskaben maa opgive at
naae Sandheden, naar den guddommelige Viden \emph{ganske}
falder udenfor den menneskelige, er da ikke noget Beviis for, at
den guddommelige Viden \emph{ikke ganske} falder udenfor den
menneskelige. $\beta$) Man kan jo negte, at der overhovedet
gives nogen guddommelig Viden, medens man samtidigt
dermed hævder Muligheden af at naae den videnskabelige
Sandhed.
\par\noindent \emph{Fremdeles}: „Kan man nu paa dette Standpunkt ikke
vide, hvad Sandhed er, saa kan man jo heller ikke vide,
hvad der hører til en sand Viden.“
\par Ja, saa faaer man virkelig undvære det!
\par\noindent \emph{Fremdeles}: „Men en Subjectivitet, der saaledes faaer det
Sande udenfor sig og i den Grad farer vild fra al Viden
af Sandheden, at den tilsidst ikke veed, hvad sand Viden
er, farer vild i sig selv, og kan som en i sig vildfarende
Subjectivitet umulig unddrage sig den Selvmodsigelse,
hvoraf al anden Selvmodsigelse udspringer.“
\par Men hvoraf følger det udenvidere, at den Subjektivitet,
der antager vor fuldstændige Ikke-Viden om den guddommelige
Viden, derfor skulde „faae det Sande udenfor sig og
i den Grad fare vild fra al Viden af Sandheden, at den tilsidst
% (footnote text on printed page 57 begins here)
ikke veed, hvad sand Viden er“? Kant \emph{kunde} svare: Jeg lader
mig ikke skræmme af en „bussemandsagtig“ Dogmatisme
--- skjøndt Kant i sin Philosophi rigtignok aldrig charakteriserer
ved Hjælp af slige Bestemmelser; ønsker man at
vide, hvad han \emph{vilde} svare, da lyder det saaledes:
\par „Was nun die \textbf{Gewißheit} betrift, so habe ich mir
selbst das Urtheil gesprochen: daß es in dieser Art von Betrachtungen
auf keine Weise erlaubt sey, zu \textbf{meinen} und
daß alles, was darin einer Hypothese nur ähnlich sieht,
verbotene Waare sey, die auch nicht vor den geringsten
Preiß feil stehen darf, sondern, so bald sie entdeckt wird,
beschlagen werden muß.“ Critik der reinen Vernunft. Riga.
1781. Vorrede. S. 9.}.

% --- p. 55 ---
% (PDF 60)
\opage{55}\setcounter{footnote}{0}%
Vor næste Indvending gjælder den Maade, hvorpaa
„Grundideernes Logik“ behandler den \emph{theoretiske} Phi%
% --- p. 56 ---
% (PDF 61)
\opage{56}\setcounter{footnote}{0}%
losophi adskilt fra den \emph{praktiske}. Selv om den theoretiske
Viden var udviklet med ti Gange saa stor Ge%
% --- p. 57 ---
% (PDF 62)
\opage{57}\setcounter{footnote}{0}%
nialitet som Tilfældet er, saa vilde „Grundideernes
Logik“ dog aldrig kunne danne Udgangspunktet for en ny
Æra i Philosophien, fordi den behandler den theoretiske
Viden uafhængig af den praktiske; det var netop den
Feil, som Kant og Fichte begik, saa at de i denne Henseende
staae for os som et Exempel --- ikke til Efterlignelse
--- men til Advarsel. Hvormegen Genialitet
vi end ellers kunne møde i „Grundideernes Logik“,
hvor skarp og træffende en Kritik den ogsaa giver af
Kant, Fichte, Schelling og Hegel, hvor stærkt den end
accentuerer Virkeligheden --- den ligger dog i Rækken
af de philosophiske Forsøg, hvis Eensidighed betegnes
derved, at de udvikle det theoretiske Problem uafhængigt
af det praktiske. Hvorledes det System, der næste Gang
vil gjøre Epoche i Philosophiens Historie, kommer til
at see ud i det Enkelte, derom er det naturligviis umuligt
at danne sig nogen Forestilling; kun Saameget er
% --- p. 58 ---
% (PDF 63)
\opage{58}\setcounter{footnote}{0}%
vist, det vil ikke fra Begyndelsen af adskille den theoretiske
Philosophi fra den praktiske.

% UNSURE: "At" measured letterspaced (A-t gap 20 px at 300 dpi, vs 6-9 px in unspaced "At" pp. 53, 54); "den" between is unspaced (9-12 px); marked \emph{At} alone.
\emph{At} den \emph{theoretiske} Philosophies Grundproblem
(Spørgsmaalet om Verdenslovenes Overeensstemmelse
med vor egen Fornufts Love, eller med andre Ord,
Spørgsmaalet om Muligheden og Gyldigheden af en
objektiv Erkjendelse) i „Grundideernes Logik“ ogsaa virkelig
er udviklet uden Hensyn til Villien, sees allerede
af Kategoriskalaen: den umiddelbare, den reflekterende
og den begribende Subjectivitet, som angive theoretiske
Stadier indenfor den rene Viden. For at see, hvorledes
dette forholder sig, gaae vi tilbage til Begyndelsen
af „Grundideernes Logik“ og undersøge Systemet i
dets Tilblivelse, i dets første Spire, idet vi dog tilføie,
at den foreliggende Opgave ikke tillader os at udføre
denne Undersøgelse i dens hele Omfang. Vi ville \emph{her}
kun søge at paavise det formeentlig Uholdbare i selve
Udgangspunktet.

I Modsætning til Hegel, som søger Philosophiens
Begyndelsesbegreb i \emph{den rene Væren}, vil Professor
Nielsen begynde med \emph{den rene Viden}. Gjennem en
Kritik af „den rene Væren“ paaviser Forfatteren, at den
som „den yderste Abstraction af alt Værende“, som
den, der indeholder Alt og dog Intet, vel kunde kaldes
Begyndelsesbegreb, men at denne Begyndelse dog er
eensidig, fordi den hviler paa en eensidig Abstraktion.
„Som den yderste Abstraction af alt Værende maa jo
den rene Væren forudsætte Abstractionen, nemlig den
Abstraction, hvorved den selv fremkommer; men at forudsætte
Abstractionen er at forudsætte Tænkningen: den
% ===== batch b07_pp59-68 =====
% --- p. 59 ---
% (PDF 64)
\opage{59}\setcounter{footnote}{0}%
rene Væren forudsætter den rene Tænken. Saa lidt
som Begyndelsen kan skee med en Tænken uden Væren,
ligesaa lidt kan den skee med en Væren uden Tænken; men,
naar Begyndelsen selv forudsætter Eenhed af Tænken og
% NOTE: speck („rene Væren“), 37 px vs full stop F=173 px, just above the baseline right after „rene“; not transcribed.
Væren, saa begyndes her jo ikke med den rene Væren,
men med den rene Viden.“\footnote{Grundideernes Logik. I. S. 2---3.}

Ved Begyndelsen af et philosophisk System gjælder
det at være opmærksom paa hvert Ord og hver
Tankevending, fordi Begyndelsen skal indeholde Anlæget
til \emph{Alt} det, som i det Følgende vil fremkomme. Har
man først indrømmet Begyndelsen, saa kan man spare
sig alle Indvendinger i det Følgende. Hvad er her
altsaa skeet i denne Begyndelse?

For det Første tager Forfatteren sit Udgangspunkt
\emph{historisk}, forsaavidt han strax indfører den faktisk
foreliggende Begyndelse i det Hegelske System, „den
rene Væren“; dernæst tager han sit Udgangspunkt \emph{psychologisk},
hvad han ogsaa ligefrem er nødsaget til. Men
med hvilken Ret kan Forfatteren nu paa eengang begynde
at tale \emph{ontologisk}? S. 4 hedder det nemlig: „her
spørges ikke om min og din ufuldkomne Viden, men om
Videns evige Væsen.“ Hertil har han aabenbart ingen Ret.
Naar det nemlig hedder, at „den rene Væren forudsætter
den Abstraction, hvorved den selv fremkommer“, saa er
denne Væren kun \emph{psychologisk} tænkt, og saaledes Eet
med Værens \emph{Begreb}; thi det hedder jo udtrykkeligt,
at den „\emph{fremkommer}“ ved Abstraktionen; fremdeles er
Abstraktionen selv \emph{psychologisk}, og den rene Tænken
% --- p. 60 ---
% (PDF 65)
\opage{60}\setcounter{footnote}{0}%
\emph{psychologisk}. Den Sætning: „den rene Væren forudsætter
den rene Tænken“ har altsaa følgende Betydning
og Værdi: Naar vi i vor \emph{psychologiske} Bevidsthed
forefinde et saadant Begreb som „den rene
Væren“, saa forudsætter dette, at der i Forveien maa
være foretaget en \emph{psychologisk} Abstraktion, altsaa en
Tænken i den \emph{psychologiske} Bevidsthed. Altsaa:
den rene Væren er et \emph{Begreb}, og vi vide endnu aldeles
Intet, hverken om en \emph{ontologisk} Væren eller
om en \emph{ontologisk} Tænken.

Det hedder derefter, at „Begyndelsen ligesaalidt
kan skee med en Tænken uden Væren som med en
Væren uden Tænken“. Vi maae her indrømme Forfatteren
Følgende: „Begyndelsen kan ikke skee med en
Tænken uden Væren“, naar det forstaaes saaledes, at
den \emph{psychologiske} Tænken maa, for at existere, have
et Indhold, som efter sit Væsen er \emph{Begreb}; den rene
Væren er altsaa \emph{Begreb}. Ligeledes indrømme vi, at
„Væren ikke kan være uden Tænken“, naar det forstaaes
saaledes, at \emph{Begrebet} Væren, som Begreb i en \emph{psychologisk}
Bevidsthed, ikke kan existere uden denne psychologiske
Bevidsthed, uden at tænkes af denne Bevidsthed.
Fremdeles indrømme vi Forfatteren Ret til at indføre
den rene Viden som \emph{Benævnelse} paa Eenheden af
denne Tænken og denne Væren, idet her bestandig kun er
Tale om den \emph{psychologiske} Viden, opfattet i dens
yderste Abstraktion.

Hos Hegel, siger Forfatteren, finder der en dobbelt
Mystifikation Sted: For det Første er den rene Væren
hos Hegel „ikke en objectiv, væsentlig og virkelig Væren,
% --- p. 61 ---
% (PDF 66)
\opage{61}\setcounter{footnote}{0}%
men en abstrakt subjectiv, en reen Tankeværen“. Dernæst
„nedsætter han den tænkende Subjectivitet til en
ørkesløs Tilskuer ved, hvorledes den objective Væren
formaaer at tænke sig selv“ o. s. v. Fra sit eget
Standpunkt gjør nu Forfatteren (S. 4) gjældende, at
„Philosophien maa blive sig tydelig bevidst, hvad det
vil sige, at den udtrykkelig begynder, ikke med den rene
Væren, men med den rene Viden. \emph{Ved at sætte Begyndelsen
i Viden, er Tænkningens subjective
og Værens objective Væsen ligelig hævdet}.“

Her er den første Tvetydighed og den første Subreption.
Hvorledes komme vi udenvidere til at tale
om „Værens \emph{objective} Væsen“; i det Foregaaende
var der \emph{kun} Tale om Væren som \emph{Begreb}, thi det
% NOTE: speck at the line end after „gjennem“, 35 px vs full stop F=151 px, low in the x-height; not transcribed.
var jo (S. 3) den rene Væren, som \emph{fremkom} gjennem
Abstraktionen. Ved at tale om „Værens objective
Væsen“ gjøres her et Forsøg paa at komme udenfor
Bevidstheden; men hertil mangler der enhver Berettigelse;
vi maae altsaa blive, hvor vi ere, indenfor den \emph{psychologiske}
Bevidsthed.

Det hedder fremdeles (S. 4): „Idet Begyndelsen
udgaaer fra Viden, udgaaer den fra en tænkende og
vidende \emph{Subjectivitet}, og har i Subjectiviteten sin
ideelle \emph{Selvstændighed ligeover for den objective
Væren}, paa hvis Bestemmelse den væsentlig
% footnote continues onto p. 62
gaaer ud\footnote{Den anførte Sætning indeholder en Sprogfeil, som vi,
skjøndt det er en Bagatel, dog maae anføre, fordi vi have
Brug for Sætningen. Som den staaer i „Grundideernes
Logik“ er det \emph{Begyndelsen}, der „i Subjectiviteten har sin
ideelle Selvstændighed ligeover for den objektive Væren.“ Meningen
er naturligviis, at det er \emph{Viden}, der i Subjektiviteten
har o. s. v. Vel er Viden ogsaa Begyndelsen, men
den tilføiede Bestemmelse er her ikke charakteristisk for Viden
qua Begyndelse, men for Viden qua Viden.}.“

% --- p. 62 ---
% (PDF 67)
\opage{62}\setcounter{footnote}{0}%
For det Første indfører Forfatteren her udenvidere
Kategorien „Subjectivitet“ uden at give nogensomhelst
Forklaring af, hvad han derved forstaaer. Men dette er
meget uheldigt, naar Opgaven netop er at paavise den \emph{hele}
Philosophies \emph{Begyndelse}, hvor hver Begrebsbestemmelse,
ja selv hver Sprogvending er prægnant og afgjørende
for \emph{alt} det Følgende. Efter almindelig philosophisk
Sprogbrug kunde der paa dette Stadium end
ikke være Tale om Subjektivitet. Vi have den rene
Væren, den rene Tænken og den rene Viden --- \emph{ellers
Intet}. Af disse maatte alt det Følgende udledes;
skulde der altsaa tales om „Subjectivitet“, maatte den
for det Første være \emph{deduceret}, for det Andet være begrebsmæssigt
\emph{begrændset}. Som den her indføres, er
den hverken det Ene eller det Andet. Forfatteren indfører
udenvidere Subjektiviteten, om hvilken Enhver
veed, at den kan betyde det Allerforskjelligste. Men
hvilken Betydning Forfatteren nu ogsaa lægger i Kategorien,
saa har han, ifølge det Foregaaende, ingen Ret
til at lade den gaae ud over det \emph{Psychologiske}.

See vi derefter paa den anførte Sætning i dens Heelhed,
saa lykkes det heller ikke her at komme udenfor den
psychologiske Tænken; thi hvor objektiv man end gjør „den
objective Væren“, overfor hvilken Viden skal hævde sin
% --- p. 63 ---
% (PDF 68)
\opage{63}\setcounter{footnote}{0}%
ideelle Selvstændighed, saa er og bliver den en objektiv
Væren \emph{indenfor} Viden. I det Hele er Udtrykket „objektiv
Væren“ paa dette Sted tvetydigt; falder den \emph{indenfor}
Viden, saa maatte det udtrykkeligt være bemærket;
falder den \emph{udenfor} Viden, saa er her et
Spring i Udviklingen, netop paa det afgjørende Punkt.
At det er det sidste Alternativ, som er tilstede, sees af
de følgende Ord: „den \emph{reale} Væren er væsentlig
hævdet“; thi den \emph{reale} Væren kan ikke betyde Andet
end den objektive Virkelighed i Modsætning til den subjektive
Tænken i Bevidstheden. Forfatteren vil maaskee
svare, at han i de samme Linier begrunder sin Sætning,
naar han siger, „at medens det ingenlunde ved sig selv
er indlysende, at Væren overalt maa forudsætte Viden,
er det derimod ved sig selv indlysende, at Viden overalt
maa forudsætte Væren.“ Men i denne Begrundelses
\emph{første} Sætning er der \emph{deels} en Feil, \emph{deels} kommer
Forfatteren i Modsigelse med sig selv. Man kan nemlig
vel sige, at \emph{det Værende}, d. v. s. de enkelte tilværende
Ting, ikke overalt maa forudsætte Viden, nemlig
\emph{psychologisk} Viden, thi om Andet er her ikke Tale; derimod
lader det sig ikke sige om \emph{Væren}; thi \emph{Væren} existerer
kun i og ved en Abstraktion, og forudsætter, efter Forfatterens
egne Ord (S. 3), denne Abstraktion; forudsætter
altsaa Tænken, og forudsætter Viden, Eenheden af
% NOTE: „Væren.“ -- above the stop a slanted sliver, 47 px vs full stop F=186 px, 10 px above it (a colon point on this page is full-stop size at x-height top); not a colon, not transcribed.
Tænken og Væren. Hvad angaaer Begrundelsens \emph{anden}
Deel, at „Viden overalt maa forudsætte Væren“, da er
dette ganske rigtigt, men i denne Sammenhæng en Tautologi.
Nemlig: Hvorfra har Viden, i Forfatterens Fremstilling,
% --- p. 64 ---
% (PDF 69)
faaet den rene Væren? Ved en Abstraktion fra det Væ\opage{64}\setcounter{footnote}{0}rende
(S. 3). Altsaa: Hvorfor maa Viden forudsætte
Væren? Fordi den maa forudsætte det, den maa forudsætte.

Med Alt dette er 1) Forfatteren ikke kommen ud
over den psychologiske Viden. 2) Ved „ganske uskyldigt“
at begynde med Hegels „rene Væren“ og i den Anledning
„omtale“ det Værende, kommer dette som „af sig
selv“ ind med i Forhandlingerne; men vi savne den Deduktion,
som viser Overgangen mellem det Værende og
Viden; thi i Philosophiens \emph{Begyndelse} bliver et Udtryk
som „Abstractionen“ kun en Betegnelse. 3) Den „objective“
og den „reale Væren“ indføres ved en Subreption.

Forfatteren vil nu begrunde „den rene Videns“ Fortrinlighed
fremfor „den rene Væren“ gjennem to Argumentationer;
den rene Viden skal være for det Første
sandt \emph{Begyndelsesbegreb}, for det Andet \emph{Bevægelsesbegreb},
medens „den rene Væren“ ikke kunde gjøre
Fordring paa nogen af Delene. Vi ville nu forsøge at
bevise, at „den rene Viden“ hverken er et \emph{sandt} Begyndelsesbegreb
eller det \emph{hele} Begyndelsesbegreb.

1) \emph{„Den rene Viden“ er ikke et} \textbf{sandt} \emph{Begyndelsesbegreb}.
„At Viden som Begyndelsesbegreb forudsætter
Alt og indeholder Intet, er da ogsaa let at
paavise. I Viden maa nemlig Alt, hvad der paa nogen
Maade skal blive Gjenstand for philosophisk Behandling,
nødvendigviis opgaae. Gaves der nu en Art Værende,
der ganske faldt udenfor Viden, da maatte dette Værende
jo ogsaa falde udenfor Philosophien.“ Er dette et Beviis?
Lad os antage, at der gaves et Værende, som
faldt udenfor Philosophien, hvad beviser saa det? Ikke
% --- p. 65 ---
% (PDF 70)
det Allermindste. Lad os antage, at der gaves et Væ\opage{65}\setcounter{footnote}{0}rende,
som ganske faldt udenfor Viden, hvad beviser saa
% NOTE: speck („den, at“), flattened, 89 px vs full stop F=185 px (0.48F), on the baseline after the comma; not transcribed.
det? Ikke det Allermindste. Men Sagen er den, at
Forfatteren her laver sig til at gaae over fra den \emph{psychologiske}
til den \emph{ontologiske} Viden, uden at Overgangen
i mindste Maade motiveres eller retfærdiggjøres.
% NOTE: speck („hvad der“), 60 px vs full stop F=185 px, just above the baseline after „hvad“; not transcribed.
I Sætningen: „I Viden maa nemlig Alt, hvad der
paa nogen Maade skal blive Gjenstand for philosophisk
Behandling, nødvendigviis opgaae“ --- er der, deels
ifølge den foregaaende Udvikling (S. 2---3), deels ifølge
Sætningen selv, ikke Tale om Andet end om den \emph{psychologiske}
Viden; og Sætningen siger ikke Andet, end at
den endelige, psychologiske Viden maa, naar den philosopherer
over Noget, vide, hvad den philosopherer over,
og at den ikke kan philosophere over det, om hvilket
den Intet veed. I den næste Sætning skeer Springet
ganske uskyldigt og uformærkt, gjennem den forbigaaende
Parenthes: „thi her spørges ikke om \emph{min} og
\emph{din ufuldkomne Viden}, men om Videns evige Væsen.“
Jo her spørges netop om min og din ufuldkomne
Viden, og her har fra Begyndelsen af slet ikke været spurgt
om Andet. Hvorledes skulde ogsaa den Viden, som et Par
Linier iforveien charakteriseres, bestemmes og begrændses ved
det, „der paa nogen Maade skal blive Gjenstand for philosophisk
Behandling“, lige paa eengang blive \emph{ontologisk}?
I den \emph{ontologiske} Viden kunde jo Meget opgaae, som
slet ikke paa nogen Maade kunde blive Gjenstand for
philosophisk Behandling i den \emph{psychologiske} Viden.
Er denne Subreption først tilstede, saa gaaer det med
\emph{Begyndelsen}, med den \emph{objective, reale Væren},
med \emph{Andetheden} og med \emph{Magten} ganske som af sig selv.
% --- p. 66 ---
% (PDF 71)
\opage{66}\setcounter{footnote}{0}%
Men Systemet hviler i sin Begyndelse paa Feilslutninger
og Subreptioner. Udviklingen naaer ikke længere tilbage
end til den \emph{psychologiske} Viden og er forsaavidt
\emph{ingen sand} Begyndelse.

2) \emph{„Den rene Viden“ er heller ikke det}
\textbf{hele} \emph{Begyndelsesbegreb}. Havde Hegel ved at begynde
med den rene Væren deels ikke gjennemført Abstraktionen,
deels ført den i eensidig Retning, saa gjælder
det Samme om „Grundideernes Logik“, naar den
begynder med den rene Viden. For det Første nemlig
fører den kun Hegels Abstraktion et Stykke videre; for
det Andet overseer den, at \emph{Villien} er ligesaa oprindelig
som Viden. Skjøndt „Grundideernes Logik“ har udviklet
den rene Viden uden Hensyn til Villien, er Villien
mærkeligt nok alligevel med, eller, om man vil, naturligviis
med. Er det sandt, at Villien maa være med i
\emph{al} Tænkning som nødvendig og integrerende Faktor, saa
maa den ogsaa være med i Professor Nielsens Tænkning,
enten han nu ikke kan see det eller ikke vil indrømme
det. I Begyndelsen af „Grundideernes Logik“
bestemmes den rene Videns Fortrinlighed fremfor den
rene Væren hos Hegel derved, at den bestaaer sin
Prøve som \emph{Bevægelses}begreb. Det var en Illusion,
naar Hegel, ved at gaae ud alene fra det Almeenlogiske,
antog, at den rene Væren kunde bevæge sig selv og
udvikle sig selv til det almene Væsen, den almene Substants
o. s. v.\footnote{Jfr. Grundideernes Logik. I. S. 36.} „Væren for sig kan ikke være Bevægelsesbegreb;
den rene Væren kan, naar der abstraheres
% --- p. 67 ---
% (PDF 72)
\opage{67}\setcounter{footnote}{0}%
fra den abstraherende Tænkning, umulig bevæge sig selv.
Bevægelsen er, ifølge Hegel, en Overgang fra Væren til
Intet, og fra Intet til Væren; en saadan Overgang er
imidlertid kun en Værensbevægelse, forsaavidt det er en
Tankebevægelse; men en ved Tanken bestemt Værensbevægelse
er jo udtrykkelig en Vidensbevægelse“ .... „Det
Sande heri er, at Viden i det Uendelige forudsætter sig
selv, bevæger sig fra sig selv gjennem alt Andet til sig
selv, saa at dens hele Udvikling væsentlig er en Selvudvikling.
At Viden --- ikke min eller din ufuldkomne
Viden, men Viden i sin Oprindelighed --- forudsætter
sig selv og begynder med sig selv, vil da med andre
Ord sige, at den har sit Princip i sig selv, at dens
% PRINTED AS IS: stray point in the footnote („I.· S. 5.“), 435 px vs full stop F=191 px, raised to upper x-height, right after „I.“.
Væsen er Selvhed“\footnote{Jfr. Grundideernes Logik. I.· S. 5.}.

Men hvorledes foregaaer en Videns\emph{bevægelse}?
Aabenbart ved det i Viden indeholdte Villiesmoment.
At Viden \emph{forudsætter} sig selv, at den \emph{bevæger} sig
fra sig selv gjennem alt Andet tilbage til sig selv, hviler
i sin sidste Grund paa det i Viden indeholdte logiske
Villieselement. Der er i Forudsætningen Noget, som
\emph{tvinger} Tænkningen hen til Konsekventsen, men der er
i Forudsætningens \emph{Indhold} i og for sig (d. v. s. naar
det ikke er \emph{bestemt som} Forudsætning) Intet, der
tvinger Tænkningen hen til Konsekventsens \emph{Indhold};
tager man Forudsætningens Indhold blot som \emph{Indhold}
(uden at \emph{bestemme} det som Forudsætning), saa
bliver Tanken staaende derved som ved et Givet, og
rører sig ikke af Stedet. Naar derimod Indholdet
% --- p. 68 ---
% (PDF 73)
% PRINTED AS IS: „be-|temmes“ (= bestemmes): p. 67 ends „be=“, p. 68 begins „temmes“; the s is missing and line 1 of p. 68 starts c. one letter-width in from the margin
\emph{be\opage{68}\setcounter{footnote}{0}temmes} til at være Forudsætning, saa følger Konsekventsen
med Nødvendighed. Idet altsaa Nødvendigheden
forudsætter en \emph{Bestemmen}, forudsætter den en
\emph{Villiesakt}; og da al Tænken efter sit inderste Væsen
foregaaer gjennem en uafbrudt Række af Forudsætninger
og Konsekventser, saa finder her bestandig et Vexelforhold
Sted mellem Tænkning og \emph{Villie}. At dette ikke foregaaer
paa endelig, udstykket Maade, naar der tales om
det ontologiske Forhold, forstaaer sig af sig selv; en
gjennemført Dialektik af Videns Mulighed og Videns
Virkelighed vil vise, at Viden og Villie paa ethvert
Punkt forudsætte hinanden. Uden Villie staaer Viden
stille og opløser sig selv, fordi den, ligesom den rene
Væren, mangler enhver \emph{Grund} til Bevægelse.

At vi i Ideelæren finde Viden \textit{clauda altero pede}
kan iøvrigt ikke forundre, naar vi erindre den skarpe
Grændse, der i Videnskabslæren blev draget mellem personlige
% PRINTED AS IS: stray point in the footnote („S. 21. ·“), 395 px vs full stop F=175 px, at x-height top, free-standing c. 186 px after „21.“.
og upersonlige Sandheder\footnote{See ovenfor S. 21. ·}. Det er den samme
Eensidighed, som her kommer igjen under en anden
Form. Afsondret fra Villien vil den rene Viden aldrig
kunne give noget Grundlag for en Villiesethik, fordi
den ikke kan give en Villiens Logik. Det Citat af
Schleiermachers Dialektik, som Biskop Martensen anfører
i sit Skrift: Om Tro og Viden, samler som i en Hovedsum
alle disse Indvendinger og gjælder ikke blot
Kant og Fichte, men ogsaa og i fuldeste Maal Forfatteren
af „Grundideernes Logik“: „Wir bedürfen eben
so wohl eines transscendentalen Grundes für unsere Ge%
% ===== batch b08_pp69-78 =====
% --- p. 69 ---
% (PDF 74)
\opage{69}\setcounter{footnote}{0}%
wißheit im \emph{Wollen} als für die im \emph{Wissen, und
beide können nicht verschieden seyn}.“\footnote{H. Martensen. Om Tro og Viden. S. 38. Anm.}

Skjøndt disse Undersøgelser ikke have til Formaal
at analysere „Grundideernes Logik“, men kun forsaavidt
berøre den, som den staaer i Forbindelse med Problemet
Tro og Viden, skulle vi dog, for at undgaae
enhver Misforstaaelse, fremstille vore Indvendinger skarpt
og tydeligt, \emph{forsaavidt de forekomme i disse Undersøgelser}:
% Layout: the objections are set as a list, not run on: 1)-4) each begin a new line, the number standing in
% from the margin with the text hanging; a)-c) under 1) each begin a new line, indented further and hanging.
% No paragraph indents. Under 3) (p. 70) the a) and b) run on in the text.
\par\noindent 1) „Grundideernes Logik“ angiver \emph{thetisk}, at Viden,
Magt og Sandhed ere Grundideerne, men\\
\textit{a)} den beviser ikke, \emph{at de ere det};\\
\textit{b)} den beviser ikke, at Grundideerne ere \emph{disse tre
og ingen andre};\\
\textit{c)} den indfører paa et senere Stadium: det Gode,
det Sande og det Skjønne som Grundideer,
uden at gjøre Rede for, \emph{hverken} hvorfra de
% PRINTED AS IS: stray point („betyde, .eller“), 107 px vs full stop F=164 px, sunk just below the baseline, immediately before the second „eller“.
komme, \emph{eller} hvad de betyde, .\emph{eller} hvorledes
de forholde sig til de oprindelig angivne Grundideer.\\
2) „Grundideernes Logik“ sætter sig som Formaal,
\emph{principielt} at paavise \emph{Philosophiens Begyndelse}; den begynder med den rene Viden,
altsaa med Videns Virkelighed, men paaviser ikke
Videns \emph{Mulighed}.\\
3) „Grundideernes Logik“ sætter sig som Formaal,
\emph{principielt} at paavise \emph{Philosophiens Begyndelse}; men Begyndelsen, den rene Viden,
% --- p. 70 ---
% (PDF 75)
\opage{70}\setcounter{footnote}{0}er \textit{a)} \emph{ingen sand} Begyndelse, fordi det, der
udgives for ontologisk Viden, kun er psychologisk
Viden; og \textit{b)} \emph{ingen heel} Begyndelse, fordi den
rene Viden uden Villie er ligesaa afmægtig som
den rene Væren.\\
4) „Grundideernes Logik“ seer i den \emph{ontologiske
Subjektivitet} Grundlaget for et \emph{Personlighedsprincip}, der nærmere bestemmes som \emph{Villiesprincip}; hvorledes kunne nu Villiens og
Personlighedens Begreber \emph{udledes af} den rene
Viden?

% "Absolut-" ends the line, "Ueensartede" begins the next with a capital:
% a hyphenated compound, hyphen kept.
Der er i Professor Nielsens Philosophi to Hovedbestræbelser,
som begge maa strande paa det Absolut-Ueensartede.

% "A." is a small antiqua capital.
\textit{A.} Den ene Bestræbelse gaaer ud paa at lade \emph{Ideelæren afgive Grundlaget for en tilsvarende
Religionsphilosophi}. At denne Bestræbelse er tilstede,
derfor have vi Forfatterens Ord: „Ihvorvel en
Grundideernes Logik ikke i noget Parti kan falde umiddelbart
sammen med Religionsphilosophien, saa er det
dog aldeles nødvendigt, at den, medmindre den da skulde
kunne hjælpe sig med et atheistisk Princip, maa udvikle
saavel Gudsideens som Verdensideens logiske Form, maa
logisk bestemme Forholdet mellem disse Ideer \emph{og saaledes afgive Grundlaget for en tilsvarende
Religionsphilosophie}.“\footnote{R. Nielsen. Om „Den gode Villie“ \&c. S. 116.}

En saadan Bestræbelse kan altid have sin Berettigelse
som Forsøg; tilmed maa Religionsphilosophien
% --- p. 71 ---
% (PDF 76)
\opage{71}\setcounter{footnote}{0}have sine Forudsætninger, hvis den ikke skal falde ned
fra Skyerne. Et ganske andet Spørgsmaal er det derimod,
hvorvidt en rationel Ideelære \emph{kan} afgive Forudsætninger
for en paa Troen bygget Religionsphilosophi,
og dernæst, hvorvidt \textit{in casu} „\emph{Grundideernes Logik}“
kan gjøre det. At den \emph{ikke kan} det vil allerede sees af
det Udkast, der gjøres i Skriftet: Om „Den gode Villie“
\&c. S. 138: „Der skal nu paavises en alsidig articuleret
Vexelbestemmelse imellem de ueensartede\footnote{Det er ret mærkeligt, at der \emph{netop paa dette Sted} staaer
\emph{ueensartede} istedetfor \emph{absolut ueensartede}.} Begreber;
det skeer saaledes, at Vidensbegreberne tjene som intellectuelt
Underlag, Troesbegreberne derimod træde i Forgrunden
og udvikles systematisk. Af Vexelforholdet mellem
Begreberne fra de to Sphærer fremlyser saa Vexelforholdet
imellem den videnskabelige og den religiøse
Theisme, den Theisme, hvis Gudsbegreb kan udvikles
med logisk Nødvendighed i en Grundideernes Logik, og
den Theisme, hvis aabenbarede Gud er Gjenstand
for Tro.“

Det vilde nu have været saare heldigt, hvis man
havde faaet lidt Oplysning om, hvorledes denne Anviisning
lader sig udføre; thi fordi det hedder, „der
skal nu paavises; det skeer saaledes; heraf fremlyser
saa“, følger jo ikke, at der paavises, skeer eller fremlyser
det Allermindste. Der opkastes strax efter (S. 138)
det Spørgsmaal, hvorledes det er muligt, at der
kan tilveiebringes Overeensstemmelse imellem de absolut
% --- p. 72 ---
% (PDF 77)
\opage{72}\setcounter{footnote}{0}ueensartede Begreber, men man venter forgjæves Besvarelsen,
thi de Ord (S. 140): „nei, Modsigelserne
(i Troesbegreberne) forklares, i Conseqvens af Principet,
udtrykkeligt deraf, at de manglende Mellembestemmelser
ligge skjulte i det absolute Mysterium“, ---
indeholde ikke Glimt af Forklaring.

At nu det ovenfor givne Udkast til en Religionsphilosophi,
med Vidensbegreberne som intellektuelt Underlag,
ikke lader sig realisere, er indlysende af de Fordringer,
der stilles; thi der skal paavises, ikke blot en
artikuleret, men en \textbf{alsidig} \emph{artikuleret} Vexelbestemmelse
mellem de \textbf{absolut} \emph{ueensartede} Begreber; og der skal
fremdeles paavises et, med \emph{logisk Nødvendighed}
udviklet, videnskabeligt Gudsbegreb. Umuligheden heraf
er indlysende; thi Eet af To: \emph{Enten} maa „Grundideernes
Logik“ tage Ordet Gud i en anden Betydning
end den, hvori Religionen tager det (i Forbindelse med
en Række positive Troesdogmer), og da er det kun en
Subreption at tale om et videnskabeligt Gudsbegreb,
der i dette Tilfælde ikke er Andet end den gamle, velbekjendte
% UNSURE: square blot fused to the foot of the final „a“ of „ogsaa“ (excess ink c. 267 px vs full stop F=173 px, baseline to 22 px above; not a separate mark, the text layer reads a full stop); read as a blot (alternative „ogsaa.“).
Omskrivning af Ideen; \emph{eller} ogsaa maa „Grundideernes
Logik“ tage Ordet i den strengt religiøse Betydning, (og
dette er aabenbart Meningen), men hvilken Opgave den
end stiller sig, og hvad Navn den end giver sine
Afspeilinger af Gudsbegrebet, om „logiske Afspeilinger“
eller „rationelle Ækvivalenter“ --- Et slipper den ikke
for: de positiv-religiøse Bestemmelser i Gudsbegrebet.
Naar nu „Grundideernes Logik“ tager sit Udgangspunkt
i den rene Viden, \emph{hvor findes der da i dette}
% --- p. 73 ---
% (PDF 78)
\opage{73}\setcounter{footnote}{0}\emph{Begreb en indre Drift, som paa nogen Maade
med logisk Nødvendighed kan føre til} Bestemmelser
som Inspiration, Inkarnation o. s. v., med
samme Nødvendighed altsaa, hvormed vi føres til Begreber
som f. Ex. Forudsætning og Konsekvents, Grund
og Følge, Væsen og Skin? Umuligheden er indlysende
ved sig selv. \emph{Hvorledes skulde den rene Viden
paa nogetsomhelst Punkt af sin Udvikling
overhovedet kunne falde paa slige Begreber
eller endog blot paa en Afspeiling af dem?}
Svaret er overmaade simpelt: Det er den gamle Historie,
som gjentager sig blot under en ny Form, at
Videnstheismen laaner slige Begreber fra Troestheismen
og behandler dem saalænge, indtil den omsider paastaaer
selv at have udviklet dem med logisk Nødvendighed.
Paaviis blot een eneste Troesbestemmelse, hvortil den
rationelle Viden med logisk Nødvendighed kan udvikle det
rationelle Ækvivalent, og vi indrømme hele Resten!

Man kan stille sig paa det Standpunkt, at man
reentud negter Gyldigheden af det Positiv-Religiøse;
det er et Standpunkt, man kan forstaae. Man kan
% "non-" (antiqua, with hyphen) ends the line, "liquet" begins the next:
% hyphen kept as part of the phrase.
fremdeles lade Spørgsmaalet staae hen under et \textit{non-liquet};
% NOTE: speck („er idetmindste“), 43 px vs full stop F=170 px, low in the x-height in the word space; not transcribed.
det er idetmindste ærligt. Men Et er der,
man ikke kan, og det er at gaae ud fra den rene Viden
og derefter med logisk Nødvendighed udvikle om
det saa kun er Afspeilinger af et theistisk Gudsbegreb
i en „Grundideernes Logik“, naar det vel at mærke skal
være Alvor med Theismen. Men maaskee vil der blive
svaret, at den rationelle Theisme ikke selv naaer eller
kan naae til Bestemmelser som Inspiration og Inkarnation,
% --- p. 74 ---
% (PDF 79)
\opage{74}\setcounter{footnote}{0}men at den som i en Afspeiling modtager dem fra
Troestheismen til rationel Bearbeidelse. I saa Fald
maa den da opgive den logiske Nødvendighed.

% "B." is a small antiqua capital.
\textit{B.} Herved komme vi til den anden Hovedbestræbelse,
at lade \emph{Troestheismen afgive Grundlaget for
Videnstheismen}. At denne Bestræbelse er tilstede,
derfor have vi Forfatterens Ord: „Det ligger i Sagens
Natur, at Charakteermærkerne i Videnstheismen maae
speile Charakteermærkerne i Troestheismen som Nødvendighedsbestemmelser
Frihedsbestemmelser, som Vidensbegreber
Troesbegreber. Under denne Forudsætning vil
man forstaae: hvorledes en christelig Verdensbetragtning
vil, ikke umiddelbart, men middelbart, kunne faae \emph{Indflydelse}
paa Philosophien som \emph{autonomisk} Videnskab;
hvorledes \emph{Henblik} til Troestheismen kan være
\emph{veiledende} ved en \emph{streng videnskabelig} Udvikling
af Videnstheismen; hvorledes den i Grundideernes Logik
„udprægede“ Theisme kan --- \emph{uden i noget Punkt}
at overskride Grændsen for det \emph{Logisk-Physiske} ---
være \emph{bestemt ved} og \emph{bestemmende for} en tilsvarende
Religionsphilosophie.“\footnote{R. Nielsen. Om „Den gode Villie“ \&c. S. 136.}

Hvorledes nu dette skal forstaaes, kan sees af de
Kategorier, som i den Anledning sættes i Bevægelse:
„Til Troesbegrebet Skabelse svarer f. Ex. et Vidensbegreb
Skabelse, hiint udviklet af Nødvendigheden, dette
af Friheden; til Troesbegrebet Gudmenneske svarer Vidensbegrebet
Gudmenneske o. s. v.“\footnote{Smstds. S. 137.} Man kan her
% --- p. 75 ---
% (PDF 80)
\opage{75}\setcounter{footnote}{0}fortsætte: Til Troesbegreberne Forløsning, Forsoning, Opstandelse
og Himmelfart svare Vidensbegreberne Forløsning,
Forsoning, Opstandelse og Himmelfart; thi med disse Bestemmelser
staaer og falder Gudmenneskets Begreb. Hvorledes
seer nu en videnskabelig Himmelfart ud, ---
navnlig naar den er udviklet med logisk Nødvendighed?
Det vilde være et frugtbart Thema for Anbringelsen af
de Godtkjøbs-Vittigheder, som saa ofte gavmildt kastes ud
til et Publikum, der hverken forstaaer sig paa Sagen
eller bekymrer sig om Grunde, men alene i Forventningen
om slige Herligheder som forslugne Haifisk følger
i Philosophiens Kjølvand.

Naar Udgangspunktet er taget i Troestheismen, forholder
det sig naturligviis anderledes med Begreber som
% NOTE: speck („v., fordi“), 85 px vs full stop F=183 px, at x-height top just before „fordi“; not transcribed.
Gudmenneske, Forløsning, Forsoning o. s. v., fordi de
her udenvidere ere forudsatte og givne. Vanskeligheden
kommer nu fra den modsatte Side, idet nemlig Troestheismen
skal være \emph{veiledende} ved en \emph{streng videnskabelig}
Undersøgelse af Videnstheismen. Hvorledes
er dette muligt? Ved en Afspeiling svares der; men
hvormegen \emph{kategorisk} Valuta er der i en Afspeiling,
og hvorledes skal man stille sig, for at formaae Troesbestemmelserne
til at lade sig afspeile? Thi hvorledes
man end stiller sig f. Ex. overfor Troesbestemmelsen
Gudmenneske, saa vender den altid en Side til, som
ikke vil lade sig afspeile. Det gaaer her som med
Manden, der forundrede sig over, at man aldrig kunde
faae Maanen at see paa Vrangen, fordi den altid
vendte den samme Side mod Jorden, men som dog var
af den Mening, at naar man med et Speil blot kom
„langt nok ud“, saa gik det.

% --- p. 76 ---
% (PDF 81)
\opage{76}\setcounter{footnote}{0}Og dog har Forfatteren et saadant Speil; det er
\emph{Meningernes} Speil. „Meningen“, siger Forfatteren,
„er et særligt psychologisk Medium, hvori de „to absolut
ueensartede Principer“ kunne mødes og skilles og paa
en Maade træde i Underhandling med hinanden . . . .
Mening er det rette Navn, den kategoriske Betegnelse for
det søgte Medium. Den vældige Modsætning mellem
Oldtidens og Nutidens Anskuelse af Himmel og Jord,
. . . . vil, ikke forsvinde, men lade sig udligne ved
fortsatte Forhandlinger paa Meningens neutrale Grund.“\footnote{R. Nielsen. Om „Den gode Villie“ \&c. S. 62.}

\emph{For det Første}: Hvad vil \emph{neutral Grund}
sige? Det vil sige en Grund, hvor der ikke gives
„Grunde“; thi, siger Forfatteren, Viden hører til Videnskaben
og er ikke neutralt Gebet, og Troeserkjendelsen
udelukker Videnskaben, og frembyder altsaa heller
ikke noget Neutralt. En saadan Grund er en mærkelig
Grund, allersomhelst naar den er Grund for Forhandlinger.

\emph{Dernæst}: Er den nu virkelig ogsaa \emph{neutral}?
Forfatteren siger, at vi kunne bevare „Fædrenes“ Tro
og \emph{forkaste} deres Mening. Men saa er Grunden jo
\emph{ikke} neutral; thi efter almindelig Menneskeforstand ere
\emph{Forkastelse} og \emph{Neutralitet} hinanden udelukkende
Begreber.

\emph{Fremdeles}: I Meningen kunne Tro og Viden
mødes og skilles og \emph{paa en Maade} træde i Underhandling
med hinanden. Men hvorledes kan den strengt
videnskabelige Viden indlade sig med Meninger og gjennem
% --- p. 77 ---
% (PDF 82)
\opage{77}\setcounter{footnote}{0}dem som Mæglere (Media) korrespondere med Troen?
Hvorledes kan Viden, uden at opgive sin \emph{Charakteer
af} Viden (som den, der \emph{konsekvent} hævder sine Forudsætninger),
gaae ind paa \emph{neutral} Grund, hvor der
kun maa menes?

\emph{Ydermere}: Hvorledes foregaaer en \emph{Forhandling}
mellem \emph{Meninger}, ja mellem Tro, Viden og
Meninger? Hvorledes foregaaer en \emph{Forhandling} paa
\emph{neutral} Grund, naar den neutrale Grund ikke er et
Territorium, et Stykke Jord, der i Egenskab af Jord
ikke kommer Forhandlingen ved, men naar Neutraliteten
selv deeltager i Forhandlingerne og dog er en \emph{Væsensbestemmelse},
der \emph{udelukker} Differentserne? Vel
bruges der Udtrykket „Forhandlinger \emph{paa} Meningens
neutrale Grund“, et Udtryk, som leder Forestillingen
hen paa Meningen som et ideelt Territorium, i Grunden
et videnskabeligt Genf, hvor der holdes Fredskongres
mellem Tro og Viden. Men denne Forestillings Realitet
udelukkes igjen ved selve Fremstillingen, ifølge hvilken
Meningerne jo selv deeltage i Kongressen. Men hvorom
Alting er, naar nu Forfatteren udtrykkeligt siger, at „den
vældige Modsætning“ \emph{vel ikke vil forsvinde, men
lade sig udligne ved fortsatte Forhandlinger
paa Meningens neutrale Grund}, saa see vi ikke
rettere, end at Biskop Martensen træffende og aldeles adækvat
har betegnet Resultatet af de Forhandlinger, Forfatteren
indleder mellem Tro og Viden, som et \textbf{Konkordat}.
Hvorledes det iøvrigt vil gaae paa denne Fredskongres,
er saare let at indsee: Saasnart Videnskritiken arriverer,
% --- p. 78 ---
% (PDF 83)
\opage{78}\setcounter{footnote}{0}vil den, som en anden videnskabelig Garibaldi, udenvidere
smide Meningerne ud, og dermed er Freden „sluttet“.

\emph{Endydermere}: Dersom „Forskjellen mellem Tro
og Mening \emph{kun} kan klares i en Troeskritik, ligesom
Forskjellen mellem Viden og Menen\footnote{Her synes at spores en svagt antydet Tendents til at skjelne
mellem \emph{Mening} og \emph{Menen}, saa at maaskee \emph{Mening}
skulde anerkjendes, \emph{Menen} forkastes; det kan dog vel aldrig
være \emph{Meningen}?} \emph{kun} kan klares i
en Videnskritik“, hvorledes kunne da Tro og Viden mødes
paa Meningernes \emph{neutrale} Grund, da jo Videnskritiken
og Troeskritiken ere \emph{absolut} ueensartede?

\emph{Tilsidst}: Hvad berettiger Forfatteren til at indføre
Mening i \emph{en} Betydning i Viden, i en \emph{anden}
Betydning i Troen? Vel hedder det, at „Fædrenes“
Mening \emph{forkastes}, men der tales tillige om en \emph{Udligning}
paa Meningens neutrale Grund, altsaa er
Forkastelsen dog kun relativ. Og naar det fremdeles
efter hvert afgjørende Spørgsmaal hedder: „Ja, derom
er der forskjellige Meninger“ --- hvorledes gaaer det
saa igjen an at forkaste Fædrenes Mening, da disse
jo kunne sige: „Vor Mening er ligesaa god som Jer'
Mening“? Jo, svares der, det gaaer an, thi paa
Culturbevidsthedens nuværende Trin kunne vi \emph{umulig}
opfatte Fortællingen om Paradiset i samme Mening som
Fædrene. Men dette „umulig“ er \emph{her} dogmatisk, eftersom
Forholdet mellem Tro og Mening \emph{kun} kan klares
i en Troeskritik, \emph{men} Troeskritiken blæser ad Kulturbevidstheden,
og \emph{dog} skal Troen mødes \emph{netop med}
% ===== batch b09_pp79-88 =====
% --- p. 79 ---
% (PDF 84)
\opage{79}\setcounter{footnote}{0}%
\emph{Viden} paa Meningens neutrale Grund, altsaa \emph{dog}
forhandle med Kulturbevidstheden, hvilket den \emph{alligevel
ikke} kan, da Troeskritik og Videnskritik ere absolut
ueensartede.

\emph{Og saa til allersidst}: Hvorledes kunne Meningerne
indføres som \emph{kategorisk Betegnelse} i Videnskaben
% UNSURE: „alene“ — the fourth letter shows two serifed stems with no visible arch (reads „aleue“); the n-hairline has dropped out elsewhere on this page (Meningerne, below), so taken as n (alternative: turned n, „aleue“ as printed).
uden netop og alene for at blive \emph{forkastede}? Thi i
\emph{Videnskaben} indføres de, deels fordi deres Anvendelse
paa Fortolkningen af Kjendsgjerningerne i den hellige
Historie dog vel er en videnskabelig Anvendelse;
deels og fornemmelig, fordi i denne Sammenhæng \emph{den
philosopherende Bevidsthed og Meningerne
falde sammen}. Dersom vi nu udenvidere henstillede en
saadan Paastand, vilde det simpelthen være en Uhøflighed
mod Forfatteren, en af de plumpe Grovheder, der, naar
de anvendes i Videnskaben, tilsidst som en Vexel vende
tilbage til Udstederen. Men ligesom vi i disse Undersøgelser
egentlig kun forholde os betragtende og lade
Forfatterens egne Udsagn kæmpe mod hinanden, saa at
de paa en Maade blive Vidner i Sagen: Nielsen \textit{contra}
Nielsen, saaledes ville vi ogsaa i nærværende Tilfælde
lade Forfatterens egne Ord tale for sig selv. Vi citere
% „ꝛc.“ (r rotunda + c, et cetera) is transcribed „\&c.“ throughout this batch.
derfor to Steder i Skriftet: Om „Den gode Villie“ \&c. Først
et Udsagn, som vi allerede eengang have citeret, men som
vi anføre igjen ifølge den gamle Sætning: \textit{opposita
juxta se posita magis illucescunt}. Altsaa:
% Layout: the two citations (S. 62, S. 136) are set as an indented block, pp. 79--80, each opening with a paragraph indent inside the block.
\begin{quote}
S. 62: Spørge vi til Slutning, om her \textbf{i en
Bevidsthed, der skal kjende og erkjende baade
Troens og Videns Princip}, da ikke er et særligt
psychologisk \textbf{Medium}, hvori de „to absolut ueens%
% --- p. 80 ---
% (PDF 85)
\opage{80}\setcounter{footnote}{0}%
artede Principer“ kunne mødes og skilles og paa en
Maade træde i Underhandling med hinanden, svare
vi: Jo! her er ganske rigtigt et saadant Medium.
Men hvad skulle vi vel kalde dette Medium? \textbf{Viden}
kunne vi da \textbf{ikke} kalde det; thi Viden hører til Videnskaben
og er ikke neutralt Gebet. Tro kunne vi
heller ikke kalde det, thi Troeserkjendelsen udelukker
Videnskaben, og frembyder altsaa heller ikke noget
Neutralt. Vi maae da kalde det \emph{Mening}. \textbf{Mening
er det rette Navn, den kategoriske Betegnelse
for det søgte Medium.}

Hermed sammenholdes S. 136: \textbf{Den philosopherende
Bevidsthed er, hvad der da forstaaer
sig af sig selv, det subjective Medium, hvori
de ueensartede Begreber mødes.}
\end{quote}

I den første Udtalelse er \emph{Meningen} det søgte
Medium; i den anden Udtalelse er \emph{den philosopherende
Bevidsthed} det søgte Medium. Mon saa ikke,
idetmindste i dette Tilfælde, \emph{den philosopherende Bevidsthed
falder sammen med Meningen}? Men
ikke engang i denne Selvmodsigelse er der Konsekvents; thi
inden i denne Selvmodsigelse opdage vi en ny Selvmodsigelse,
saa at den første Selvmodsigelse ikke engang faaer
Lov at blive i Charakteren af Selvmodsigelse, men modsiger
sig selv \textit{qua} Selvmodsigelse. Thi medens det i
den anden Udtalelse hedder, at den \textbf{philosopherende}
% PRINTED AS IS: „Bevidsthed- er“, a stray hyphen after „Bevidsthed“ (mid-line, not a line-end break).
Bevidsthed- er Mediet, saa hedder det i den første Udtalelse,
at \textbf{Viden} kunne vi da \textbf{ikke} kalde Mediet.

Endelig erindre vi, at det i Skriftet: Om „Den
gode Villie“ \&c. S. 33 hedder: „Som negativ Midte
% --- p. 81 ---
% (PDF 86)
\opage{81}\setcounter{footnote}{0}%
\textbf{forener Mysteriet} derimod \textbf{de adskilte Extremer} og
viser sig som Regionernes Gjennemgangsmedium: \textbf{gjennem
Mysteriet gaaer Veien fra Tro til Viden, fra
% PRINTED AS IS: stray apostrophe-shaped sort („Troende ’og“), 13x30 px, 196 px vs quote half c. 300 px on the page, above the x-height, immediately before „og“.
Viden til Tro.} Idet den Troende ’og den Vidende
begge paa een Gang rette deres Blik mod \textbf{Mysteriet,
udlignes Striden}, og \textbf{Extremerne enes.}“
% Layout: „Altsaa:“ begins a new line flush left (no paragraph indent); the three statements are set as a list with hanging indents, each on its own line.
\par\noindent Altsaa: Først er \textbf{Meningen} „det rette Navn, den
kategoriske Betegnelse for“ Mediet.\\
Derefter er den \textbf{philosopherende} Bevidsthed,
„hvad der da forstaaer sig af sig selv“,
Mediet.\\
Endelig er \textbf{Mysteriet} Regionernes Gjennemgangsmedium.

Til Udligningen af Differentserne mellem Tro
og Viden behøves der altsaa den eiendommelige Forudsætning,
at \textbf{Meningen} = den \textbf{philosopherende} Bevidsthed
= \textbf{Mysteriet}.

Idet vi imidlertid støtte os til Forfatterens egne
Ord, at „Videnskabens Opgave netop er at vide det
Ueensartede som ueensartet“\footnote{R. Nielsen. Om „Den gode Villie“ \&c. S. 137.}, og idet vi denne Gang,
følgende Forfatterens Exempel, lade Troeskritiken afspeile
Videnskritiken og Videnskritiken igjen afspeile
Troeskritiken, dog saaledes, at vi indføre Videnskritiken
som Grundlag for Troeskritiken, og idet vi tilføie,
at vi ikke blot gjøre det denne Gang, men at vi altid
% PRINTED AS IS: „saa“ (before „gaae“) is set in bold Fraktur, probably a wrong-fount sort.
og under alle Omstændigheder ville gjøre det: \textbf{saa} gaae
vi videre i Texten.

% --- p. 82 ---
% (PDF 87)
\opage{82}\setcounter{footnote}{0}%
Det hedder i et „oplysende“ Exempel, at „hvad
der i første Mosebog fortælles om Paradiset, om Livets
Træ, Kundskabens Træ o. s. v., ganske vist er fortalt
i den Mening, at det Altsammen hører under det
udvortes Virkelige, at Livets Træ og Kundskabens Træ
ligesaavist have havt deres Rødder i Havens sandselige
materielle Jordbund, som enten Viintræet eller Figentræet.“
Ja, ganske vist var det deres \emph{Mening}, men
det var ogsaa deres \emph{Tro}. Men saa bliver her jo ingen
\emph{Forskjel} mellem Tro og Mening, og det var dog
en saadan, Troeskritiken lovede. Jo der gjør, thi
Troen kunde jo holde sig til Fortællingens \emph{indre} Virkelighed.
Men den kunde da ogsaa holde sig til den
\emph{ydre}; thi herom, siger Forfatteren, er der forskjellige
Meninger. „Men saa er det jo kun en Mythe? Ja,
derom kan der igjen være forskjellige Meninger“. Med
denne Fart paa Udligningernes Skraaplan vilde vi
% NOTE: speck („dersom ikke“), 15 px vs full stop F=198 px, near the baseline right after „dersom“; not transcribed.
være om en Hals, dersom ikke Forfatteren standsede
Farten ved at nedramme to afgjørende Sætninger.
Den \emph{første} Sætning lyder saaledes: „Dersom der nu
ingen væsentlig Forskjel var mellem Tro og Mening,
saa vilde en Forsoning af Religion og Cultur ganske
vist være umulig.“\footnote{R. Nielsen. Om „Den gode Villie“ \&c. S. 63.} Man spørger uvilkaarligt: Hvorfor?
--- men faaer intet Svar. Ligesom Forfatteren
har forelsket sig i Biskop Martensens \emph{Dogmatik},
saaledes have vi igjen forelsket os i Forfatterens \emph{Dogmatisme}.
Som vi mødte den i den \emph{første} Sæt%
% --- p. 83 ---
% (PDF 88)
\opage{83}\setcounter{footnote}{0}%
ning, saaledes møde vi den igjen i den \emph{anden} Sætning:
„Men er her en Forskjel, en principiel Forskjel,
saa er Forsoningen mulig“. Beviset for Dogmatismens
Tilstedeværelse ligger deri, at Forfatteren \emph{hverken}
godtgjør, at der er nogen principiel Forskjel, \emph{eller}
at Forsoningen under denne Forudsætning er mulig.
Nemlig:

1) „At her \emph{virkelig} er en principiel Forskjel mellem
Tro og Mening, skal kjendes derpaa, at vi kunne bevare Fædrenes
Tro og forkaste deres Mening.“ Det vilde nu være
meget godt, naar man blot kunde stole paa, at Fædrene rykkede
ud med Meningen; men fordi det hedder, at det ganske
vist er Meningen i første Mosebog, og ganske vist
er Meningen hos vore luthersk-orthodoxe Fædre, saa
følger da ingenlunde deraf, at det ganske vist \emph{er}
Meningen; og da vi tilmed ere paa Meningernes neutrale
Grund, saa kan man jo med samme Ret sige,
at det ganske vist slet ikke var Meningen. Men da
vi saaledes paa ingen optænkelig Maade kunne faae
Vished om, hvad Meningen nu virkelig var, saa kunne
vi heller ikke \emph{forkaste} den; men kunne vi ikke forkaste
den, saa kunne vi heller ikke paavise nogen \emph{principiel}
Forskjel mellem Tro og Mening. \textit{Quod erat demonstrandum.}

2) Men selv om vi kunde paavise denne principielle
Forskjel, hvorfor følger saa deraf, at Forsoningen mellem
Religion og Kultur er mulig? Derom oplyses der
ikke det Allermindste. At henstille en saadan Sætning
uden Beviis, er Dogmatisme. \textit{Quod erat demonstrandum.}

% --- p. 84 ---
% (PDF 89)
\opage{84}\setcounter{footnote}{0}%
Idet nu Forfatteren, som sagt, udtrykkeligt\footnote{R. Nielsen. Om „Den gode Villie“ \&c. S. 63.} henfører
de nyligt omtalte Exempler paa Forskjellen mellem
Tro og Mening under Troeskritiken, saa maae vi
--- skjøndt i en anden Mening --- tiltræde Forfatterens
Anskuelse, at „Troeskritiken som Kritik ikke er \emph{opbyggelig}.“\footnote{Smstds. S. 62.}
Ved imidlertid at svinge mellem
Skriftet: Om „Den gode Villie“ \&c. og „Grundideernes
Logik“ II. S. 194---201 kommer Tanken som et philosophisk
Kompensationspendul til at regulere sig selv.
Med andre Ord: \emph{Hvad Forfatteren indfører og
hævder i det førstnævnte Skrift, gjendriver
og forkaster han i det sidste}. Den mulige
Indvending, at vi paa det ene Sted have en Troeskritik,
paa det andet en Videnskritik, betyder Intet, da der
aldeles ikke er beviist, at Meningerne \emph{kunne} have nogen
anden Betydning i den ene Slags Kritik end i den anden.
Omkvædet: „Ja, derom er der forskjellige Meninger“,
% NOTE: stray grave-like stroke above the x-height between „Samme“ and „og“, 90 px vs full stop F=190 px; not transcribed.
lyder paa det Samme og betyder det Samme, enten vi
tale om Tro eller Viden. Hvorledes skulde ogsaa Meningerne,
det Løseste af alt det Løse, pludselig faae Betydning,
fordi der tales om Troesbestemmelser? Altsaa: I
„Grundideernes Logik“ II. S. 195 siges der, at det
er Culturstadiet, de subjective Reflexioners Stadium, der
egentlig er at opfatte som Meningernes Stadium, hvor
vi faae en Hærskare af politiske, æsthetiske og \emph{religiøse}
Meninger. Det første Stadium er det individuelle
(S. 196): „Hvert Individ, selv en Skoledreng, har jo
% --- p. 85 ---
% (PDF 90)
\opage{85}\setcounter{footnote}{0}%
paa dette Trin en uforgribelig Mening“. Paa det
næste Stadium stiger Meningen til \emph{Paastand}: „Idet
den Menende bliver paastaaelig, seer det unegtelig ud,
som om han vilde \emph{paatrænge Andre sin tilfældige
Subjectivitet}; men dette er nu ikke Meningen.
Meningen er, at den Paastaaelige ivrer for Sagen, \emph{er den
Eneste, der seer Sagen}, som den virkelig er. Den Paastaaelige
mener kun, at hans Mening ubetinget bør være \emph{Alles}
Mening, og det af den simple Grund, at det er den
% NOTE: speck („Mening.“ Idet“), 35 px vs full stop F=163 px, at mid x-height in the word space; not transcribed.
eneste rigtige Mening.“ Idet nu de mange ulige Meninger
„geraade“ i Strid med hinanden, er det „Opgaven
at finde den for Alle gjældende Mening, den almindelige
Mening, nærmere bestemt den offentlige Mening, \emph{Opinionen}“.
Herefter charakteriserer Forfatteren Kampene
% „Kunst og“: printed with a very narrow space (c. 2 mm, a third of the other word spaces on the line).
for Opinionen i Politik, Kunst og \emph{Religion} og opstiller
enkelte Kriterier for de politiske, æsthetiske og \emph{religiøse}
Meningskampe. Kun den sidste interesserer os her; om de
theologiske Meningsudtalelser yttrer Forfatteren sig (S. 201)
saa grumme haardt: „Slige Paastande ere \emph{forkastelige},
ja \emph{foragtelige}, naar der strides for en Sandhed“.
Men naar nu Forfatteren \emph{selv} indfører Meningerne i
Forhandlingen mellem Tro og Viden, og naar han,
hvor han taler om sig selv, i et ganske anderledes mildt
og forsonligt Sprog tilføier,\footnote{R. Nielsen. Om „Den gode Villie“ \&c. S. 64.} at han, „ved deres Anvendelse
paa Fortolkningen af Kjendsgjerningerne i den
hellige Historie, \emph{har øinet en Udsigt til}, at de
gamle Stridigheder om den hellige Skrifts Troværdighed,
maatte, vistnok ikke ende, men dog indtræde i en ny
% --- p. 86 ---
% (PDF 91)
\opage{86}\setcounter{footnote}{0}%
Phase“ --- \emph{hvad er saa dette}? Der var nok dem,
som vilde bruge ligesaa skarpe Udtryk om Forfatteren,
som han selv nyligt brugte; vi indskrænke os til følgende
Forespørgsel: Naar det i „Grundideernes Logik“ II. S.
197 hedder: „Det Charakteristiske er, at hvormeget her
end tales om Principer og Principspørgsmaal, saa kommer
det dog \emph{aldrig, og kan, i Meningernes,
Paastandenes, Forsikkringernes, de assertoriske
Dommes skarpt afgrændsede Sphære,
aldrig komme til nogen egentlig Principernes
Undersøgelse}“ --- hvorledes forholder det sig saa i
Skriftet: Om „Den gode Villie“ \&c. S. 62---64? Er
\emph{der} da ikke Tale om Princip og Principspørgsmaal?
Eller skulle maaskee Meninger, Paastande og Forsikkringer
nok gjælde i \emph{Troeskritiken}, men ikke i \emph{Videnskritiken}?

Ved saaledes at see \emph{Meningen} indført i \emph{Videnskaben},
indført, saa at sige, som Mægler mellem Himmel
og Jord, bliver man tilmode som den, der befinder sig
paa et Observatorium, netop som der gjøres en stor
Opdagelse, vi ville antage Opdagelsen af en \textit{nova stella},
thi \emph{Meningen} er en \textit{nova stella} i Forfatterens \emph{System};
den har de ægte Kriterier: den kommer pludselig
tilsyne, forbauser Samtiden og forsvinder igjen, efterladende
de Lærde i mangehaande Tvivl om Phænomenets
Betydning og i Forventning om, naar der igjen skulde
vise sig Noget. Ingen kan derfor fortænke os i, at vi
have taget de nøiagtigst mulige Observationer og derved
sikkret os en paalidelig Ephemeride for paakommende
Tilfælde.

% --- p. 87 ---
% (PDF 92)
\opage{87}\setcounter{footnote}{0}%
Imidlertid ere vi dog under al denne Menings-Forskjellighed
enige med Forfatteren i to Ting. For
det Første indrømme vi ham, „at det, naar virkelige
Analyser foretages, maa paa Meningernes Omraade
komme til \emph{mangfoldige Conflicter}“. For det Andet
see vi heller ikke rettere, end „at de gamle Stridigheder
om den hellige Skrifts Troværdighed, ved Meningernes
Anvendelse paa Fortolkningen af Kjendsgjerningerne i
den hellige Historie, maatte, vistnok ikke ende, men dog
\textbf{indtræde i en ny Phase}“, --- idet vi dog nærmere
bestemme denne Phase som \emph{sidste Kvarteer}.

Men i eet Punkt ere vi komplet uenige med Forfatteren.
Thi om Nogen, i Anledning af Forfatterens
Udvikling af Meningerne, spurgte os: Men er dette
Videnskab? Er dette Kritik? Er dette Begrundelse?
% UNSURE: „Philosophi“ — the second o is open on the right and reads as c („Philoscphi“); taken as a damaged o.
Er dette Philosophi? --- da maatte vi, tvungne ved
Grundenes Magt, svare: Nei Kjære, derom er der kun
een Mening.

Skulde Læseren muligen, i Anledning af den vidtløftige
Exkursion om Meningen have glemt, hvor vi
vare, saa ville vi komme hans Hukommelse tilhjælp:
Talen var om at sætte Troestheismen som Grundlag
for Videnstheismen, i hvilken Anledning vi fremhævede,
at dette ikke lod sig gjøre, fordi Troesbestemmelserne slet
ikke vilde lade sig afspeile. Vi skulle i den Anledning
endnu kun tilføie, at Troestheismen paa sin Viis er
ligesaa nøieregnende som den rationelle Videnskab og
ikke taaler, at dens positive Dogmer omsættes i Videns
fortyndede Medium og saa endda kaldes theistiske.

Hvorledes og hvorvidt endelig de i Troesbegreberne
% --- p. 88 ---
% (PDF 93)
\opage{88}\setcounter{footnote}{0}%
selv indeholdte Modsigelser kunne overvindes, det bliver
den Troendes Sag; den rationelle Viden kan i saa
Henseende spare sig alle Forklaringer; den kan aldrig,
med noget for Videnskaben væsentligt Udbytte, afspeile
Troesbegreberne, og hvad Troestheismen angaaer, da
betakker den sig for Videnstheismens „gode Villie“. Men
vil man endelig forsøge sig paa en Forklaring af de i Troesbegreberne
selv indeholdte Modsigelser, saa maa man
levere Noget, der har virkelig Betydning, Noget, som
er klart og tydeligt, og ikke istedetfor en \emph{saadan} Efterviisning
--- som her er det Eneste, noget Menneske
bryder sig om --- give en Forklaring som denne:
„At Troesbegreberne, naar de udvikles af deres eget
Princip, indeholde Modsigelser, kan \emph{ikke}, som fra et
scholastisk-orthodox Standpunkt, forklares deraf, at den
syndige Fornufts oprindelige og nødvendige Kategorier
skulde staae i Strid med den aabenbarede Sandhed,
\emph{heller ikke}, som fra et kritisk-videnskabeligt Standpunkt
% NOTE: speck („deraf, at“), 36 px vs full stop F=187 px, on the baseline between „deraf“ and the comma; not transcribed.
deraf, at Troesbegreberne selv skulde være usande; \emph{nei},
(nu kommer Forklaringen) \emph{Modsigelserne forklares,
i Conseqvens af Principet}, \textbf{udtrykkeligt} \emph{deraf,
at de manglende Mellembestemmelser ligge
skjulte i det absolute Mysterium}.“ Men dette
er jo ingen Forklaring, det er simpelthen en dogmatisk
Paastand; thi Forfatteren kjender ligesaalidt som vi
Mellembestemmelserne, kan altsaa ikke i Konsekvents af
Principet(?) drage nogensomhelst Slutning angaaende
Mysteriet. I Mysteriet svimler det for Bevidstheden,
og ingen menneskelig Forstand øiner Muligheden af ad
denne Vei at naae til rationelle Ækvivalenter i Viden
% ===== batch b10_pp89-93 =====
% --- p. 89 --- (PDF 94)
\opage{89}\setcounter{footnote}{0}%
for det, man har maattet opgive i Troen. Saalænge
dette Spørgsmaal ikke er besvaret i sammenhængende
Totalitet, og saa længe Afspeilingen og de rationelle
Ækvivalenter aldrig udtydes anderledes end ved en
Henviisning til Mysteriet, saalænge maae vi nøies med
at see dem ikke blot som i et Speil, men nærmest som
i en mørk Tale.

% p. 89: a centred rule (c. 1.8 in.) closes section III; the head follows,
% with a shorter centred rule below it.
\begin{center}\rule{4em}{0.4pt}\end{center}

\begin{center}\textbf{\textit{IV.} Den philosophiske Methode.}\end{center}

\begin{center}\rule{2em}{0.4pt}\end{center}

Methoderne i den nyere Tids Philosophi charakteriseres ved de to Hovedretninger: den \emph{kritiske} og den
\emph{dialektiske}.

Den \emph{kritiske} Methode, som især blev udviklet af
Kant, lægger an paa saa skarpt som muligt at adskille
Begreberne og at fastholde dem i Adskillelsen, for paa
denne Maade at opnaae den størst mulige Klarhed og
Tydelighed.

Den \emph{dialektiske} Methode, taget i moderne Forstand, blev opdaget og indført af Hegel; den optager
den kritiske Methode i sig som Moment. Vel skal Tænkningen begynde med en skarp Adskillelse af Begreberne,
men den maa ikke blive staaende herved som sidste Resultat. I de sondrede Begreber ligge nemlig \emph{Modsigelser} skjulte, d. v. s. Momenter, som paa eengang
% --- p. 90 --- (PDF 95)
udelukke og bekræfte hinanden. Først naar disse Mod\opage{90}\setcounter{footnote}{0}sigelser vækkes, begynder det egentlige Tankeliv; først
derved, at Vexelforholdet mellem de modsigende Momenter bringes istand, bliver Begrebets indre Selvudvikling
mulig.

Det er denne sidste Methode Professor Nielsen anvender i sit philosophiske System. „Som den Videnskab,
der overalt i det Tvetydige opdager det Tvesidige, og
ved Hjælp af Modsigelsen klarer Fordoblingen, er Philosophien \emph{Dialektik}“.\footnote{R. Nielsen. Forelæsninger over „Philosophisk Propædeutik.“
% NOTE: speck („overflødigt i“), 53 px vs full stop F=173 px, at x-height top in the word space; not transcribed.
Andet Cursus. S. 15.} Det vil her ikke være overflødigt i al Korthed at angive dens Forudsætninger.
Da Dialektiken hviler paa en Modsigelse, udfordres der
først to \emph{modsatte} Bestemmelser (f. Ex. formelt: Noget
og Andet; reelt: Ilt og Brint); dernæst fordres der, at
disse Bestemmelser, ikke ifølge ydre Tilfældighed, men
ifølge indre, begrebsmæssig Nødvendighed, \emph{skulle forudsætte} og \emph{betinge} hinanden (saaledes lader Noget
sig ikke gjennemtænke uden Andet, og omvendt. Ilten
forudsætter, for at virke, Brintens Nærværelse, og omvendt). I denne gjensidige Udelukkelse og Bekræftelse
fremkommer der et objektivt Reflexionsforhold (Speilingsforhold), hvori de modsatte Led forvandles til
\emph{Momenter} i den dialektiske Eenhed (den dialektiske
Eenhed, hvori Noget og Andet beherskes som Momenter,
er \emph{Forandringen}; den dialektiske Eenhed af Ilt og
Brint bliver af Naturen selv realiseret i den chemiske
Forbindelse: \emph{Vand}). Dette er i al Korthed Dialektikens
Forudsætninger, som i de forskjellige Anvendelser vel
% --- p. 91 --- (PDF 96)
\opage{91}\setcounter{footnote}{0}%
kunne modificeres paa mangfoldige Maader, men som i
Hovedsagen dog altid blive de samme.

Nu til Spørgsmaalet om denne Methodes Anvendelighed i Problemet Tro og Viden. Naar Troestheisme
og Videnstheisme skulle speile sig i hinanden, saa er
hermed Dialektiken \emph{fordret}; det kommer nu blot an
paa, om den \emph{er} gjennemført, og om den \emph{lader} sig
gjennemføre. Til Bedømmelsen heraf behøve vi kun at
see tilbage paa Forudsætningerne. Den første Betingelse er fuldkomment tilfredsstillet: Tro og Viden ere
saa \emph{modsatte} som vel muligt: de ere \textbf{absolut} \emph{ueensartede}. Men den anden Betingelse mangler, da Tro
og Viden hverken forudsætte eller betinge hinanden: „Med
sin iboende Grændse, sin indre Begrændsning, er Viden
ligeoverfor Troen aldeles uindskrænket, saa uindskrænket,
% p. 91n: printed „ꝛc.“ (Fraktur r rotunda + c, i.e. etc.), set here as \&c.
som om der slet ingen Tro var“ o. s. v.\footnote{R. Nielsen. Om „Den gode Villie“ \&c. S. 35.} Maaskee
vil der blive svaret, at Videns Grændse alligevel er
dialektisk, fordi den, foruden sin egen iboende Grændse,
har en Grændse i Mysteriet; men Mysteriet kan aldrig
hæves til Kategori; og naar Tro og Viden --- paa
Grund af Mysteriet --- umulig kunne indskrænke hinanden, bliver enhver Anvendelse af Dialektiken en Umulighed.

At opstille Tro og Viden som \textbf{absolut} \emph{ueensartede} Principer er ikke blot \emph{udialektisk}, men Udtryk
for en saadan \emph{Kriticisme}, at selv de Kantiske Antinomier træde i Skyggen; og samtidigt dermed at fordre
begge Sphærers indbyrdes Afspeiling --- hvis Mulighed
% --- p. 92 --- (PDF 97)
\opage{92}\setcounter{footnote}{0}%
intetsteds \textit{in concreto} paavises --- er Udtryk for en saa
umiddelbar \emph{Dogmatisme}, at den endog slaaer over i
det reent Populære. At det nu er umuligt af et eensidig
kritisk og et eensidig dogmatisk Element at tilveiebringe
nogen \emph{dialektisk Eenhed}, forstaaer sig af sig selv;
og man behøver kun at anvende den her angivne Synsmaade paa hvert eneste af de Punkter, som ovenfor bleve
udviklede i Anledning af Tro og Viden, for at overbevise sig om den dialektiske Methodes Uanvendelighed.
Det er overalt det Absolut-Ueensartede, der som en
snigende Gift gjennemtrænger det hele System og myrder
alt Tankeliv derved, at det myrder Dialektiken.

Vi kunne ikke nøies med idelige Forsikkringer om,
at Videnskaben \emph{kan} give rationelle Ækvivalenter for
det, man har maattet opgive i Troen, eller at Troestheisme og Videnstheisme \emph{kunne} afspeile sig i hinanden; ligesaalidt kunne vi tilfredsstilles ved en formeentlig Efterviisning af disse Muligheder i en Række
vilkaarligt valgte Exempler, som i deres Udførelse vel
kunne tiltale uphilosophiske Dannede ved den Overflødighed af Billeder og den Veltalenhed, hvoraf de bæres
op, men som for Enhver, der forstaaer at skjelne mellem
Billede og Begreb, forudsætte saa uendelig Meget, som
slet ikke bliver undersøgt, forudsætte det, hvorpaa det
netop kommer an. Hvad vi trænge til og ønske, ikke
vilkaarligt, men fordi det ligger i Sagens Natur, er en
ligefrem Indgaaen paa Religionsphilosophiens Problem,
en Indgaaen fra Grunden af, som liberalt giver Plads
for Skepticisme saaledes, at den selv gjør \emph{Muligheden} af den hele Antagelses fuldstændige Ugyldighed
% --- p. 93 --- (PDF 98)
\opage{93}\setcounter{footnote}{0}%
til Gjenstand for den omhyggeligste Undersøgelse --- thi
kun paa denne Betingelse bliver der i Virkeligheden
% NOTE: speck on the baseline at the end of the dash („--- og“), 29 px vs full stop F=207 px; not transcribed.
\emph{philosopheret}, --- og som endelig giver en Udvikling
i sammenhængende, systematisk og afsluttet \emph{Totalitet}.
Det er en Hovedmangel ved Professor Nielsens Udviklinger, at de i Reglen give Brudstykker, ikke Totalitet.

De meest paafaldende Exempler paa denne Mangel
findes blandt Forfatterens tidligste og blandt hans seneste
Arbeider. I Aaret 1841 udkom et Skrift under Titelen „Den speculative Logik i dens Grundtræk af Lic.
Rasmus Nielsen.“ Paa Omslaget af det første Hefte
meddeles der, at „Disse Grundtræk ere at betragte som
Fragment af en philosophisk Methodologi, hvis første
Deel skulde indeholde Logiken med en forudskikket Indledning. Nødvendigheden af at have en trykt Veiledning
for det mundtlige Foredrag har fremskyndet Udgivelsen. De
% p. 93: „Hæfter“ here, „Hefte“ two lines above and „Hefter“ below, as printed.
øvrige Hæfter ville følge, efterhaanden som den i Lectionscatalogen annoncerede Forelæsning rykker frem.“ Dette
Skrift har, i Henseende til \emph{Indholdet}, \emph{nu} kun forsaavidt Interesse, som det begynder Logiken med den
rene Væren og præsenterer sin Forfatter som den rene
% NOTE: stray vertical tick below the baseline after „Ulige“, 9x24 px, 149 px vs comma 335 px on the page (head below the baseline); not a comma, not transcribed.
Hegelianer. Ulige større Interesse har det i Henseende
til \emph{Formen}; naar Forfatteren selv betegner det som
Fragment, saa er dette at forstaae saaledes, at det er
blevet et Fragment af et Fragment; thi den belovede
„forudskikkede Indledning“ kom aldrig, ligesaalidt som de
% NOTE: a stray oblique stroke fused to the f of „fulgte“ (p. 93); not a separate sort, not transcribed.
manglende Hefter fulgte. Skriftet, som gaaer fra Side
1 til Side 192, begynder til en Forandring med § 11
og ender inde i en Sætning med Ordene „Forsaavidt
% ===== batch b11_pp94-95 =====
% --- p. 94 ---
% (PDF 99)
\opage{94}\setcounter{footnote}{0}%
den“. Ved et besynderligt Spil af Skjæbnen --- thi
af Forsynet kan det vel neppe være --- hedder det
sidste Afsnit: \textit{c)} \textbf{Totalitet.} Under denne Overskrift
(S. 180) læser man: „\emph{Den med sine ydre Former
incommensurable Inderlighed ytrer sig i
uudtømmelig Productivitet;} Identiteten af Indre
og Ydre reflecteres i Kraften, og Totaliteten udvikler sig
til et Forhold mellem Kraft og Product. Som Formationen
af en ny Gestalt forudsætter den producerende
Kraft, saa forudsætter Kraftens specifiske Productivitet
\emph{producerede Gestalter i det Uendelige}, saa at
enhver \emph{Realformation} reflecterer sin \emph{Præformation}.“
Tage vi nu dette embryonale Værk som \emph{Præformationen}
og „Grundideernes Logik“ som \emph{Realformationen},
saa slaae Forfatterens Ord mærkværdig godt
til, at nemlig den med sine ydre Former inkommensurable
Inderlighed yttrer sig i uudtømmelig Produktivitet;
thi efter den Inddeling, som angives i Indholdsfortegnelserne
i „Grundideernes Logik“, udgjør Anden Deel $\frac{2}{243}$
af det Hele, eller tydeligere: Anden Deel indeholder \emph{to
Underafdelinger af anden Afdeling af andet
Kapitel af første Afsnit af første Deel af det
hele Værk}. Det Tilbagestaaende af Værket, udført efter
denne Maalestok, vil udgjøre 114$\frac{1}{2}$ Deel. Nemlig:
% TABLE (p. 94). Printed layout: a centred bold heading „(Grundideernes
% Logik. I.)“ (Fraktur bold; „I.“ bold antiqua; full stop after „Logik“),
% then three rows set to the full measure, no rules. Measured at 300 dpi
% (measure = 100): (1) label column, 0--58; in row 1 „Videns Idee“ then,
% after a gap of one point-pitch, eight spaced full stops (connected
% components; pitch 4.6) ending at the column's right edge; in rows 2--3
% the label is spread by wide word spaces to fill the column; (2) the
% computation „=1/3 × 1/3“, 59--71 (rows 2--3); in row 1 only a stray
% „-“ (a single short raised stroke, not the Fraktur double hyphen) at
% 68--69; (3) the result „=“, 72--75 in rows 2--3, but 74--77 in row 1,
% where it is set tight against its fraction; (4) the result fractions,
% one under the other (centres at 78); (5) „af hele Værket“ (row 1, at a
% word space after 1/3, running to the right margin), a ditto mark
% centred under it in rows 2--3. Fractions are small built-up case
% fractions. Here: tabular* to \linewidth with column (1) as a p-column of
% the printed width, the spread of the labels by \hfill between words; the
% leaders are the printed count, at the printed pitch (\hspace{1.25em}).
\begin{center}
\textbf{(Grundideernes Logik. \textit{I.})}\\[1.5ex]
\renewcommand{\arraystretch}{1.3}%
\begin{tabular*}{\linewidth}{@{}p{0.575\linewidth}@{\extracolsep{\fill}\ }r@{\extracolsep{0pt}\ }r@{\,}c@{\ }c@{}}
% PRINTED AS IS: stray „-“ before „=“ in the „Videns Idee“ row.
Videns Idee\hfill.\hspace{1.25em}.\hspace{1.25em}.\hspace{1.25em}.\hspace{1.25em}.\hspace{1.25em}.\hspace{1.25em}.\hspace{1.25em}. & -\kern0.5em & $=$ & $\frac{1}{3}$ & af hele Værket\\
% ditto marks: printed „〃“ (two short slanted strokes, the sort used for
% the closing quotation mark), standing for „af hele Værket“; here ''.
Første \hfill Afsnit: \hfill Den \hfill subjective \hfill Viden & $=\frac{1}{3}\times\frac{1}{3}$ & $=$\hspace{0.45em} & $\frac{1}{9}$ & ''\\
Andet \hfill Kapitel: \hfill Objectivitetens \hfill Logik & $=\frac{1}{3}\times\frac{1}{9}$ & $=$\hspace{0.45em} & $\frac{1}{27}$ & ''\\
\end{tabular*}
\end{center}
% --- p. 95 ---
% (PDF 100)
% TABLE (p. 95). Printed layout: a centred bold heading „(Grundidecrnes
% Logik II.)“, then five rows to the full measure, no vertical rules.
% Measured at 300 dpi (measure = 100): (1) label, leaders and computation
% together, 0--72: the label at the left; the spaced full stops stand on
% a common grid (pitch 4.7; B., a., (I), (II) share positions) and run up
% to about half a pitch before the first „=“ of the computation;
% connected components give B.: five; a.: „hver.“ with its point at
% punctuation distance plus three (four in all); (I): three; (II): seven;
% row b.: none; the computations end flush one under another at 72;
% (2) the result „=“ one under the other at 72.5, set tight against its
% fraction; (3) „af hele Værket“ (row B., at a word space after 1/81, to
% the right margin), a ditto mark centred under it in the other four
% rows. Rows a. and b. are indented by about half an em under „B.“ (1.9),
% and a single right brace „}“ spanning both rows joins „a.
% Formalfunctioner“ / „b. Modalfunctioner“ to „hver“ / „altsaa
% tilsammen“, which begin together after it. Row b. and row (II) have no
% computation, only „= 2/243“ in the result column. In row (I) the second
% term is a stacked compound fraction: „(1/27)“ over „3“. A full-width
% rule follows the table, then the totals line „Resten = …“ set across the
% full measure at the left margin (not indented). Labels B., a., b., (I),
% (II) and all numerals are antiqua („B.“ bold). The minus signs in the
% totals line are printed as rules as long as „=“; here $-$. Here: a
% tabular* with column (1) a p-column of the printed width, \hfill between
% label and leaders, the leaders at the printed pitch (\hspace{1.25em})
% and \enspace before the computation.
\begin{center}
\opage{95}\setcounter{footnote}{0}%
% PRINTED AS IS: heading „(Grundideernes Logik II.)“ has no full stop
% after „Logik“, unlike p. 94 „(Grundideernes Logik. I.)“.
% PRINTED AS IS: „Grundidecrnes“ (for Grundideernes): the second e of the
% heading prints as a complete bold Fraktur c (wrong sort): the head's
% right arm tapers to a formed pointed terminal (stroke 12--11--9--7--5 px
% at 300 dpi) and stops, with no stub of a crossbar on the arm or on the
% stem, whereas both e's beside it are closed by a crossbar as heavy as
% the stem; the glyph has 835 black pixels against 949 and 956 for those
% e's, same set width.
\textbf{(Grundidecrnes Logik \textit{II.})}\\[1ex]
\renewcommand{\arraystretch}{1.3}%
\begin{tabular*}{\linewidth}{@{}p{0.70\linewidth}@{\extracolsep{\fill}\,}l@{\extracolsep{0pt}\ }c@{}}
\textbf{\textit{B.}} Logiske Functioner\hfill.\hspace{1.25em}.\hspace{1.25em}.\hspace{1.25em}.\hspace{1.25em}.\enspace$=\frac{1}{3}\times\frac{1}{27}$ & $=\frac{1}{81}$ & af hele Værket\\
\hspace*{0.5em}\textit{a.} Formalfunctioner\ \rlap{\smash{\raisebox{-0.62\baselineskip}{$\left.\rule[-0.95\baselineskip]{0pt}{1.9\baselineskip}\right\}$}}}\hspace{0.7em}hver.\hfill.\hspace{1.25em}.\hspace{1.25em}.\enspace$=\frac{1}{3}\times\frac{1}{81}$ & $=\frac{1}{243}$ & ''\\
\hspace*{0.5em}\rlap{\textit{b.} Modalfunctioner}\phantom{\textit{a.} Formalfunctioner}\ \hspace{0.7em}altsaa tilsammen & $=\frac{2}{243}$ & ''\\
% speck below the „=“ in row (I), not a character.
\textit{(I)} Første Deel bliver da\hfill.\hspace{1.25em}.\hspace{1.25em}.\enspace$=\frac{1}{27}+\frac{(\frac{1}{27})}{3}$ & $=\frac{12}{243}$ & ''\\
% (II): the leaders stop where row B.'s computation begins, as printed.
\textit{(II)} Anden Deel\hfill.\hspace{1.25em}.\hspace{1.25em}.\hspace{1.25em}.\hspace{1.25em}.\hspace{1.25em}.\hspace{1.25em}.\enspace\phantom{$=\frac{1}{3}\times\frac{1}{27}$} & $=\frac{2}{243}$ & ''\\
\end{tabular*}
\end{center}
\noindent\rule[0.5ex]{\linewidth}{0.8pt}\\
Resten $=$ Det hele Værk $-$ \textit{(I$+$II)} $= 1 - (\frac{12}{243}+\frac{2}{243}) = \frac{229}{243}$.

Forholdet mellem Resten \textit{(R)} og Anden Deel \textit{(II)}
bliver da:
% Display: not centred but indented: it begins at 22 % of the measure and
% ends at 61 % (300 dpi: measure 446--2984, display 1008--1991); reproduced
% with flushleft and \hspace*. R/II full size antiqua (\mathrm in math),
% the inner fractions small case fractions in parentheses, „229/2“ a small
% case fraction.
\begin{flushleft}
\hspace*{0.22\linewidth}$\displaystyle\frac{\mathrm{R}}{\mathrm{II}}=\frac{(\frac{229}{243})}{(\frac{2}{243})}=\tfrac{229}{2}=114\tfrac{1}{2}.$
\end{flushleft}

% PRINTED AS IS: „ndgivet“ (for udgivet; turned u).
% NOTE: speck beside the comma after „nu“ (0.40 F, F = 174 px), not transcribed.
Antage vi nu, at der blev ndgivet een Deel aarligt,
saa vilde Værket fra nu af kunne endes i 114$\frac{1}{2}$
Aar. Hermed ere vi dog endnu ikke færdige, thi efter
„Grundideernes Logik“ maatte følge en tilsvarende Naturphilosophi
og Aandsphilosophi, og naar vi saa have
naaet Methusalems Alder, kan Diskussionen begynde.

At anlægge et Værk efter en Plan, der fornuftigviis
ikke lader sig gjennemføre, er en Selvmodsigelse; og
% „Forfatteren“: the o is a damaged sort (broken bowl); reading o.
dog har Forfatteren en Opfordring til at gaae videre,
thi ved at udgive de to første Dele af „Grundideernes
Logik“ har han, ideelt forstaaet, skrevet sig selv et
Tvangspas, der forbyder ham at standse. At regne paa
de 114$\frac{1}{2}$ Aar vilde, efter den Probabilitet, som gjælder
for Levealderen i den nuværende Generation, være overilet.
At give den store Rest i sammentrængt Fremstilling,
vilde være uheldigt, deels paa Grund af den derved
indtrædende Mangel paa Egalitet, deels fordi en
stærk Sammendragning vilde gjøre Udviklingen vanskelig
at forstaae. Skjøndt det naturligviis maa beroe paa
% ===== batch b12_pp96-98 =====
% ==== batch b12: printed pp. 96--98 (PDF 101--103) ====
% --- p. 96 ---
% (PDF 101) p. 96 continues section IV mid-sentence („... beroe paa“ / „Forfatterens“).
\opage{96}\setcounter{footnote}{0}Forfatterens egen Afgjørelse, hvorledes denne Plan skal
gjennemføres, saa synes det dog fra Læserens Standpunkt
at være det Ønskeligste, om Forfatteren reent opgav
denne Plan, lod „Grundideernes Logik“ ligesom
„Den speculative Logik“ staae hen som interessant Ruin,
og istedetfor Alt dette, indenfor rimelige Grændser udviklede
sit System heelt, klart og afrundet, at man dog
% „komme“: the last minim of the second m printed detached (reads like „komnie“); worn sort, not a different word.
eengang kunde komme til en Forstaaelse af og med
Forfatteren.

% Section IV ends here with a centred rule (c. 47 mm) below its last line.
% Then the head „Efterskrift.“, centred, in bold Fraktur (stroke c. 13 px
% against 9 px for body text at 300 dpi, and larger), with a shorter centred
% rule (c. 23 mm) below it. A small ink blot to the left of that rule is dirt.
\begin{center}\textbf{Efterskrift.}\end{center}

Idet jeg hermed udsender denne Kritik, er jeg mig
bevidst paa ethvert Punkt at have holdt mig til Sagen.

Det har været mig magtpaaliggende at ledsage enhver
Udtalelse med Grunde; det forstaaer sig af sig selv,
at jeg kan have taget feil eller valgt eensidige Synspunkter
i Argumentationen; men hvorledes det nu ogsaa
forholder sig, saa har Undersøgelsen overalt Begrundelsens
Charakteer. Saafremt der altsaa kommer et
Svar, forlanger jeg som min Ret, \emph{at ethvert Punkt
deri belægges med Grunde}.

Jeg har fremdraget \emph{afgjørende Hovedpunkter}
og accentueret dem som Hovedpunkter. Jeg fordrer
derfor, at man i det eventuelle Svar \emph{ikke nøies med}
at henvise til hele „Grundideernes Logik“ eller til store
Afsnit deraf --- hvad der jo i en Undersøgelse som den
% --- p. 97 ---
% (PDF 102) Sheet signature „7“ at the foot of p. 97 dropped. Much worn type on
% pp. 97--98: several e's print without the eye and read like c („frie“,
% „respektere“, „Grunde.“, „kjender“, „Studerende“, „være“), transcribed as e.
\opage{97}\setcounter{footnote}{0}foreliggende vilde være aldeles intetsigende --- fremdeles,
at man \emph{ikke indskrænker sig} til at fremdrage enkelte
uvæsentlige Bestemmelser paa en saadan Maade, at
Hovedpunkterne glemmes, \emph{eller lader}, som om Hovedpunkterne
slet ikke vare fremdragne. Om jeg indlader
mig med andre Indlæg end dem, jeg muligen kunde
vente fra Professor Nielsen selv, lader jeg være afhængigt
af deres Beskaffenhed.

Kan man med Grunde overbevise mig om Feil,
Eensidigheder og Misforstaaelser, \emph{skal jeg strax og
resolut indrømme det}, for at Undersøgelsen kan
gaae videre. Det er mig aldeles ikke magtpaaliggende
at „faae Ret“; det evige Pindehuggeri om Bagateller
sløver Interessen for Videnskaben og hæmmer den Spændkraft,
som ligger i den frie Aandsudvikling. Men saavist
som jeg ubetinget skal respektere Grunde, saavist vil
% NOTE: speck („Intet respektere“), 59 px vs full stop F=141 px, low in the x-height right after „Intet“; not transcribed.
jeg Intet respektere uden Grunde. Meningstilkjendegivelser,
Autoritetsdomme o. desl. anseer jeg i en videnskabelig
Undersøgelse for intetsigende og --- komiske.

Naar den Prøvelse af Professor Nielsens Philosophi,
jeg her har anstillet, gaaer i modsat Retning af
mine tidligere Udtalelser, saa vil jeg baade for Læserens
og for min egen Skyld lade Forklaringen medfølge.
% „Det vil være ...“ is set flush left after a full line: not a new paragraph.
Det vil være en Deel af Læserne bekjendt, at jeg skylder
Professor Nielsen min philosophiske Udvikling, og jeg
maa tilføie, at jeg herfor er ham stor Tak skyldig. For
Enhver, der kjender Professor Nielsens Begavelse, hans
Produktivitet og Veltalenhed, vil det være let forstaaeligt,
at den yngre Studerende kan føle sig saa stærkt
tiltrukken, at han, naar Indtrykkene stadigt fornyes, og
% --- p. 98 ---
% (PDF 103)
\opage{98}\setcounter{footnote}{0}naar Læren ubetinget og peremptorisk fremstilles som den
ene sande, omsider kan komme til at see omtrent Alt gjennem
sin Lærers Briller. Men alt som den yngre Studerende
bliver en ældre Studerende, begynder han at ville see
Tingene med sine egne Øine og anlægge sin egen
Maalestok. Paa dette Stadium i Udviklingen gives der
mangfoldige Overgange, hvor det Videnskabelige og det
Personlige, Grundenes Magt og Autoritetsmagten, kunne
% „eller“: both e's without the eye and the first l short (worn sorts; reads like „ciler“).
kæmpe med hinanden; men tidligt eller sildigt maa det
Vendepunkt indtræde, hvor det at hævde sin Selvstændighed
og følge sin egen Overbeviisning fremtræder som
Pligt. Mig er det gaaet paa denne Maade. Det er
navnlig efter et længere Ophold i Udlandet, at jeg
er kommen til at betragte de tilvante Begrebsanalyser
og Tankeformer paa Afstand. Vistnok nærede jeg allerede
inden min Afreise stærke Tvivl om Holdbarheden
af Professor Nielsens Lære om Tro og Viden; nu er
denne Tvivl bleven til Vished. Naar hertil kommer,
at jeg tilsidst befandt mig i det energiske Dilemma,
enten at fornegte min Overbeviisning eller at bryde med
den Theori, jeg tidligere hyldede, saa var der for mig
intet Valg; \emph{thi af min Overbeviisning offrer
jeg Intet for nogensomhelst Priis}. Idet jeg
% PRINTED AS IS: „til ikke udtale“ (no „at“ before „udtale“).
føler mig for stolt til ikke udtale, at jeg er min gamle
Lærer taknemmelig for de mange Aars videnskabelige
Samliv, undseer jeg mig ikke ved at tilføie, at et Offer
af den nævnte Art vilde være, hvad Ordsproget kalder
„at kjøbe Guld for dyrt“.

Hermed udsendes disse Linier til Læserens, ikke
velvillige, men skarpest mulige Prøvelse.

% Closing rule, centred, below the last line of p. 98: a double rule (two thin
% parallel rules, 806 and 789 px at 300 ppi, c. 68 mm in the scan, 0.32 of the
% 2530 px measure, 18 px apart centre to centre, both broken along their
% length; set like the double rule closing the Forord, p. 9). Nothing else on the page:
% no date, place, signature or printer's line. End of the book.
% PDF 104 = blank end-paper leaf (no printing; scan noise and the spine edge only).
% PDF 105 = back board, exterior: dark green marbled paper with cloth spine and
% corners (colour photograph, 72 ppi).
\begin{center}\rule[2.2pt]{0.32\textwidth}{0.4pt}\hspace*{-0.32\textwidth}\rule{0.32\textwidth}{0.4pt}\end{center}

\end{document}
